% Calcul littéral — Mathématiques 3ème — Cours signé OnBuch+
% Programme officiel MINESEC. Compiler avec : tectonic N05-calcul-litteral.tex
\newcommand{\DOCMATIERE}{Mathématiques}
\newcommand{\DOCNIVEAU}{3ème}
\newcommand{\DOCMODULE}{Mathématiques – classe de 3ème}
\newcommand{\DOCLECON}{N05}
\newcommand{\DOCTITRE}{Calcul littéral}
% OnBuch+ — préambule « cours signés » ENRICHI (compilé avec tectonic / XeLaTeX)
% Système visuel « Pop » d'OnBuch+ : fond crème (jamais blanc), bordures encre
% épaisses, ombres « dures » décalées, titres Archivo Black en capitales.
% Version MATHÉMATIQUES : graphes (pgfplots/tikz), boîtes pédagogiques
% (définition, propriété, méthode, attention, le savais-tu, exercice résolu,
% exercice, TP, bilan), page de garde et signature OnBuch+ ancrée en bas.
%
% Macros attendues dans main.tex AVANT \input{preamble} :
%   \DOCMATIERE  \DOCNIVEAU  \DOCMODULE  \DOCLECON (numéro)  \DOCTITRE
\documentclass[11pt]{article}

\usepackage{fontspec}
\usepackage[french]{babel}
\usepackage[a4paper,margin=2cm,top=3.4cm,bottom=2.8cm,headheight=1.4cm,footskip=1.3cm]{geometry}
\usepackage{amsmath,amssymb,amsfonts}
\usepackage{mathtools} % \xmapsto, \coloneqq... souvent utilisés par les modèles
\usepackage{cancel} % simplifications barrées
% \abs{z} (module, valeur absolue) et \norm{u} (norme), utilisés par les modèles
\providecommand{\abs}{}\let\abs\relax\DeclarePairedDelimiter{\abs}{\lvert}{\rvert}
\providecommand{\norm}{}\let\norm\relax\DeclarePairedDelimiter{\norm}{\lVert}{\rVert}
\usepackage[table]{xcolor} % [table] : fournit aussi \rowcolors (alternance de lignes)
\usepackage{pagecolor}
\usepackage{tikz}
\usetikzlibrary{arrows.meta,positioning,calc,shapes.geometric,shapes.misc,shapes.arrows,decorations.pathmorphing,decorations.pathreplacing,decorations.markings,patterns,shadows,mindmap,trees,fit,backgrounds,matrix,intersections,angles,quotes}
\usepackage{pgfplots}
\usepackage{pgf-pie} % diagrammes circulaires \pie (statistiques)
\pgfplotsset{compat=1.18}
\usepgfplotslibrary{fillbetween,groupplots} % aires entre courbes (calcul intégral)
\usepackage{enumitem}
\usepackage{fancyhdr}
% [most] charge toutes les bibliothèques tcolorbox (listings, minted, raster,
% documentation...) et gonfle inutilement la mémoire TeX sur un document long
% (~300 boîtes) : on ne charge que ce qui est réellement utilisé.
\usepackage[skins,breakable]{tcolorbox}
\usepackage{graphicx}
\usepackage{colortbl}
\usepackage{booktabs}
\usepackage{tabularx}
\usepackage{diagbox} % case d'en-tête diagonale des tableaux à double entrée
% Colonnes extensibles centrée / gauche / droite (C, L, R), souvent utilisées par les modèles
\newcolumntype{C}{>{\centering\arraybackslash}X}
\newcolumntype{L}{>{\raggedright\arraybackslash}X}
\newcolumntype{R}{>{\raggedleft\arraybackslash}X}
\usepackage{array}
\usepackage{multicol}
\usepackage{siunitx}
% parse-numbers=false : accepte aussi \SI{2,5 \times 10^{-3}}{...}
\sisetup{locale=FR,output-decimal-marker={,},group-minimum-digits=4,parse-numbers=false}
\DeclareSIUnit\ohm{\ensuremath{\symup{\Omega}}} % U+2126 absent de la police
\usepackage{qrcode}
% \degree : commande fréquemment utilisée par les modèles pour un angle ou une
% température en dehors de \SI{}{} (non fournie par les paquets chargés).
\providecommand{\degree}{^{\circ}}
% \qed (fin de démonstration), utilisé par les modèles sans amsthm
\providecommand{\qed}{\hfill$\square$}
% \barre : double barre des valeurs interdites dans un tableau de variations (tkz-tab)
\providecommand{\barre}{\ensuremath{\|}}
% environnement proof (amsthm non chargé) utilisé par les modèles
\ifdefined\proof\else\newenvironment{proof}[1][Démonstration]{\par\noindent\textit{#1.}\ }{\qed\par}\fi
\usepackage{lastpage}
\usepackage{hyperref}
\hypersetup{hidelinks}

% ============================================================
% Palette « Pop » OnBuch+ — mêmes hex que la classe P (plus_ui.dart)
% ============================================================
\definecolor{poporange}{HTML}{F59321}
\definecolor{popdark}{HTML}{E07A0C}
\definecolor{popcream}{HTML}{FFF6E8}
\definecolor{popink}{HTML}{1C1714}
\definecolor{poppurple}{HTML}{7A5AE0}
\definecolor{popgreen}{HTML}{2E9E6A}
\definecolor{poppink}{HTML}{E0457B}
\definecolor{popblue}{HTML}{2F7DE1}
\definecolor{popgold}{HTML}{C98A12}
\definecolor{popmuted}{HTML}{8A7F76}
\definecolor{popline}{HTML}{EADFD2}
\definecolor{popgray}{HTML}{8A7F76}
\colorlet{popgrey}{popgray}
% Alias tolérants (le modèle nomme parfois les couleurs différemment)
\colorlet{popwhite}{white}
\colorlet{popblack}{popink}
\colorlet{popred}{poppink}
\colorlet{popyellow}{popgold}
% Teintes claires pour fonds de tableaux / zones
\colorlet{popblueL}{popblue!10!white}
\colorlet{poporangeL}{poporange!14!white}
\colorlet{popgreenL}{popgreen!12!white}
\colorlet{poppurpleL}{poppurple!12!white}
\colorlet{poppinkL}{poppink!12!white}
\colorlet{popgoldL}{popgold!14!white}

\pagecolor{popcream}

% ============================================================
% Typographie
% ============================================================
\defaultfontfeatures{Ligatures=TeX}
\setmainfont{PlusJakartaSans-Medium.ttf}[Path=../../../fonts/,BoldFont=PlusJakartaSans-Bold.ttf]
% Maths en sans-serif (Fira Math) avec chiffres et lettres droites de Plus
% Jakarta, pour que les formules se fondent dans le texte.
\usepackage[math-style=ISO,bold-style=ISO]{unicode-math}
\setmathfont{FiraMath-Regular.otf}[Path=../../../fonts/]
\setmathfont{PlusJakartaSans-Medium.ttf}[Path=../../../fonts/,range={up/num,up/{Latin,latin},\mathup}]
\usepackage{icomma} % virgule décimale sans espace en mode math
% Caractères mathématiques Unicode absents des polices de texte (Plus Jakarta,
% Archivo Black) : rendus par leur équivalent mathématique, en texte comme en
% formule — sinon ils s'affichent en carré vide (ex. « Arithmétique dans ℤ »).
\usepackage{newunicodechar}
\newunicodechar{ℝ}{\ensuremath{\mathbb{R}}}
\newunicodechar{ℤ}{\ensuremath{\mathbb{Z}}}
\newunicodechar{ℂ}{\ensuremath{\mathbb{C}}}
\newunicodechar{ℕ}{\ensuremath{\mathbb{N}}}
\newunicodechar{ℚ}{\ensuremath{\mathbb{Q}}}
\newunicodechar{𝔻}{\ensuremath{\mathbb{D}}}
\newunicodechar{α}{\ensuremath{\alpha}}
\newunicodechar{β}{\ensuremath{\beta}}
\newunicodechar{γ}{\ensuremath{\gamma}}
\newunicodechar{δ}{\ensuremath{\delta}}
\newunicodechar{ε}{\ensuremath{\varepsilon}}
\newunicodechar{θ}{\ensuremath{\theta}}
\newunicodechar{λ}{\ensuremath{\lambda}}
\newunicodechar{μ}{\ensuremath{\mu}}
\newunicodechar{σ}{\ensuremath{\sigma}}
\newunicodechar{τ}{\ensuremath{\tau}}
\newunicodechar{φ}{\ensuremath{\varphi}}
\newunicodechar{ψ}{\ensuremath{\psi}}
\newunicodechar{ω}{\ensuremath{\omega}}
\newunicodechar{Σ}{\ensuremath{\Sigma}}
% majuscules produites par \MakeUppercase dans les titres (α-aminés -> Α)
\newunicodechar{Α}{\ensuremath{\alpha}}
\newunicodechar{Β}{\ensuremath{\beta}}
\newunicodechar{Γ}{\ensuremath{\Gamma}}
\newunicodechar{Δ}{\ensuremath{\Delta}}
\newunicodechar{Ε}{\ensuremath{\varepsilon}}
\newunicodechar{Θ}{\ensuremath{\Theta}}
\newunicodechar{Λ}{\ensuremath{\Lambda}}
\newunicodechar{Μ}{\ensuremath{\mu}}
\newunicodechar{Τ}{\ensuremath{\tau}}
\newunicodechar{Φ}{\ensuremath{\Phi}}
\newunicodechar{Ψ}{\ensuremath{\Psi}}
\newunicodechar{Ω}{\ensuremath{\Omega}}
\newunicodechar{↦}{\ensuremath{\mapsto}}
\newunicodechar{∈}{\ensuremath{\in}}
\newunicodechar{∉}{\ensuremath{\notin}}
\newunicodechar{∧}{\ensuremath{\wedge}}
\newunicodechar{⇔}{\ensuremath{\Leftrightarrow}}
\newunicodechar{⟺}{\ensuremath{\Longleftrightarrow}}
\newunicodechar{⇒}{\ensuremath{\Rightarrow}}
\newunicodechar{⟹}{\ensuremath{\Longrightarrow}}
\newunicodechar{∘}{\ensuremath{\circ}}
\newunicodechar{∪}{\ensuremath{\cup}}
\newunicodechar{∩}{\ensuremath{\cap}}
\newunicodechar{≡}{\ensuremath{\equiv}}
\newunicodechar{⊂}{\ensuremath{\subset}}
\newunicodechar{⊥}{\ensuremath{\perp}}
\newunicodechar{∃}{\ensuremath{\exists}}
\newunicodechar{∀}{\ensuremath{\forall}}
\newunicodechar{∛}{\ensuremath{\sqrt[3]{\,}}}
\newunicodechar{∜}{\ensuremath{\sqrt[4]{\,}}}
\newunicodechar{⃗}{\ensuremath{{}^{\to}}}
\newunicodechar{ⁿ}{\ifmmode^{n}\else\textsuperscript{\ensuremath{n}}\fi}
\newunicodechar{ˣ}{\ifmmode^{x}\else\textsuperscript{\ensuremath{x}}\fi}
\newunicodechar{ᵇ}{\ifmmode^{b}\else\textsuperscript{\ensuremath{b}}\fi}
\newunicodechar{ʳ}{\ifmmode^{r}\else\textsuperscript{\ensuremath{r}}\fi}
\newunicodechar{ᵘ}{\ifmmode^{u}\else\textsuperscript{\ensuremath{u}}\fi}
\newunicodechar{ʸ}{\ifmmode^{y}\else\textsuperscript{\ensuremath{y}}\fi}
\newunicodechar{ᵃ}{\ifmmode^{a}\else\textsuperscript{\ensuremath{a}}\fi}
\newunicodechar{ᵉ}{\ifmmode^{e}\else\textsuperscript{\ensuremath{e}}\fi}
\newunicodechar{ᵏ}{\ifmmode^{k}\else\textsuperscript{\ensuremath{k}}\fi}
\newunicodechar{ᵖ}{\ifmmode^{p}\else\textsuperscript{\ensuremath{p}}\fi}
\newunicodechar{ᴺ}{\ifmmode^{N}\else\textsuperscript{\ensuremath{N}}\fi}
\newunicodechar{ᶜ}{\ifmmode^{c}\else\textsuperscript{\ensuremath{c}}\fi}
\newunicodechar{ʰ}{\ifmmode^{h}\else\textsuperscript{\ensuremath{h}}\fi}
\newunicodechar{ᵠ}{\ifmmode^{q}\else\textsuperscript{\ensuremath{q}}\fi}
\newunicodechar{ₙ}{\ifmmode_{n}\else\textsubscript{\ensuremath{n}}\fi}
\newunicodechar{ₐ}{\ifmmode_{a}\else\textsubscript{\ensuremath{a}}\fi}
\newunicodechar{ᵢ}{\ifmmode_{i}\else\textsubscript{\ensuremath{i}}\fi}
\newunicodechar{ᵣ}{\ifmmode_{r}\else\textsubscript{\ensuremath{r}}\fi}
\newunicodechar{ₖ}{\ifmmode_{k}\else\textsubscript{\ensuremath{k}}\fi}
\newunicodechar{ᵦ}{\ifmmode_{\beta}\else\textsubscript{\ensuremath{\beta}}\fi}
\newunicodechar{ₓ}{\ifmmode_{x}\else\textsubscript{\ensuremath{x}}\fi}
\newfontfamily\popdisplay{ArchivoBlack-Regular.ttf}[Path=../../../fonts/]
\newfontfamily\poplogo{SpaceGrotesk-Bold.ttf}[Path=../../../fonts/]
\newfontfamily\popsemi{PlusJakartaSans-SemiBold.ttf}[Path=../../../fonts/]
\newfontfamily\popxbold{PlusJakartaSans-ExtraBold.ttf}[Path=../../../fonts/]
\newfontfamily\popxboldls{PlusJakartaSans-ExtraBold.ttf}[Path=../../../fonts/,LetterSpace=3.0]

\setlist[enumerate]{leftmargin=1.7em,itemsep=5pt,topsep=4pt}
\setlist[itemize]{leftmargin=1.7em,itemsep=4pt,topsep=4pt,label={\color{poporange}\textbullet}}
\setlength{\parindent}{0pt}
\setlength{\parskip}{5pt}
\renewcommand{\arraystretch}{1.25}

% pgfplots : style Pop par défaut
\pgfplotsset{
  every axis/.append style={axis line style={popink,line width=1pt},tick style={popink},
    grid style={popline,line width=0.5pt},label style={font=\small},tick label style={font=\footnotesize}},
  popcurve/.style={poporange,line width=1.8pt,no markers,smooth},
}
% Même style disponible dans un \draw TikZ simple (hors pgfplots)
\tikzset{popcurve/.style={poporange,line width=1.8pt}}

% ============================================================
% Pill / squiggle
% ============================================================
\newcommand{\poppill}[3]{%
  \tikz[baseline=-0.55ex]\node[draw=popink,line width=1.2pt,rounded corners=2.6pt,%
    fill=#2,inner xsep=6.5pt,inner ysep=3pt]{%
    \popxboldls\fontsize{7}{7}\selectfont\color{#3}\MakeUppercase{#1}};%
}
\newcommand{\popsquiggle}[1]{%
  \tikz[baseline=-0.4ex]{
    \draw[#1,line width=2.6pt,line cap=round]
      plot [smooth] coordinates {(0,0.09) (0.34,0.01) (0.68,0.21) (1.02,0.01) (1.36,0.10)};
  }%
}
% Mot-clé surligné (marqueur orange sous le texte)
\newcommand{\cle}[1]{{\popxbold\color{popdark}#1}}

% ============================================================
% En-tête / pied
% ============================================================
\pagestyle{fancy}
\fancyhf{}
\renewcommand{\headrulewidth}{0pt}
\renewcommand{\footrulewidth}{0pt}
\lhead{\raisebox{-0.15ex}{\poplogo\fontsize{13}{13}\selectfont\color{popink}OnBuch\color{poporange}+}}
\rhead{\poppill{\DOCMATIERE\ \textbullet\ \DOCNIVEAU\ \textbullet\ Leçon \DOCLECON}{white}{popink}}
\lfoot{\popxbold\fontsize{7.6}{7.6}\selectfont\color{popmuted}\MakeUppercase{Cours signé OnBuch+}\ \textbullet\ {\popsemi\DOCTITRE}}
\rfoot{\tikz[baseline=-0.6ex]\node[rounded corners=8.5pt,fill=popink,minimum height=17pt,minimum width=17pt,inner xsep=5pt,inner sep=0]{\popxbold\fontsize{8}{8}\selectfont\color{popcream}\thepage\,/\,\pageref*{LastPage}};}

% ============================================================
% Titres
% ============================================================
\newcommand{\chaptitre}[2]{%
  \begin{tcolorbox}[enhanced,colback=poporange,colframe=popink,boxrule=2.6pt,arc=11pt,
      drop shadow=popink,left=15pt,right=15pt,top=12pt,bottom=11pt,before skip=3pt,after skip=18pt]
    \poppill{#2}{popink}{popcream}\par\vspace{9pt}
    {\raggedright\hyphenpenalty=10000\exhyphenpenalty=10000\popdisplay\fontsize{22}{24}\selectfont\color{popcream}\MakeUppercase{#1}\par}
  \end{tcolorbox}
}
% Section numérotée : pastille orange avec le numéro + titre Archivo
\newcounter{popsec}
\newcommand{\coursec}[1]{%
  \stepcounter{popsec}\setcounter{popsub}{0}%
  \par\vspace{16pt}\noindent
  \begin{minipage}{\linewidth}
  \tikz[baseline=-0.7ex]\node[draw=popink,line width=1.6pt,fill=poporange,rounded corners=4pt,minimum size=22pt,inner sep=0]{\popdisplay\fontsize{12}{12}\selectfont\color{popcream}\thepopsec};%
  \hspace{8pt}{\popdisplay\fontsize{15}{17}\selectfont\color{popink}\MakeUppercase{#1}}\par
  \vspace{4pt}\noindent{\color{poporange}\rule{36pt}{3.6pt}}
  \end{minipage}\par\nopagebreak\vspace{8pt}
}
\newcounter{popsub}
\newcommand{\courssub}[1]{%
  \stepcounter{popsub}%
  \par\vspace{9pt}\noindent{\popxbold\fontsize{11.5}{13}\selectfont\color{popink}{\color{poporange}\thepopsec.\thepopsub}\hspace{6pt}#1}\par\nopagebreak\vspace{4pt}
}

% ============================================================
% Boîtes pédagogiques — même recette : fond blanc, bordure épaisse colorée,
% ombre dure encre, badge plein. Titre optionnel [#1].
% ============================================================
\tcbset{popbase/.style={breakable,enhanced,colback=white,boxrule=2.3pt,arc=9pt,
  shadow={3pt}{-3pt}{0pt}{popink},left=14pt,right=14pt,top=11pt,bottom=11pt,before skip=11pt,after skip=15pt,
  fontupper=\popsemi\fontsize{10.6}{14.5}\selectfont\color{popink},
  fontlower=\popsemi\fontsize{10.6}{14.5}\selectfont\color{popink}}}
% Coche dessinée (le glyphe ✓ n'existe pas dans les polices du document)
\newcommand{\popcheck}[1]{\tikz[baseline=-0.5ex]\draw[#1,line width=1.6pt,line cap=round,line join=round](-0.13,0.0)--(-0.04,-0.09)--(0.13,0.1);}
\newcommand{\popboxhead}[3]{\poppill{#1}{#2}{white}\ifx\relax#3\relax\else\hspace{6pt}{\popxbold\fontsize{10.6}{12}\selectfont\color{popink}#3}\fi\par\vspace{7pt}}

\newtcolorbox{aretenir}[1][]{popbase,colframe=popgreen,before upper={\popboxhead{À retenir}{popgreen}{#1}}}
\newtcolorbox{exemplebox}[1][]{popbase,colframe=poporange,before upper={\popboxhead{Exemple}{poporange}{#1}}}
\newtcolorbox{definition}[1][]{popbase,colframe=popblue,before upper={\popboxhead{Définition}{popblue}{#1}}}
\newtcolorbox{propriete}[1][]{popbase,colframe=poppurple,before upper={\popboxhead{Propriété}{poppurple}{#1}}}
\newtcolorbox{methode}[1][]{popbase,colframe=popgold,before upper={\popboxhead{Méthode}{popgold}{#1}}}
\newtcolorbox{attention}[1][]{popbase,colframe=poppink,before upper={\popboxhead{Attention}{poppink}{#1}}}
\newtcolorbox{experience}[1][]{popbase,colframe=popblue,colback=popblueL,before upper={\popboxhead{Expérience}{popblue}{#1}}}
% « Le savais-tu ? » : carte violette pleine (contexte, culture, Cameroun)
\newtcolorbox{savaistu}[1][]{popbase,colback=poppurple,colframe=popink,
  fontupper=\popsemi\fontsize{10.4}{14.2}\selectfont\color{white},
  before upper={\poppill{Le savais-tu\,?}{white}{poppurple}\ifx\relax#1\relax\else\hspace{6pt}{\popxbold\color{white}#1}\fi\par\vspace{7pt}}}
% Exercice résolu : énoncé en haut, \tcblower puis la solution sur fond crème
\newcounter{popexo}\newcounter{popexr}
\newtcolorbox{exoresolu}[1][]{popbase,colframe=poporange,colbacklower=popcream,
  segmentation style={popink,line width=1.2pt,dashed},
  before upper={\stepcounter{popexr}\popboxhead{Exercice résolu \thepopexr}{poporange}{#1}},
  before lower={\poppill{Solution}{popink}{popcream}\par\vspace{6pt}}}
% Exercice (non résolu en place) — la correction est dans la partie Corrigés
\newtcolorbox{exercice}[1][]{popbase,colframe=popink,
  before upper={\stepcounter{popexo}\popboxhead{Exercice \thepopexo}{popink}{#1}}}
% Corrigé
\newtcolorbox{corrige}[1][]{popbase,colframe=popgreen,colback=popgreenL,
  before upper={\popboxhead{Corrigé}{popgreen}{#1}}}
% Objectifs / prérequis / bilan
\newtcolorbox{objectifs}{popbase,colframe=popink,colback=poporangeL,
  before upper={\popboxhead{Objectifs de la leçon}{popink}{}}}
\newtcolorbox{prerequis}{popbase,colframe=popink,colback=popblueL,
  before upper={\popboxhead{Prérequis}{popblue}{}}}
\newtcolorbox{bilan}{popbase,colframe=popink,colback=white,boxrule=2.8pt,
  before upper={\popboxhead{Fiche bilan}{poporange}{L'essentiel à connaître pour l'examen}}}
% Figure encadrée avec légende
\newtcolorbox{popfigure}[1][]{enhanced,breakable=false,colback=white,colframe=popline,boxrule=1.4pt,arc=8pt,
  left=10pt,right=10pt,top=10pt,bottom=8pt,before skip=10pt,after skip=12pt,halign=center,
  lower separated=false,halign lower=center,
  fontlower=\popsemi\fontsize{9}{11.5}\selectfont\color{popmuted},
  #1}
\newcommand{\legende}[1]{\par\vspace{6pt}{\popsemi\fontsize{9}{11.5}\selectfont\color{popmuted}#1}}

% ============================================================
% Signature OnBuch+
% ============================================================
\newcommand{\poplogoinline}[1]{{\poplogo\fontsize{#1}{#1}\selectfont\color{popink}OnBuch\color{poporange}+}}
\newcommand{\signatureonbuch}{%
  \begin{tcolorbox}[enhanced,colback=popink,colframe=popink,boxrule=2.2pt,arc=10pt,
      drop shadow=poporange,left=14pt,right=14pt,top=10pt,bottom=10pt,before skip=0pt,after skip=0pt]
    \tikz[baseline=-0.7ex]\node[circle,fill=popgreen,draw=popcream,line width=1.3pt,minimum size=22pt,inner sep=0]{\popcheck{white}};%
    \hspace{9pt}\begin{minipage}[c]{0.62\linewidth}
      {\popxbold\fontsize{12}{13}\selectfont\color{popcream}Signé\ \textbullet\ {\poplogo OnBuch\color{poporange}+}}\par\vspace{2pt}
      {\raggedright\popsemi\fontsize{8}{10.5}\selectfont\color{popline}\MakeUppercase{Équipe pédagogique OnBuch+}\\\MakeUppercase{Conforme au programme officiel MINESEC}\par}
    \end{minipage}\hfill
    {\poplogo\fontsize{22}{22}\selectfont\color{popcream}OnBuch\color{poporange}+}
  \end{tcolorbox}%
}

% ============================================================
% Page de garde
% #1 titre · #2 matière · #3 niveau · #4 module · #5 n° leçon · #6 durée ·
% #7 description / objectifs (texte ou itemize)
% ============================================================
\newcommand{\pagegarde}[7]{%
  \thispagestyle{empty}
  \newgeometry{a4paper,margin=1.7cm,top=1.6cm,bottom=1.4cm}%
  \begin{tikzpicture}[remember picture,overlay]
    \fill[poporange!18] ([xshift=-3.2cm,yshift=-3.6cm]current page.north east) circle (4.6cm);
    \fill[poppurple!14] ([xshift=2.4cm,yshift=7.6cm]current page.south west) circle (3.8cm);
  \end{tikzpicture}%
  {\poplogo\fontsize{30}{30}\selectfont\color{popink}OnBuch\color{poporange}+}\hfill
  \raisebox{0.6ex}{\poppill{Cours signé \textbullet\ Premium}{popink}{popcream}}\par
  \vspace{1pt}\popsquiggle{poppurple}\par
  \vspace{8pt}
  \begin{tcolorbox}[enhanced,colback=poporange,colframe=popink,boxrule=2.8pt,arc=13pt,
      drop shadow=popink,left=18pt,right=18pt,top=12pt,bottom=13pt,after skip=10pt]
    \poppill{#2 \textbullet\ #3}{popink}{popcream}\hspace{5pt}\poppill{#4}{white}{popink}\par\vspace{8pt}
    {\popdisplay\fontsize{14}{14}\selectfont\color{popink}LEÇON #5}\par\vspace{4pt}
    {\raggedright\hyphenpenalty=10000\exhyphenpenalty=10000\popdisplay\fontsize{25}{27}\selectfont\color{popcream}\MakeUppercase{#1}\par}
    \vspace{8pt}
    {\popxbold\fontsize{9.5}{9.5}\selectfont\color{popink}\MakeUppercase{Durée conseillée : #6}}
  \end{tcolorbox}
  \begin{tcolorbox}[enhanced,colback=white,colframe=popink,boxrule=2.2pt,arc=10pt,
      drop shadow=popink,left=16pt,right=16pt,top=10pt,bottom=10pt,after skip=8pt]
    \poppill{Dans cette leçon}{poppurple}{white}\par\vspace{5pt}
    \popsemi\fontsize{9.3}{12.6}\selectfont\color{popink}#7
  \end{tcolorbox}
  \noindent\begin{minipage}[t]{0.487\linewidth}
    \begin{tcolorbox}[enhanced,colback=popgreen,colframe=popink,boxrule=2.2pt,arc=10pt,
        drop shadow=popink,left=11pt,right=11pt,top=10pt,bottom=10pt,height=3.1cm,valign=center]
      \tcbox[colback=white,colframe=popink,boxrule=1.3pt,arc=5pt,boxsep=0pt,nobeforeafter,
        left=3pt,right=3pt,top=3pt,bottom=3pt,valign=center]{\qrcode[height=1.9cm]{https://play.google.com/store/apps/details?id=com.luvvix.onbuch&hl=fr}}%
      \hspace{8pt}\begin{minipage}[c]{0.5\linewidth}
        {\popdisplay\fontsize{11}{12.5}\selectfont\color{white}\MakeUppercase{Télécharge OnBuch}}\par\vspace{4pt}
        {\raggedright\popsemi\fontsize{8.4}{11}\selectfont\color{white}Gratuit \textbullet\ Google Play\\Tuteur IA, annales, concours, résultats\par}
      \end{minipage}
    \end{tcolorbox}
  \end{minipage}\hfill
  \begin{minipage}[t]{0.487\linewidth}
    \begin{tcolorbox}[enhanced,colback=poppurple,colframe=popink,boxrule=2.2pt,arc=10pt,
        drop shadow=popink,left=11pt,right=11pt,top=10pt,bottom=10pt,height=3.1cm,valign=center]
      \tikz[baseline=-0.7ex]\node[circle,fill=white,draw=popink,line width=1.3pt,minimum size=1.6cm,inner sep=0]{\popdisplay\fontsize{24}{24}\selectfont\color{poppurple}+};%
      \hspace{8pt}\begin{minipage}[c]{0.6\linewidth}
        {\popdisplay\fontsize{11}{12.5}\selectfont\color{white}\MakeUppercase{Passe à OnBuch+}}\par\vspace{4pt}
        {\raggedright\popsemi\fontsize{8.4}{11}\selectfont\color{white}Tous les cours signés, Tuteur IA illimité, fiches PDF — onglet OnBuch+ de l'app\par}
      \end{minipage}
    \end{tcolorbox}
  \end{minipage}
  \vspace{8pt}
  \popcontenu
  \vfill
  \signatureonbuch
  \restoregeometry
  \newpage
}

% Rangée « Ce cours contient » (page de garde)
\newcommand{\poptile}[3]{% #1 couleur #2 gros texte #3 légende
  \begin{tcolorbox}[enhanced,colback=#1,colframe=popink,boxrule=2pt,arc=9pt,shadow={2.5pt}{-2.5pt}{0pt}{popink},
      left=6pt,right=6pt,top=9pt,bottom=9pt,halign=center,valign=center,height=1.75cm,width=\linewidth,nobeforeafter]
    {\popdisplay\fontsize{12.5}{13}\selectfont\color{white}\MakeUppercase{#2}}\par\vspace{3pt}
    {\centering\popsemi\fontsize{7.8}{9.5}\selectfont\color{white}#3\par}
  \end{tcolorbox}}
\newcommand{\popcontenu}{%
  \par\noindent{\popxboldls\fontsize{7.5}{7.5}\selectfont\color{popmuted}CE COURS SIGNÉ CONTIENT}\par\vspace{7pt}
  \noindent\begin{minipage}[t]{0.235\linewidth}\poptile{popblue}{Cours}{complet, illustré, pas à pas}\end{minipage}\hfill
  \begin{minipage}[t]{0.235\linewidth}\poptile{poporange}{Méthodes}{et exercices résolus}\end{minipage}\hfill
  \begin{minipage}[t]{0.235\linewidth}\poptile{poppink}{Exercices}{type examen + corrigés}\end{minipage}\hfill
  \begin{minipage}[t]{0.235\linewidth}\poptile{popgreen}{Fiche bilan}{l'essentiel à retenir}\end{minipage}\par}

% Sommaire
\newtcolorbox{sommaire}{enhanced,colback=white,colframe=popink,boxrule=2.2pt,arc=10pt,
  drop shadow=popink,left=18pt,right=18pt,top=13pt,bottom=13pt,before skip=6pt,after skip=20pt,
  before upper={\poppill{Sommaire}{poppurple}{white}\par\vspace{9pt}\popsemi\fontsize{10.6}{14.5}\selectfont\color{popink}}}

% Signature de fin de cours — poussée tout en bas de la dernière page
\newcommand{\findecours}{%
  \par\vfill\nopagebreak
  \begin{center}\popsquiggle{poporange}\end{center}\vspace{4pt}
  \signatureonbuch
}
\AtBeginDocument{\def\llbracket{\mathopen{[\mkern-3.5mu[}}\def\rrbracket{\mathclose{]\mkern-3.5mu]}}} % intervalles d'entiers (glyphe ⟦ absent de la police)
% Glyphes absents de Fira Math : redéfinis à partir de glyphes disponibles
\AtBeginDocument{%
  \def\setminus{\mathbin{\backslash}}\let\smallsetminus\setminus
  \def\vdots{\mathord{\vbox{\baselineskip3.2pt\lineskiplimit0pt\kern4pt\hbox{.}\hbox{.}\hbox{.}}}}%
  \def\ddots{\mathinner{\mkern1mu\raise6.5pt\hbox{.}\mkern2mu\raise3.6pt\hbox{.}\mkern2mu\raise0.7pt\hbox{.}\mkern1mu}}%
  \def\triangle{\mathord{\scalebox{0.9}{$\symup{\Delta}$}}}%
  \def\nabla{\mathord{\raisebox{\depth}{\scalebox{1}[-1]{$\symup{\Delta}$}}}}%
  \def\checkmark{\popcheck{popdark}}%
  \def\star{\mathbin{\ast}}\def\bigstar{\mathord{\ast}}%
  \def\diamond{\mathbin{\scalebox{0.6}{$\Diamond$}}}%
  \def\bigcup{\mathop{\mathchoice{\raisebox{-0.3em}{\scalebox{1.7}{$\cup$}}}{\scalebox{1.25}{$\cup$}}{\cup}{\cup}}\displaylimits}%
  \def\bigcap{\mathop{\mathchoice{\raisebox{-0.3em}{\scalebox{1.7}{$\cap$}}}{\scalebox{1.25}{$\cap$}}{\cap}{\cap}}\displaylimits}%
  \def\lesseqqgtr{\mathrel{\lessgtr}}\def\bigoplus{\mathop{\oplus}\displaylimits}%
}
\providecommand{\ssi}{si et seulement si} % abréviation fréquente
\AtBeginDocument{\providecommand{\wideparen}{\overparen}} % arc de cercle

%% FIGFIX-BEGIN v5 — correctifs globaux des figures (docs/onbuch-generation/tools/figfix)
\usepackage{adjustbox}\usepackage{etoolbox}\tcbuselibrary{raster}
% toute figure TikZ/pgfplots plus large que la ligne est réduite à la largeur disponible
\BeforeBeginEnvironment{tikzpicture}{\begin{adjustbox}{max width=\linewidth}}
\AfterEndEnvironment{tikzpicture}{\end{adjustbox}}
% légendes pgfplots lisibles par-dessus les courbes
\pgfplotsset{every axis/.append style={legend style={fill=white,fill opacity=0.92,text opacity=1,draw=black!15,inner sep=3pt}}}
% étiquettes d'arêtes (syntaxe edge["..."]) avec fond blanc
\tikzset{every edge quotes/.append style={fill=white,inner sep=1.2pt}}
%% FIGFIX-END
\begin{document}

\pagegarde{Calcul littéral}{Mathématiques}{3ème}{Mathématiques – classe de 3ème}{N05}{5 séances (fiche de progression harmonisée)}{Cette leçon consolide et approfondit le calcul littéral : manipulation des monômes et polynômes, développement, réduction et factorisation d'expressions, jusqu'aux fractions rationnelles. Elle fournit les outils algébriques indispensables à la résolution d'équations, à la modélisation de situations concrètes et à la réussite au BEPC.
\vspace{4pt}
\begin{multicols}{2}\small
\begin{itemize}
\item À la fin de cette leçon, je sais reconnaître et évaluer une expression littérale pour des valeurs données des variables
\item À la fin de cette leçon, je sais identifier un monôme, déterminer son coefficient et son degré, et additionner ou multiplier des monômes
\item À la fin de cette leçon, je sais réduire et ordonner un polynôme, et déterminer son degré
\item À la fin de cette leçon, je sais développer un produit d'expressions en utilisant la distributivité simple et double
\item À la fin de cette leçon, je sais utiliser les identités remarquables pour développer ou factoriser rapidement
\item À la fin de cette leçon, je sais factoriser une expression littérale par facteur commun ou par identité remarquable
\end{itemize}\end{multicols}}

\begin{sommaire}
\begin{multicols}{2}
\begin{enumerate}
\item Expressions littérales
\item Monômes et polynômes
\item Développement et réduction
\item Factorisation
\item Fractions rationnelles
\item Synthèse et applications à la vie courante
\item Activité d'intégration
\item Méthodes et exercices résolus
\item Exercices
\item Corrigés des exercices
\item Fiche bilan
\end{enumerate}
\end{multicols}
\end{sommaire}

\begin{objectifs}
\begin{itemize}
\item À la fin de cette leçon, je sais reconnaître et évaluer une expression littérale pour des valeurs données des variables
\item À la fin de cette leçon, je sais identifier un monôme, déterminer son coefficient et son degré, et additionner ou multiplier des monômes
\item À la fin de cette leçon, je sais réduire et ordonner un polynôme, et déterminer son degré
\item À la fin de cette leçon, je sais développer un produit d'expressions en utilisant la distributivité simple et double
\item À la fin de cette leçon, je sais utiliser les identités remarquables pour développer ou factoriser rapidement
\item À la fin de cette leçon, je sais factoriser une expression littérale par facteur commun ou par identité remarquable
\item À la fin de cette leçon, je sais simplifier une fraction rationnelle après avoir déterminé ses valeurs interdites
\item À la fin de cette leçon, je sais modéliser une situation concrète de la vie courante par une expression littérale et justifier ma démarche par écrit
\end{itemize}
\end{objectifs}

\begin{prerequis}
\begin{itemize}
\item Calcul avec les nombres relatifs et les fractions (4ème)
\item Notion d'expression littérale et substitution d'une valeur à une lettre (4ème)
\item Distributivité simple : k(a + b) = ka + kb (4ème)
\item Puissances d'un nombre et règles de calcul sur les exposants (4ème)
\item Résolution d'équations simples du premier degré (4ème)
\end{itemize}
\end{prerequis}

\newpage

\coursec{Expressions littérales}

Imagine un instant le bruit vibrant du marché Mokolo : les appels des vendeurs, l'odeur du ndolé qui mijote, les tissus colorés qui volent d'un étal à l'autre. Au milieu de cette effervescence, \textbf{Mama Ngo Bell} vend des pagnes. Elle achète chaque pagne à \cle{x} francs (son prix d'achat varie selon les arrivages) et le revend avec un bénéfice fixe de 500 FCFA. Pour fidéliser ses clientes, elle offre une remise de 10 \% à toute personne qui en prend au moins trois.

Chaque matin, elle pourrait refaire ses calculs pour chaque cliente : « Si j'achète à 2 500 FCFA, je vends à 3 000 FCFA, et si la cliente en prend trois, ça fait… ». Mais Mama Ngo Bell est maligne. Elle écrit \textbf{une seule formule} avec la lettre \cle{x}. Cette formule, c'est une \emph{expression littérale}. Elle lui permet de connaître instantanément sa recette pour \emph{n'importe quel} prix d'achat, sans recommencer le calcul.

Mais comment écrire cette formule proprement ? Comment la simplifier pour qu'elle soit utilisable en un clin d'œil ? Et si on connaît la recette finale, comment retrouver le prix d'achat \cle{x} ? C'est tout l'art du \textbf{calcul littéral} : calculer avec des lettres pour résoudre d'un coup une infinité de problèmes concrets.

\courssub{Définition, vocabulaire et conventions d'écriture}

\begin{definition}[Expression littérale, variable, inconnue, paramètre]
Une \textbf{expression littérale} est une écriture mathématique qui mélange des \textbf{nombres} (constantes) et des \textbf{lettres}, liées par des signes d'opérations ($+$, $-$, $\times$, $\div$, puissances).
\begin{itemize}
 \item La lettre représente un \textbf{nombre} : on l'appelle \textbf{variable} (elle peut prendre plusieurs valeurs), \textbf{inconnue} (on cherche sa valeur dans une équation) ou \textbf{paramètre} (elle est fixe dans le problème mais peut changer d'un problème à l'autre).
 \item Exemples : $3x + 5$, $x^2 - 2xy + y^2$, $\dfrac{P}{2} - L$, $2\pi r$.
\end{itemize}
\end{definition}

\begin{propriete}[Conventions d'écriture (simplification de l'écriture)]
Pour alléger les expressions, on adopte les règles universelles suivantes :
\begin{enumerate}
 \item \textbf{Suppression du signe $\times$} devant une lettre ou une parenthèse : $3 \times x$ s'écrit \cle{3x} ; $a \times (b+c)$ s'écrit \cle{a(b+c)}.
 \item \textbf{Ordre alphabétique} pour le produit de lettres : $x \times y$ s'écrit \cle{xy} (et non $yx$).
 \item \textbf{Exposant} pour les produits identiques : $x \times x$ s'écrit \cle{$x^2$} (se lit « x carré ») ; $x \times x \times x$ s'écrit \cle{$x^3$} (« x cube »).
 \item \textbf{Coefficient numérique} : dans $5x^2y$, le nombre $5$ est le \cle{coefficient}, $x^2y$ est la \cle{partie littérale}.
 \item \textbf{Priorité des opérations} : puissances $\rightarrow$ multiplications/divisions $\rightarrow$ additions/soustractions. Les parenthèses priment sur tout.
\end{enumerate}
\end{propriete}

\begin{exemplebox}[Écriture simplifiée]
\begin{itemize}
 \item $7 \times a \times b$ s'écrit \cle{7ab}.
 \item $x \times x \times 3 \times y$ s'écrit \cle{$3x^2y$} (coefficient 3, puis $x^2$, puis $y$).
 \item $(a+b) \times (a-b)$ s'écrit \cle{(a+b)(a-b)}.
 \item $\dfrac{1}{2} \times h \times (B+b)$ s'écrit \cle{$\dfrac{h(B+b)}{2}$} ou \cle{$\frac{1}{2}h(B+b)$}.
\end{itemize}
\end{exemplebox}

\begin{savaistu}[L'origine du mot « algèbre »]
Le mot \textbf{algèbre} vient de l'arabe \emph{al-jabr} , qui signifie « restauration » ou « réduction ». Il est tiré du titre de l'ouvrage fondateur \emph{Al-Kitāb al-mukhtaṣar fī ḥisāb al-jabr wal-muqābala} (\emph{Le livre compendieux sur le calcul par la restauration et la comparaison}), écrit vers 820 par le grand mathématicien perse \textbf{Al-Khawarizmi} (d'où vient aussi le mot « algorithme »). « Al-jabr » désignait l'opération qui consiste à passer un terme soustrait de l'autre côté de l'égalité en le faisant devenir ajouté (ex: $x - 5 = 10 \rightarrow x = 10 + 5$). C'est la naissance du calcul littéral tel que nous le connaissons !
\end{savaistu}

\courssub{Valeur numérique d'une expression littérale}

Une expression littérale n'est qu'une « recette de calcul » en attente. Pour obtenir un nombre, il faut \textbf{remplacer la (les) lettre(s) par des valeurs numériques précises}. C'est le calcul de la \textbf{valeur numérique}.

\begin{methode}[Calculer la valeur numérique d'une expression (Substitution)]
Pour calculer la valeur d'une expression pour une valeur donnée de la variable :
\begin{enumerate}
 \item \textbf{Écrire l'expression} en remplaçant \textbf{systématiquement} la variable par la valeur numérique \textbf{entre parenthèses}. (C'est la règle d'or pour éviter les erreurs de signe !)
 \item \textbf{Calculer} en respectant scrupuleusement l'ordre des priorités (puissances, puis multiplications/divisions, puis additions/soustractions).
 \item \textbf{Donner le résultat final} (entier, décimal ou fraction simplifiée).
\end{enumerate}
\end{methode}

\begin{attention}[Le piège du signe moins sans parenthèses --- Erreur n°1 au BEPC !]
Soit l'expression $A = -x^2$. Calculons sa valeur pour $x = -3$.
\begin{itemize}
 \item \textbf{Faux (sans parenthèses) :} $-(-3)^2$ ? Non, l'élève écrit souvent $- -3^2 = -9$ ou pire $9$.
 \item \textbf{Vrai (avec parenthèses) :} $A = -(x)^2$. Pour $x=-3$ : $A = -(-3)^2 = -(9) = \mathbf{-9}$.
 \item \textbf{Autre piège :} $B = (-x)^2$. Pour $x=-3$ : $B = (-(-3))^2 = (3)^2 = \mathbf{9}$.
\end{itemize}
\textbf{Règle d'or :} \emph{Jamais de substitution sans parenthèses.} Écrivez toujours $(\text{valeur})$ à la place de la lettre.
\end{attention}

\begin{exoresolu}[Calcul de valeurs numériques]
Évaluer les expressions suivantes pour $x = -2$ :
\begin{enumerate}
 \item $A = 2x^2 - 3x + 1$
 \item $B = (2x - 3)^2$
 \item $C = \dfrac{x+5}{x-1}$
\end{enumerate}

\tcblower
\textbf{Solution détaillée :}

\textbf{1. Expression $A$ :}
On remplace $x$ par $(-2)$ \textbf{entre parenthèses} :
\[ A = 2(-2)^2 - 3(-2) + 1 \]
Priorité aux puissances : $(-2)^2 = 4$.
\[ A = 2 \times 4 - 3 \times (-2) + 1 \]
Multiplications : $2 \times 4 = 8$ ; $-3 \times (-2) = +6$.
\[ A = 8 + 6 + 1 = \mathbf{15} \]

\textbf{2. Expression $B$ :}
\[ B = (2(-2) - 3)^2 = (-4 - 3)^2 = (-7)^2 = \mathbf{49} \]
\emph{Attention :} $(2x-3)^2 \neq 4x^2 - 9$ ! On calcule la parenthèse \textbf{avant} le carré.

\textbf{3. Expression $C$ :}
\[ C = \dfrac{(-2)+5}{(-2)-1} = \dfrac{3}{-3} = \mathbf{-1} \]
\end{exoresolu}

\begin{popfigure}
\centering
\begin{tabularx}{\linewidth}{|c|*{4}{>{\centering\arraybackslash}X|}}\hline
\rowcolor{popblueL}
\textbf{Valeur de $x$} & \textbf{-2} & \textbf{0} & \textbf{1} & \textbf{3} \\ \hline
\textbf{Expression $A = 2x^2 - 3x + 1$} & 
$2(-2)^2 - 3(-2) + 1$ \\ $= 8 + 6 + 1 = \mathbf{15}$ & 
$2(0)^2 - 3(0) + 1$ \\ $= 0 - 0 + 1 = \mathbf{1}$ & 
$2(1)^2 - 3(1) + 1$ \\ $= 2 - 3 + 1 = \mathbf{0}$ & 
$2(3)^2 - 3(3) + 1$ \\ $= 18 - 9 + 1 = \mathbf{10}$ \\ \hline
\textbf{Expression $B = (2x-3)^2$} & 
$(2(-2)-3)^2$ \\ $= (-7)^2 = \mathbf{49}$ & 
$(2(0)-3)^2$ \\ $= (-3)^2 = \mathbf{9}$ & 
$(2(1)-3)^2$ \\ $= (-1)^2 = \mathbf{1}$ & 
$(2(3)-3)^2$ \\ $= (3)^2 = \mathbf{9}$ \\ \hline
\end{tabularx}
\legende{Tableau de valeurs des expressions $A$ et $B$ pour différentes valeurs de $x$. Remarquez l'importance des parenthèses pour $x=-2$.}
\end{popfigure}

\courssub{Traduire une situation concrète en expression littérale}

C'est l'étape cruciale de la \textbf{modélisation} : passer du français au langage mathématique. On suit généralement trois étapes :
\begin{enumerate}
 \item \textbf{Choisir la (les) variable(s)} : quelle grandeur varie ? Quelle lettre la représente ? (Souvent $x$, $y$, $t$, $p$, $n$...).
 \item \textbf{Exprimer les grandeurs intermédiaires} en fonction de la variable.
 \item \textbf{Écrire l'expression finale} demandée (prix total, périmètre, aire, recette, bénéfice...).
\end{enumerate}

\begin{exemplebox}[Modélisation classique : Géométrie]
Un rectangle a pour dimensions : \textbf{Longueur} $L = x + 5$ (en mètres) et \textbf{Largeur} $\ell = x$ (en mètres).
\begin{itemize}
 \item \textbf{Périmètre} $P = 2 \times (L + \ell) = 2 \times ((x+5) + x) = 2(2x + 5) = \mathbf{4x + 10}$ mètres.
 \item \textbf{Aire} $\mathcal{A} = L \times \ell = (x+5) \times x = \mathbf{x(x+5)}$ ou $\mathbf{x^2 + 5x}$ m².
\end{itemize}
\end{exemplebox}

\begin{popfigure}
\begin{tikzpicture}[scale=0.7, every node/.style={font=\small}]
 % Rectangle
 \draw[popblue, thick] (0,0) rectangle (6,3.5);
 % Dimensions
 \draw[<->, poporange, thick] (0,-0.5) -- (6,-0.5) node[midway, below=2pt] {$x+5$ m};
 \draw[<->, poporange, thick] (-0.5,0) -- (-0.5,3.5) node[midway, left=2pt] {$x$ m};
 % Points
 \foreach \pt in {(0,0),(6,0),(6,3.5),(0,3.5)} \fill[popdark] \pt circle (2pt);
 % Formules inside
 \node[align=center, text=popdark, fill=popcream, rounded corners, inner sep=4pt] at (3,1.75) {
 $\mathcal{A} = x(x+5) = x^2+5x$ \\ $P = 2(x+5+x) = 4x+10$
 };
\end{tikzpicture}
\legende{Parcelle rectangulaire de dimensions $(x+5)$ m et $x$ m. Modélisation littérale du périmètre et de l'aire.}
\end{popfigure}

\courssub{Le cas de Mama Ngo Bell : Modélisation complète}

Revenons à notre situation d'accroche pour construire l'expression de la recette.

\begin{exoresolu}[Modélisation : La recette de Mama Ngo Bell]
\textbf{Énoncé :} Mama Ngo Bell achète un pagne à $x$ FCFA. Elle le revend avec un bénéfice de 500 FCFA. Si une cliente achète au moins 3 pagnes, elle bénéficie d'une remise de 10 \% sur le total.
\begin{enumerate}
 \item Exprimer le \textbf{prix de vente unitaire} $PV$ en fonction de $x$.
 \item Exprimer le \textbf{prix total} $PT_3$ pour 3 pagnes \textbf{sans remise}.
 \item Exprimer le \textbf{montant de la remise} $R$ (10 \% de $PT_3$).
 \item Exprimer le \textbf{prix final} $PF$ que paie la cliente pour 3 pagnes \textbf{avec remise}.
 \item Simplifier l'expression de $PF$.
 \item Calculer $PF$ si $x = 2\,500$ FCFA.
\end{enumerate}

\tcblower
\textbf{Solution :}

\begin{enumerate}
 \item \textbf{Prix de vente unitaire :} Prix d'achat + Bénéfice.
 \[ PV = x + 500 \]
 \item \textbf{Prix total pour 3 pagnes (sans remise) :}
 \[ PT_3 = 3 \times PV = 3(x + 500) \]
 \item \textbf{Montant de la remise (10 \%) :} 10 \% = $\frac{10}{100} = {0,1}$.
 \[ R = {0,1} \times PT_3 = {0,1} \times 3(x + 500) = {0,3}(x + 500) \]
 \item \textbf{Prix final avec remise :} Total $-$ Remise.
 \[ PF = PT_3 - R = 3(x + 500) - {0,3}(x + 500) \]
 \item \textbf{Simplification (factorisation anticipée) :} On met $3(x+500)$ en facteur.
 \[ PF = (3 - {0,3})(x + 500) = {2,7}(x + 500) \]
 On peut aussi développer : $PF = {2,7}x + 1\,350$.
 \item \textbf{Calcul numérique pour $x = 2\,500$ :}
 \[ PF = {2,7} \times (2\,500 + 500) = {2,7} \times 3\,000 = \mathbf{8\,100 \text{ FCFA}} \]
 \emph{Vérification :} $PV = 3\,000$. $3 \times 3\,000 = 9\,000$. Remise $= 900$. $9\,000 - 900 = 8\,100$. \checkmark
\end{enumerate}
\end{exoresolu}

\courssub{Égalité de deux expressions littérales : Vérifier $\neq$ Prouver}

Deux expressions littérales $E_1(x)$ et $E_2(x)$ sont \textbf{égales} si elles donnent la \textbf{même valeur numérique} pour \textbf{toutes} les valeurs possibles de $x$ (dans leur ensemble de définition).

\begin{propriete}[Test de valeurs vs Preuve]
\begin{itemize}
 \item \textbf{Tester des valeurs} (ex: $x = 0, 1, -1, 2$) permet de \textbf{se convaincre} ou de \textbf{trouver une erreur} (si les valeurs diffèrent, les expressions sont \textbf{différentes}).
 \item \textbf{Un seul contre-exemple suffit pour affirmer que deux expressions ne sont pas égales.}
 \item \textbf{Mais tester 10, 100, 1000 valeurs ne prouve JAMAIS l'égalité.} Pour \textbf{prouver} l'égalité, il faut \textbf{transformer l'une en l'autre par des opérations algébriques valides} (développement, réduction, factorisation, identités remarquables) — ce que nous ferons dans les sections suivantes.
\end{itemize}
\end{propriete}

\begin{exemplebox}[Égal ou pas égal ?]
Soit $E_1 = (x+2)^2$ et $E_2 = x^2 + 4$.
\begin{itemize}
 \item Test $x=0$ : $E_1 = 4$, $E_2 = 4$. \emph{(Égal)}
 \item Test $x=1$ : $E_1 = 9$, $E_2 = 5$. \emph{(Différent !)}
\end{itemize}
\cle{Contre-exemple trouvé pour $x=1$}. Les expressions ne sont \textbf{pas égales}.
\emph{Note :} $E_1 = x^2 + 4x + 4 \neq x^2 + 4$.
\end{exemplebox}

\courssub{Autres applications camerounaises}

Le calcul littéral n'est pas qu'un exercice scolaire : c'est l'outil quotidien du commerçant, de l'artisan, de l'ingénieur, du développeur d'applications (MTN MoMo, Orange Money), du géomètre-topographe.

\begin{exercice}[Applications variées — À vous de jouer !]
\begin{enumerate}
 \item \textbf{Transport (Taxi) :} Un taxi prend 200 FCFA pour la prise en charge, puis 100 FCFA par kilomètre. Écrire l'expression du prix $P$ d'une course de $x$ km. Calculer pour 12 km.
 \item \textbf{Agriculture :} Un champ carré de côté $x$ mètres est agrandi : on augmente la longueur de 10 m et la largeur de 5 m. Exprimer la nouvelle aire $\mathcal{A}_{new}$ en fonction de $x$. Développer.
 \item \textbf{Télécom (Forfait Data) :} Un forfait coûte 1 000 FCFA pour 1 Go. Au-delà, chaque Mo supplémentaire coûte 5 FCFA. Si $x$ représente le nombre de Mo consommés \textbf{au-delà} du Go inclus, écrire le coût total $C$. Calculer pour 350 Mo supplémentaires.
 \item \textbf{Électricité (EDC/ENEO) :} La facture mensuelle comprend un abonnement fixe de 500 FCFA et un coût au kWh de 50 FCFA. Si $x$ est le nombre de kWh consommés, écrire la facture $F$. Si la facture est de 12 500 FCFA, quelle expression permet de trouver $x$ ? (On ne résout pas, on écrit l'équation).
\end{enumerate}
\end{exercice}

\begin{corrige}[Exercice n°1 -- Transport]
$P(x) = 200 + 100x$. Pour $x=12$ : $P = 200 + 1200 = 1\,400$ FCFA.
\end{corrige}
\begin{corrige}[Exercice n°2 -- Agriculture]
Ancienne aire $= x^2$. Nouvelles dimensions : $(x+10)$ et $(x+5)$.
$\mathcal{A}_{new} = (x+10)(x+5) = x^2 + 15x + 50$ m².
\end{corrige}
\begin{corrige}[Exercice n°3 -- Télécom]
$C(x) = 1\,000 + 5x$. Pour $x=350$ : $C = 1\,000 + 1\,750 = 2\,750$ FCFA.
\end{corrige}
\begin{corrige}[Exercice n°4 -- Électricité]
$F(x) = 500 + 50x$. Si $F = 12\,500$, l'équation est $500 + 50x = 12\,500$.
\end{corrige}

\begin{aretenir}[Bilan : Ce qu'il faut maîtriser avant la suite]
\begin{itemize}
 \item Une \textbf{expression littérale} mélange nombres et lettres. La lettre est une \textbf{variable}.
 \item \textbf{Conventions :} $3x$, $x^2$, $xy$, parenthèses pour produit.
 \item \textbf{Valeur numérique :} Substitution \textbf{entre parenthèses} + respect des priorités.
 \item \textbf{Modélisation :} Variable $\rightarrow$ Grandeurs intermédiaires $\rightarrow$ Expression finale.
 \item \textbf{Égalité :} Un contre-exemple suffit pour dire « non ». La preuve algébrique (développement/réduction/factorisation) est nécessaire pour dire « oui ».
\end{itemize}
\end{aretenir}

\coursec{Monômes et polynômes}

Après avoir découvert les \emph{expressions littérales} — ces écritures qui mêlent nombres et lettres pour modéliser des situations concrètes (prix d'un panier de fruits, périmètre d'un champ, facture d'électricité CAMTEL) — nous allons maintenant les \emph{structurer}. Toute expression littérale peut être décomposée en \textbf{blocs élémentaires} appelés \textbf{monômes}. Comprendre ces blocs, savoir les manipuler et les assembler en \textbf{polynômes} est la clé pour développer, factoriser, résoudre des équations et, demain, étudier des fonctions.

\courssub{Les monômes : définition, coefficient, partie littérale et degré}

\begin{definition}[Monôme]
Un \textbf{monôme} est une expression littérale qui ne contient \textbf{que des multiplications et des divisions par des nombres non nuls}, ainsi que des puissances d'entiers naturels des variables. Il s'écrit sous la forme canonique :
\[
k \cdot x_1^{\alpha_1} x_2^{\alpha_2} \cdots x_n^{\alpha_n}
\]
où :
\begin{itemize}
    \item $k \in \mathbb{R}$ (souvent $\mathbb{Q}$ ou $\mathbb{Z}$) est le \textbf{coefficient} (ou coefficient numérique) ;
    \item $x_1^{\alpha_1} x_2^{\alpha_2} \cdots x_n^{\alpha_n}$ est la \textbf{partie littérale} ;
    \item $\alpha_i \in \mathbb{N}$ sont les \textbf{exposants} des variables.
\end{itemize}
\end{definition}

\begin{exemplebox}[Exemples et contre-exemples]
\begin{center}
\begin{tabularx}{\linewidth}{|l|X|c|}\hline
\textbf{Expression} & \textbf{Analyse} & \textbf{Monôme ?} \\ \hline
$7x^2y$ & $k=7$, partie littérale $x^2y$, exposants $2,1$ & \textbf{Oui} \\ \hline
$-3a^5$ & $k=-3$, partie littérale $a^5$ & \textbf{Oui} \\ \hline
$\dfrac{2}{3}uv^2$ & $k=\frac{2}{3}$, partie littérale $uv^2$ & \textbf{Oui} \\ \hline
$x + 2$ & Contient une \textbf{addition} & \textbf{Non} (binôme) \\ \hline
$\dfrac{5}{x}$ & Division \textbf{par une variable} ($x^{-1}$) & \textbf{Non} (fraction rationnelle) \\ \hline
$\sqrt{x}$ & Exposant $\frac{1}{2} \notin \mathbb{N}$ & \textbf{Non} \\ \hline
$4$ & $k=4$, partie littérale $1$ (degré $0$) & \textbf{Oui} (monôme constant) \\ \hline
\end{tabularx}
\end{center}
\end{exemplebox}

\begin{propriete}[Degré d'un monôme]
Le \textbf{degré} d'un monôme à une variable $x$ est l'exposant de $x$.
Le \textbf{degré} d'un monôme à plusieurs variables est la \textbf{somme des exposants} de toutes ses variables.
Le monôme non nul constant (ex. $7$) a pour degré $0$. Le monôme nul $0$ n'a pas de degré défini.
\end{propriete}

\begin{popfigure}
\centering
\begin{tikzpicture}[
    level distance=1.8cm,
    sibling distance=2.5cm,
    every node/.style={font=\small, align=center},
    edge from parent/.style={draw=popdark, thick, ->},
    monome/.style={rectangle, draw=popblue, thick, fill=popblueL, rounded corners=4pt, minimum width=2.2cm},
    detail/.style={rectangle, draw=poporange, fill=poporangeL, rounded corners=3pt, minimum width=2.5cm},
    deg/.style={rectangle, draw=popgreen, fill=popgreenL, rounded corners=3pt, minimum width=2cm},
]
\node[monome] (M) {$5x^2y^3$}
    child { node[detail] (C) {Coefficient : $5$} }
    child { node[detail] (PL) {Partie littérale : $x^2y^3$}
        child { node[detail] (X) {Variable $x$ : exposant $2$} }
        child { node[detail] (Y) {Variable $y$ : exposant $3$} }
        child { node[deg] (D) {Degré = $2 + 3 = 5$} }
    };
\node[draw=poppurple, fill=poppurpleL, rounded corners=4pt, align=center, font=\small\bfseries, below=1.2cm of M] {
    Règle : \textdegree(monôme) = somme des exposants des variables
};
\end{tikzpicture}
\legende{Lecture structurée du degré d'un monôme à deux variables.}
\end{popfigure}

\begin{attention}[Piège classique : coefficient « invisible »]
L'expression $x^3$ est un monôme dont le coefficient est \textbf{$1$} (et non $0$). De même, $-xy^2$ a pour coefficient \textbf{$-1$}. Ne jamais oublier le signe $-$ qui fait partie du coefficient !
\end{attention}

\courssub{Opérations sur les monômes}

\courssub{Somme et différence : les monômes semblables}

\begin{definition}[Monômes semblables]
Deux monômes sont \textbf{semblables} s'ils ont \textbf{exactement la même partie littérale} (mêmes variables aux mêmes exposants). Seuls leurs coefficients diffèrent.
\end{definition}

\begin{propriete}[Réduction de monômes semblables]
La somme (ou la différence) de deux monômes semblables est un monôme \textbf{semblable} dont le coefficient est la somme (ou la différence) des coefficients.
\[
ax^\alpha y^\beta + bx^\alpha y^\beta = (a+b)x^\alpha y^\beta
\]
Si les monômes ne sont \textbf{pas} semblables, la somme \textbf{ne se réduit pas} : on laisse l'expression telle quelle (c'est le début d'un polynôme).
\end{propriete}

\begin{attention}[Erreur fatale : additionner les exposants !]
\textbf{Jamais} on n'additionne les exposants lors d'une addition.
\begin{align*}
\textbf{Faux :} \quad & 2x^2 + 3x^2 = 5x^4 \quad (\text{confusion avec le produit}) \\
\textbf{Vrai :} \quad & 2x^2 + 3x^2 = (2+3)x^2 = 5x^2 \\
\textbf{Faux :} \quad & 4ab + 2ab = 6a^2b^2 \\
\textbf{Vrai :} \quad & 4ab + 2ab = 6ab
\end{align*}
\textit{Astuce :} Vérifie que la partie littérale ne change \textbf{pas} après l'addition.
\end{attention}

\begin{exoresolu}[Réduire des sommes de monômes]
Réduire si possible :
\begin{enumerate}
    \item $A = 7x^3y - 2x^3y + 4x^2y$
    \item $B = 5a^2b - 3ab^2 + 2a^2b$
    \item $C = -4uv + 6vu - uv^2$
\end{enumerate}
\tcblower
\textbf{Solution.}
\begin{enumerate}
    \item $7x^3y$ et $-2x^3y$ sont semblables (partie littérale $x^3y$). $4x^2y$ est différent (exposant de $x$).
    \[
    A = (7-2)x^3y + 4x^2y = \boxed{5x^3y + 4x^2y}
    \]
    \item $5a^2b$ et $2a^2b$ sont semblables. $-3ab^2$ est différent (exposants inversés).
    \[
    B = (5+2)a^2b - 3ab^2 = \boxed{7a^2b - 3ab^2}
    \]
    \item $vu = uv$ (commutativité). $-4uv$ et $6uv$ sont semblables. $-uv^2$ est différent.
    \[
    C = (-4+6)uv - uv^2 = \boxed{2uv - uv^2}
    \]
\end{enumerate}
\end{exoresolu}

\courssub{Produit de monômes et puissance d'un monôme}

\begin{propriete}[Produit de monômes]
Le produit de deux monômes est un monôme dont :
\begin{itemize}
    \item le coefficient est le \textbf{produit des coefficients} ;
    \item la partie littérale s'obtient en \textbf{multipliant les puissances de même base} (on \textbf{additionne les exposants}).
\end{itemize}
\[
(ax^\alpha y^\beta) \times (bx^{\alpha'} y^{\beta'}) = (ab) \, x^{\alpha+\alpha'} y^{\beta+\beta'}
\]
\end{propriete}

\begin{propriete}[Puissance d'un monôme]
Pour $n \in \mathbb{N}$ :
\[
(ax^\alpha y^\beta)^n = a^n \, x^{n\alpha} \, y^{n\beta}
\]
On \textbf{multiplie} les exposants par $n$ (et on élève le coefficient à la puissance $n$).
\end{propriete}

\begin{exemplebox}[Produits et puissances]
\begin{align*}
(3x^2y) \times (4xy^3) &= (3\times4) \, x^{2+1} y^{1+3} = 12x^3y^4 \\
(-2a^3b)^2 &= (-2)^2 \, a^{3\times2} b^{1\times2} = 4a^6b^2 \\
(5u^2v) \times (-3uv^2) \times (2u) &= (5\times-3\times2) \, u^{2+1+1} v^{1+2} = -30u^4v^3
\end{align*}
\end{exemplebox}

\courssub{Les polynômes : définition, forme réduite et ordonnée}

\begin{definition}[Polynôme]
Un \textbf{polynôme} est une somme finie de monômes. Chaque monôme est appelé un \textbf{terme} du polynôme.
Un polynôme est \textbf{réduit} si \textbf{aucun de ses termes n'est semblable à un autre} (tous les monômes semblables ont été regroupés).
Un polynôme à une variable est \textbf{ordonné} (généralement par degrés \textbf{décroissants}) si ses termes sont rangés du plus grand degré au plus petit.
\end{definition}

\begin{propriete}[Degré d'un polynôme réduit]
Le \textbf{degré} d'un polynôme réduit non nul est le \textbf{plus grand degré} de ses termes.
\begin{itemize}
    \item Un polynôme de degré $1$ est un \textbf{binôme} (ou trinôme) du premier degré.
    \item Un polynôme de degré $2$ est un \textbf{trinôme} du second degré (s'il a 3 termes).
    \item Le terme de degré $0$ (sans variable) s'appelle le \textbf{terme constant}.
\end{itemize}
\end{propriete}

\begin{popfigure}
\centering
\begin{tikzpicture}
\begin{scope}[every node/.style={font=\small}]
\matrix (M) [matrix of nodes,
    nodes={draw=popline, minimum height=0.9cm, text width=3.2cm, align=center, anchor=center},
    column sep=0.3cm, row sep=0.3cm,
    row 1/.style={nodes={fill=popblue, text=white, font=\small\bfseries}},
    row 2/.style={nodes={fill=popblueL}},
    row 3/.style={nodes={fill=poporangeL}},
    row 4/.style={nodes={fill=popgreenL}},
    row 5/.style={nodes={fill=poppurpleL}},
    row 6/.style={nodes={fill=poppinkL}},
] {
    Expression & Forme réduite & Type & Degré & Terme constant \\
    $4x^3 - 2x + x^2 - 7x^3 + 5x - 1$ & $-3x^3 + x^2 + 3x - 1$ & Trinôme (4 termes) & 3 & $-1$ \\
    $2a^2b + 3ab^2 - 5a^2b + ab^2$ & $-3a^2b + 4ab^2$ & Binôme & 3 & $0$ \\
    $7 + 2x - x^2 + 4x^2 - 3$ & $3x^2 + 2x + 4$ & Trinôme & 2 & $4$ \\
    $5xy - 2xy + 3$ & $3xy + 3$ & Binôme & 2 & $3$ \\
};
\end{scope}
\end{tikzpicture}
\legende{Classification d'expressions : réduction, identification du type (binôme/trinôme), degré et terme constant.}
\end{popfigure}

\begin{exemplebox}[Réduire et ordonner un polynôme]
Soit $P(x) = 4x^3 - 2x + x^2 - 7x^3 + 5x - 1$.
\begin{enumerate}
    \item \textbf{Regrouper les termes semblables} (même degré en $x$) :
    \begin{align*}
    \text{Degré 3 :} \quad & 4x^3 - 7x^3 = -3x^3 \\
    \text{Degré 2 :} \quad & x^2 \quad (\text{seul}) \\
    \text{Degré 1 :} \quad & -2x + 5x = 3x \\
    \text{Degré 0 :} \quad & -1 \quad (\text{terme constant})
    \end{align*}
    \item \textbf{Écrire la forme réduite} (somme des résultats) :
    \[
    P(x) = -3x^3 + x^2 + 3x - 1
    \]
    \item \textbf{Vérifier l'ordre} : degrés $3, 2, 1, 0$ \rightarrow déjà décroissant.
    \item \textbf{Caractéristiques} : degré $3$, terme constant $-1$, $4$ termes (polynôme à 4 termes).
\end{enumerate}
\end{exemplebox}

\courssub{Addition, soustraction et opposé de polynômes}

\begin{methode}[Addition / Soustraction de deux polynômes]
\begin{enumerate}
    \item Écrire les deux polynômes \textbf{sous forme réduite et ordonnée} (aligner les termes de même degré verticalement aide à ne rien oublier).
    \item \textbf{Additionner} (ou soustraire) les coefficients des termes de même degré.
    \item Écrire le résultat sous forme réduite et ordonnée.
\end{enumerate}
\end{methode}

\begin{propriete}[Opposé d'un polynôme]
L'opposé d'un polynôme s'obtient en changeant le signe de \textbf{chaque coefficient} (on multiplie le polynôme par $-1$).
\[
-(a_nx^n + \cdots + a_1x + a_0) = -a_nx^n - \cdots - a_1x - a_0
\]
Cela permet de transformer une soustraction en addition : $P - Q = P + (-Q)$.
\end{propriete}

\begin{exoresolu}[Calculer $P(x) + Q(x)$ et $P(x) - Q(x)$]
Soit $P(x) = 2x^3 - 5x^2 + 4x - 7$ et $Q(x) = -x^3 + 3x^2 - 2x + 5$.
\tcblower
\textbf{Solution.}
On aligne par degrés décroissants :
\[
\begin{array}{r|cccc}
     & x^3 & x^2 & x^1 & x^0 \\ \hline
P(x) & 2 & -5 & 4 & -7 \\
Q(x) & -1 & 3 & -2 & 5 \\ \hline
P+Q & 1 & -2 & 2 & -2 \\
P-Q & 3 & -8 & 6 & -12
\end{array}
\]
\begin{align*}
P(x) + Q(x) &= (2-1)x^3 + (-5+3)x^2 + (4-2)x + (-7+5) \\
            &= \boxed{x^3 - 2x^2 + 2x - 2} \\[1em]
P(x) - Q(x) &= (2-(-1))x^3 + (-5-3)x^2 + (4-(-2))x + (-7-5) \\
            &= \boxed{3x^3 - 8x^2 + 6x - 12}
\end{align*}
\end{exoresolu}

\courssub{Multiplication d'un polynôme par un monôme (premier pas vers le développement)}

C'est l'application directe de la \textbf{distributivité} de la multiplication sur l'addition : $k \times (A + B) = kA + kB$.

\begin{methode}[Multiplier un polynôme par un monôme]
\begin{enumerate}
    \item Multiplier le monôme par \textbf{chaque terme} du polynôme (en utilisant les règles du produit de monômes).
    \item Réduire les termes semblables s'il y en a (généralement il n'y en a pas si le polynôme de départ était réduit).
    \item Ordonner le résultat.
\end{enumerate}
\end{methode}

\begin{exemplebox}[Produit monôme $\times$ polynôme]
Calculer $M \times P$ avec $M = 3x^2y$ et $P = 2x^3 - 4xy + 5y^2 - 1$.
\tcblower
On distribue $3x^2y$ sur chaque terme de $P$ :
\begin{align*}
3x^2y \times (2x^3) &= 6x^{5}y \\
3x^2y \times (-4xy) &= -12x^{3}y^{2} \\
3x^2y \times (5y^2) &= 15x^{2}y^{3} \\
3x^2y \times (-1) &= -3x^{2}y
\end{align*}
Aucun terme n'est semblable. On ordonne par degré total décroissant (degrés $6, 5, 5, 3$) :
\[
\boxed{6x^5y - 12x^3y^2 + 15x^2y^3 - 3x^2y}
\]
\end{exemplebox}

\begin{savaistu}[Le vocabulaire des polynômes au Cameroun et ailleurs]
\begin{itemize}
    \item En classe de 3ème, on travaille surtout les polynômes \textbf{à une variable} (souvent notée $x$).
    \item Un polynôme de degré $1$ s'appelle aussi \textbf{expression affine} (lien avec les fonctions affines de 4ème/3ème).
    \item Un polynôme de degré $2$ est un \textbf{trinôme du second degré} s'il a trois termes (ex. $ax^2+bx+c$). C'est l'objet principal de la résolution d'équations $ax^2+bx+c=0$ en 2nde.
    \item Le \textbf{terme constant} est ce qui reste quand on remplace la variable par $0$ : $P(0) = a_0$. C'est l'ordonnée à l'origine de la courbe représentative.
    \item Dans les concours (BEPC, Probatoire), on demande souvent : « \textit{Réduire et ordonner...} » ou « \textit{Déterminer le degré et le terme constant de...} ». Sois précis dans le vocabulaire !
\end{itemize}
\end{savaistu}

\begin{exercice}[Entraînement : Monômes et polynômes]
\begin{enumerate}
    \item Parmi les expressions suivantes, lesquelles sont des monômes ? Pour chacun, donnez son coefficient, sa partie littérale et son degré.
    \[
    A = -7x^2y \quad;\quad B = 3x + 2 \quad;\quad C = \frac{4}{5}a^3b^2 \quad;\quad D = \frac{2}{x} \quad;\quad E = -\sqrt{2}\,u^4v
    \]
    \item Réduire les sommes suivantes :
    \begin{enumerate}
        \item $5x^2y - 3xy^2 + 2x^2y - 4xy^2$
        \item $-2a^3b^2 + 7a^2b^3 - 5a^3b^2 + a^2b^3$
    \end{enumerate}
    \item Soit $P(x) = 3x^4 - 2x^3 + 5x - 7$ et $Q(x) = -x^4 + 4x^3 - 3x^2 + 2x + 1$.
    Calculer $P(x) + Q(x)$ et $P(x) - Q(x)$. Donner le degré et le terme constant de chaque résultat.
    \item Développer et réduire : $(2x - 3)(x^2 + 4x - 5)$ \textit{(indice : distribue $2x$ puis $-3$)}.
    \item \textbf{Problème concret.} Un cultivateur de maïs près de Bafoussam a un champ rectangulaire de longueur $(3x+2)$ mètres et de largeur $(x+4)$ mètres.
    \begin{enumerate}
        \item Exprimer l'aire $A(x)$ de ce champ sous forme de polynôme réduit et ordonné.
        \item Si $x=5$, quelle est l'aire numérique du champ (en m²) ?
    \end{enumerate}
\end{enumerate}
\end{exercice}

\begin{corrige}[Exercice n°1 -- Monômes et polynômes]
\begin{enumerate}
    \item \textbf{Monômes :} $A$, $C$, $E$. \quad \textbf{Non monômes :} $B$ (somme), $D$ (division par variable).
    \begin{itemize}
        \item $A = -7x^2y$ : coeff $-7$, partie littérale $x^2y$, degré $2+1=3$.
        \item $C = \frac{4}{5}a^3b^2$ : coeff $\frac{4}{5}$, partie littérale $a^3b^2$, degré $3+2=5$.
        \item $E = -\sqrt{2}\,u^4v$ : coeff $-\sqrt{2}$, partie littérale $u^4v$, degré $4+1=5$.
    \end{itemize}
    \item \begin{enumerate}
        \item $5x^2y + 2x^2y = 7x^2y$ ; $-3xy^2 - 4xy^2 = -7xy^2$. Résultat : $\boxed{7x^2y - 7xy^2}$.
        \item $-2a^3b^2 - 5a^3b^2 = -7a^3b^2$ ; $7a^2b^3 + a^2b^3 = 8a^2b^3$. Résultat : $\boxed{-7a^3b^2 + 8a^2b^3}$.
    \end{enumerate}
    \item Alignement par degrés :
    \[
    \begin{array}{c|ccccc}
         & x^4 & x^3 & x^2 & x^1 & x^0 \\ \hline
    P & 3 & -2 & 0 & 5 & -7 \\
    Q & -1 & 4 & -3 & 2 & 1 \\ \hline
    P+Q & 2 & 2 & -3 & 7 & -6 \\
    P-Q & 4 & -6 & 3 & 3 & -8
    \end{array}
    \]
    $P+Q = 2x^4 + 2x^3 - 3x^2 + 7x - 6$ (degré 4, terme constant $-6$).
    $P-Q = 4x^4 - 6x^3 + 3x^2 + 3x - 8$ (degré 4, terme constant $-8$).
    \item $(2x-3)(x^2+4x-5) = 2x(x^2+4x-5) - 3(x^2+4x-5)$
    $= 2x^3+8x^2-10x - 3x^2-12x+15 = \boxed{2x^3 + 5x^2 - 22x + 15}$.
    \item \begin{enumerate}
        \item $A(x) = (3x+2)(x+4) = 3x^2 + 12x + 2x + 8 = \boxed{3x^2 + 14x + 8}$.
        \item $x=5 \Rightarrow A(5) = 3\times25 + 14\times5 + 8 = 75 + 70 + 8 = \boxed{153 \text{ m}^2}$.
    \end{enumerate}
\end{enumerate}
\end{corrige}

\begin{aretenir}[Résumé : Monômes et Polynômes]
\begin{itemize}
    \item \textbf{Monôme} : $k x^\alpha y^\beta\ldots$ (produit nombre $\times$ puissances). Degré = somme des exposants.
    \item \textbf{Monômes semblables} = même partie littérale. \textbf{Somme} = somme des coefficients $\times$ partie littérale commune.
    \item \textbf{Produit monômes} : coeffs $\times$, exposants $+$. \textbf{Puissance} : coeff $^n$, exposants $\times n$.
    \item \textbf{Polynôme} = somme de monômes. \textbf{Réduit} = pas de termes semblables. \textbf{Ordonné} = degrés décroissants.
    \item \textbf{Degré polynôme} = max degré des termes. \textbf{Terme constant} = terme degré 0.
    \item \textbf{Addition/Soustraction} : aligner par degrés, opérer sur les coefficients.
    \item \textbf{Produit monôme $\times$ polynôme} : distributivité (multiplier le monôme par chaque terme).
\end{itemize}
\end{aretenir}

Cette maîtrise des monômes et polynômes est ton \textbf{outil de base} pour la suite : le \textbf{développement} (produit de deux polynômes), la \textbf{factorisation} (inverse du développement) et le travail sur les \textbf{fractions rationnelles}. Entraîne-toi à réduire et ordonner vite et sans faute — c'est la condition pour ne pas perdre de points bêtes au BEPC !

\coursec{Développement et réduction}

Après avoir identifié les monômes et les polynômes, nous allons apprendre à transformer un produit de sommes en une somme de produits. C'est l'opération inverse de la factorisation et elle est au cœur de la simplification des expressions littérales.

\courssub{Rappel : la distributivité simple}

La distributivité de la multiplication sur l'addition est la règle fondamentale qui permet de « distribuer » un facteur sur chaque terme d'une somme.

\begin{propriete}[Distributivité simple]
Pour tous nombres réels $k$, $a$, $b$ :
\[
k(a+b)=ka+kb \qquad\text{et}\qquad (a+b)k=ak+bk.
\]
\end{propriete}

\emph{Justification géométrique~.} Considérons un rectangle de largeur $k$ et de longueur $a+b$. Son aire est $k(a+b)$. Si l'on coupe ce rectangle en deux rectangles de largeurs $k$ et de longueurs $a$ et $b$, les aires sont $ka$ et $kb$. La somme des deux aires donne $ka+kb$, d'où l'égalité.

\courssub{Double distributivité}

Lorsque deux sommes sont multipliées, on applique la distributivité simple deux fois de suite.

\begin{propriete}[Double distributivité]
Pour tous nombres réels $a$, $b$, $c$, $d$ :
\[
(a+b)(c+d)=ac+ad+bc+bd.
\]
\end{propriete}

\begin{experience}[Démonstration guidée par les aires]
\begin{enumerate}
\item Soit un rectangle de dimensions $(a+b)$ et $(c+d)$. Son aire est $(a+b)(c+d)$.
\item On trace une ligne verticale qui sépare la largeur $a+b$ en $a$ et $b$, et une ligne horizontale qui sépare la hauteur $c+d$ en $c$ et $d$.
\item On obtient quatre rectangles dont les aires sont respectivement $ac$, $ad$, $bc$, $bd$.
\item La somme de ces quatre aires est $ac+ad+bc+bd$, ce qui prouve l'égalité.
\end{enumerate}
\end{experience}

\begin{popfigure}
\begin{tikzpicture}[scale=0.9, font=\small]
% rectangle total
\draw[thick, popblue] (0,0) rectangle (3,2.5);
% lignes de séparation
\draw[dashed, poporange] (2,0) -- (2,2.5);
\draw[dashed, poporange] (0,1.5) -- (3,1.5);
% étiquettes des dimensions
\node[left] at (0,1.25) {$c+d$};
\node[left] at (0,0.75) {$c$};
\node[left] at (0,2) {$d$};
\node[below] at (1,0) {$a$};
\node[below] at (2.5,0) {$b$};
\node[below] at (1.5,0) {$a+b$};
% aires des sous-rectangles
\node at (1,0.75) {$ac$};
\node at (2.5,0.75) {$bc$};
\node at (1,2) {$ad$};
\node at (2.5,2) {$bd$};
\end{tikzpicture}
\legende{Rectangle $(a+b)\times(c+d)$ pavé en quatre rectangles : illustration de la double distributivité.}
\end{popfigure}

\courssub{Développement avec des signes négatifs}

La règle des signes s'applique exactement comme pour les nombres :
\begin{itemize}
\item $(+)\times(+) = +$ ; $(+)\times(-) = -$ ; $(-)\times(+) = -$ ; $(-)\times(-) = +$.
\end{itemize}
Ainsi, pour développer $(2x-3)(x+4)$, on traite $-3$ comme un terme négatif et on applique la double distributivité :
\[
(2x-3)(x+4)=2x\cdot x + 2x\cdot4 + (-3)\cdot x + (-3)\cdot4 = 2x^2+8x-3x-12.
\]

\courssub{Réduction et ordonnancement}

Après le développement, on regroupe les termes semblables (même partie littérale) et on ordonne le polynôme par degré décroissant.

\begin{definition}[Réduction d'une expression littérale]
Réduire une expression littérale, c'est effectuer toutes les opérations indiquées (développement, regroupement des termes semblables) et écrire le résultat sous forme canonique : somme de monômes ordonnés par degré décroissant.
\end{definition}

\emph{Exemple~.} $2x^2+8x-3x-12 = 2x^2+5x-12$.

\courssub{Les trois identités remarquables}

Ces identités sont des cas particuliers de la double distributivité qui reviennent très souvent. Les connaître par cœur et savoir les démontrer est indispensable pour le BEPC.

\begin{propriete}[Identités remarquables]
Pour tous réels $a$ et $b$ :
\begin{align*}
(a+b)^2 &= a^2 + 2ab + b^2, \\
(a-b)^2 &= a^2 - 2ab + b^2, \\
(a+b)(a-b) &= a^2 - b^2.
\end{align*}
\end{propriete}

\begin{experience}[Démonstrations par développement direct]
\begin{enumerate}
\item $(a+b)^2 = (a+b)(a+b) = a\cdot a + a\cdot b + b\cdot a + b\cdot b = a^2 + 2ab + b^2$.
\item $(a-b)^2 = (a-b)(a-b) = a\cdot a + a\cdot(-b) + (-b)\cdot a + (-b)\cdot(-b) = a^2 -2ab + b^2$.
\item $(a+b)(a-b) = a\cdot a + a\cdot(-b) + b\cdot a + b\cdot(-b) = a^2 - b^2$ (les termes $ab$ et $-ab$ s'annulent).
\end{enumerate}
\end{experience}

\begin{popfigure}
\begin{tikzpicture}[scale=1.1, font=\small]
% carré de côté a+b = 3.5 (a=2, b=1.5)
\draw[thick, popblue] (0,0) rectangle (3.5,3.5);
% lignes de séparation
\draw[dashed, poporange] (2,0) -- (2,3.5);
\draw[dashed, poporange] (0,2) -- (3.5,2);
% dimensions
\node[left] at (0,1) {$a$};
\node[left] at (0,2.75) {$b$};
\node[below] at (1,0) {$a$};
\node[below] at (2.75,0) {$b$};
\node[below] at (1.75,0) {$a+b$};
% aires
\node at (1,1) {$a^2$};
\node at (2.75,1) {$ab$};
\node at (1,2.75) {$ab$};
\node at (2.75,2.75) {$b^2$};
% flèches pour montrer les deux rectangles ab
\draw[<->, popgreen, thick] (2.2,0.5) -- (2.8,0.5);
\draw[<->, popgreen, thick] (0.5,2.2) -- (0.5,2.8);
\end{tikzpicture}
\legende{Carré de côté $(a+b)$ pavé en $a^2$, $b^2$ et deux rectangles $ab$ : illustration de $(a+b)^2 = a^2+2ab+b^2$.}
\end{popfigure}

\courssub{Utilisation astucieuse des identités remarquables}

Ces identités permettent de calculer mentalement certains produits ou carrés.

\begin{exemplebox}[Calculs astucieux]
\begin{itemize}
\item $101^2 = (100+1)^2 = 100^2 + 2\cdot100\cdot1 + 1^2 = 10000 + 200 + 1 = 10201$.
\item $99 \times 101 = (100-1)(100+1) = 100^2 - 1^2 = 10000 - 1 = 9999$.
\item $48^2 = (50-2)^2 = 50^2 - 2\cdot50\cdot2 + 2^2 = 2500 - 200 + 4 = 2304$.
\end{itemize}
\end{exemplebox}

\courssub{Méthode : développer puis réduire $(a+b)(c+d)$}

\begin{methode}[Développer et réduire un produit de deux sommes]
\begin{enumerate}
\item Écrire le produit en appliquant la double distributivité : multiplier chaque terme de la première somme par chaque terme de la seconde.
\item Respecter scrupuleusement la règle des signes pour chaque produit.
\item Regrouper les termes semblables (même partie littérale et même degré).
\item Ordonner le polynôme obtenu par degré décroissant.
\end{enumerate}
\end{methode}

\courssub{Exemples résolus}

\begin{exoresolu}[Développer et réduire $(2x-3)(x+4)$]
Énoncé~: Développer et réduire l'expression $(2x-3)(x+4)$.

\tcblower
\textbf{Solution.}
\begin{align*}
(2x-3)(x+4) &= 2x\cdot x + 2x\cdot4 + (-3)\cdot x + (-3)\cdot4 \quad\text{(double distributivité)}\\
&= 2x^2 + 8x - 3x - 12 \quad\text{(produits, règle des signes)}\\
&= 2x^2 + 5x - 12 \quad\text{(réduction : $8x-3x=5x$)}.
\end{align*}
L'expression réduite est $2x^2+5x-12$.
\end{exoresolu}

\begin{exoresolu}[Développer et réduire $(3x-5)^2$]
Énoncé~: Développer et réduire $(3x-5)^2$.

\tcblower
\textbf{Solution.}
On utilise l'identité remarquable $(a-b)^2 = a^2 - 2ab + b^2$ avec $a=3x$ et $b=5$ :
\begin{align*}
(3x-5)^2 &= (3x)^2 - 2\cdot(3x)\cdot5 + 5^2 \\
&= 9x^2 - 30x + 25.
\end{align*}
Vérification par double distributivité :
\[
(3x-5)(3x-5)=9x^2-15x-15x+25=9x^2-30x+25.
\]
L'expression réduite est $9x^2-30x+25$.
\end{exoresolu}

\courssub{Piège à éviter}

\begin{attention}[Erreur fréquente : $(a+b)^2 = a^2+b^2$]
Beaucoup d'élèves oublient le terme du milieu $2ab$. \textbf{Jamais} on n'écrit $(a+b)^2 = a^2+b^2$ ; le double produit $2ab$ est indispensable. De même, $(a-b)^2 \neq a^2-b^2$ ; le terme $-2ab$ doit apparaître. Pour ne pas tomber dans ce piège, visualisez toujours le pavage du carré (figure ci‑dessus) ou refaites le développement pas à pas.
\end{attention}

\courssub{Le savais‑tu ?}

\begin{savaistu}[Les identités remarquables chez Euclide]
Les trois identités remarquables étaient déjà démontrées géométriquement par Euclide dans les \emph{Éléments} (Livre II, propositions 4, 5 et 6). Il les présentait sous forme de figures : « Si une ligne est coupée en deux parties, le carré sur la ligne totale est égal à la somme des carrés sur les parties plus deux fois le rectangle contenu par les parties. » C'est exactement $(a+b)^2 = a^2+b^2+2ab$, mais sans notation algébrique ! Cette approche géométrique a traversé les siècles et reste la meilleure façon de les comprendre intuitivement.
\end{savaistu}

\coursec{Factorisation}

Dans la section précédente, nous avons appris à \textbf{développer} : transformer un produit en somme pour simplifier une expression. Mais en mathématiques, comme dans la vie, il est souvent tout aussi utile de savoir faire le chemin inverse. Imagine que tu as construit un mur en briques (le développement) ; factoriser, c'est retrouver les palettes de briques identiques pour les ranger proprement dans le dépôt. C'est l'opération inverse, et elle est \emph{indispensable} pour résoudre des équations, simplifier des fractions ou calculer mentalement plus vite.

\courssub{Qu'est-ce que factoriser ?}

\begin{definition}[Factorisation]
\emph{Factoriser} une expression littérale, c'est la réécrire sous la forme d'un \textbf{produit} de facteurs.
C'est l'opération \textbf{inverse du développement} :
\[
\text{Somme (ou différence)} \quad \xrightarrow{\text{factorisation}} \quad \text{Produit}
\]
\[
\text{Produit} \quad \xrightarrow{\text{développement}} \quad \text{Somme (ou différence)}
\]
\end{definition}

\begin{popfigure}
\begin{tikzpicture}[>=Stealth, font=\small, node distance=1.5cm]
    \node[draw=popblue, thick, rounded corners, fill=popblueL, minimum width=4cm, minimum height=2.5cm, align=center] (dev) {
        \textbf{Développement}\\[4pt]
        $k(a+b) = ka + kb$\\[4pt]
        $(a+b)(a-b) = a^2 - b^2$
    };
    \node[draw=poporange, thick, rounded corners, fill=poporangeL, minimum width=4cm, minimum height=2.5cm, align=center, right=of dev] (fact) {
        \textbf{Factorisation}\\[4pt]
        $ka + kb = k(a+b)$\\[4pt]
        $a^2 - b^2 = (a+b)(a-b)$
    };
    \draw[popblue, very thick, <->] (dev.east) -- node[above, font=\bfseries\large, text=popdark] {Opérations inverses} (fact.west);
    
    \node[draw=popgreen, thick, rounded corners, fill=popgreenL, minimum width=6cm, align=center, below=1cm of dev.south, anchor=north] {
        \textbf{Exemple :} $3x(x+2) = 3x^2 + 6x$ \quad $\Longleftrightarrow$ \quad $3x^2 + 6x = 3x(x+2)$
    };
\end{tikzpicture}
\legende{Le développement et la factorisation sont deux opérations inverses l'une de l'autre.}
\end{popfigure}

\begin{aretenir}[Comment reconnaître une factorisation ?]
Une expression est \textbf{factorisée} si l'opération \textbf{principale} (la dernière qu'on effectue en respectant la priorité des opérations) est une \textbf{multiplication}.
\begin{itemize}
    \item $3(x+5)$ est factorisé (produit de $3$ et $x+5$).
    \item $(x-2)(x+4)$ est factorisé (produit de deux sommes).
    \item $3x + 15$ n'est \textbf{pas} factorisé (somme de deux termes).
    \item $x^2 + 2x - 8$ n'est \textbf{pas} factorisé.
\end{itemize}
\end{aretenir}

\courssub{Recherche d'un facteur commun}

C'est la technique de base, celle qu'on utilise \textbf{en premier} systématiquement. Elle repose sur la distributivité lue de droite à gauche : $ka + kb = k(a+b)$.

\begin{propriete}[Mise en évidence d'un facteur commun]
Si une expression est une somme (ou différence) de termes qui possèdent tous un même facteur $k$, on peut « sortir » ce facteur devant une parenthèse :
\[
k \cdot A + k \cdot B = k(A + B)
\]
Le facteur $k$ peut être :
\begin{itemize}
    \item un \textbf{nombre} : $6x + 9 = 3(2x + 3)$
    \item une \textbf{lettre} (ou une puissance de lettre) : $x^2 + 3x = x(x + 3)$
    \item une \textbf{expression entre parenthèses} : $5(x+2) + 3x(x+2) = (x+2)(5+3x)$
\end{itemize}
\end{propriete}

\begin{methode}[Factoriser par facteur commun]
\begin{enumerate}
    \item \textbf{Identifie} le plus grand facteur commun à \textbf{tous} les termes (coefficient numérique $\times$ partie littérale de plus petit degré).
    \item \textbf{Écris} ce facteur devant une parenthèse.
    \item \textbf{Divise} chaque terme par ce facteur pour remplir la parenthèse.
    \item \textbf{Vérifie} en développant mentalement.
\end{enumerate}
\end{methode}

\begin{exemplebox}[Facteur commun simple et littéral]
Factoriser $E = 12x^3 - 8x^2 + 4x$.
\begin{enumerate}
    \item \textbf{Coefficients :} $12, -8, 4$ $\rightarrow$ PGCD = $4$.
    \item \textbf{Partie littérale :} $x^3, x^2, x$ $\rightarrow$ plus petit exposant = $x$.
    \item \textbf{Facteur commun :} $4x$.
    \item \textbf{Factorisation :} $E = 4x \left( \dfrac{12x^3}{4x} - \dfrac{8x^2}{4x} + \dfrac{4x}{4x} \right) = 4x(3x^2 - 2x + 1)$.
\end{enumerate}
\textbf{Vérification :} $4x \times 3x^2 = 12x^3$, $4x \times (-2x) = -8x^2$, $4x \times 1 = 4x$. \checkmark
\end{exemplebox}

\begin{attention}[Le piège du « facteur caché » : ne pas oublier le $1$ !]
Quand un terme \emph{est exactement} le facteur commun, il reste un \textbf{$1$} dans la parenthèse, pas un $0$ ni « rien ».
\[
\textbf{Exemple :} \quad 5x + 5 = 5(x + \mathbf{1}) \quad \text{et non } 5(x) \text{ ou } 5(x+0).
\]
\[
\textbf{Exemple :} \quad x^2 - x = x(x - \mathbf{1}).
\]
\begin{center}
\textbf{Si tu oublies le $1$, ton développement ne retombera pas sur l'expression de départ.}
\end{center}
\end{attention}

\begin{exemplebox}[Facteur commun « caché » : réécrire un terme]
Parfois, le facteur commun n'est pas visible directement. Il faut réécrire un terme pour le faire apparaître.
Factoriser $F = 2x + 4$.
\begin{itemize}
    \item Je vois $2x$ et $4$. Le PGCD des coefficients est $2$.
    \item J'écris $4 = 2 \times 2$ (je « casse » le $4$ pour faire apparaître le facteur $2$).
    \item $F = 2x + 2 \times 2 = 2(x + 2)$.
\end{itemize}
Factoriser $G = (x+3)^2 + 2(x+3)$.
\begin{itemize}
    \item Le facteur commun est la parenthèse $(x+3)$.
    \item Le premier terme est $(x+3)(x+3)$, le second est $2 \times (x+3)$.
    \item $G = (x+3)\big[(x+3) + 2\big] = (x+3)(x+5)$.
\end{itemize}
\end{exemplebox}

\courssub{Factorisation par les identités remarquables}

Quand il n'y a pas de facteur commun (ou plus), on regarde si l'expression correspond à une identité remarquable lue \textbf{de droite à gauche}. En 3ème, la seule identité de factorisation au programme est la \textbf{différence de deux carrés}.

\begin{propriete}[Différence de deux carrés]
\[
a^2 - b^2 = (a - b)(a + b)
\]
\textbf{Condition :} C'est une \textbf{différence} (soustraction) de \textbf{deux carrés parfaits}.
\end{propriete}

\begin{popfigure}
\begin{tikzpicture}[scale=1.1, font=\small, >=Stealth]
    % --- Figure de gauche : a^2 - b^2 ---
    \begin{scope}[xshift=0cm]
        % Grand carré a x a
        \draw[popdark, very thick] (0,0) rectangle (4,4);
        \node[popdark, font=\bfseries\Large] at (2, 4.3) {$a$};
        \node[popdark, font=\bfseries\Large, rotate=90] at (-0.3, 2) {$a$};
        
        % Petit carré b x b (coin sup droit) - zone enlevée
        \fill[white] (2.5, 2.5) rectangle (4,4); 
        \draw[poporange, thick, dashed] (2.5, 2.5) rectangle (4,4);
        \node[poporange, font=\bfseries] at (3.25, 3.25) {$b^2$};
        \node[poporange, font=\bfseries] at (3.25, 3.6) {$b$};
        \node[poporange, font=\bfseries, rotate=90] at (2.15, 3.25) {$b$};
        
        % Hachures sur le L restant (deux rectangles)
        \fill[pattern=north east lines, pattern color=popblue] (0,0) rectangle (2.5, 4); 
        \fill[pattern=north east lines, pattern color=popblue] (2.5, 0) rectangle (4, 2.5); 
        
        % Cotes restantes
        \node[popdark, font=\bfseries] at (1.25, -0.4) {$a-b$};
        \node[popdark, font=\bfseries, rotate=90] at (-0.4, 2) {$a$};
        \node[popdark, font=\bfseries] at (3.25, -0.4) {$b$};
        \node[popdark, font=\bfseries, rotate=90] at (4.3, 1.25) {$a-b$};
        
        \node[align=center, text=popdark] at (2, -1.2) {Carré $a \times a$ \\ moins carré $b \times b$ \\ $= a^2 - b^2$};
    \end{scope}

    % Flèche
    \draw[popgreen, very thick, ->] (5, 2) -- (6.5, 2) node[midway, above] {Découpage \& Réarrangement};

    % --- Figure de droite : (a-b)(a+b) ---
    \begin{scope}[xshift=8cm]
        % Rectangle (a-b) x (a+b) -> largeur 5.5, hauteur 2.5
        \draw[popgreen, very thick] (0,0) rectangle (5.5, 2.5);
        \node[popgreen, font=\bfseries\Large] at (2.75, 2.8) {$a+b$};
        \node[popgreen, font=\bfseries\Large, rotate=90] at (-0.3, 1.25) {$a-b$};
        
        % Ligne de séparation pour montrer les deux morceaux du L réassemblés
        \draw[poppurple, thick, dashed] (2.5, 0) -- (2.5, 2.5);
        \node[poppurple, font=\small] at (1.25, 1.25) {$(a-b)\times b$};
        \node[poppurple, font=\small] at (4.0, 1.25) {$(a-b)\times a$};
        
        \node[align=center, text=popdark] at (2.75, -1.2) {Rectangle $(a-b) \times (a+b)$ \\ $= a^2 - b^2$};
    \end{scope}
\end{tikzpicture}
\legende{Preuve géométrique de $a^2 - b^2 = (a-b)(a+b)$ : on découpe le L restant pour former un rectangle.}
\end{popfigure}

\begin{exemplebox}[Cas classiques de différence de deux carrés]
\begin{enumerate}
    \item $9x^2 - 25 = (3x)^2 - 5^2 = (3x - 5)(3x + 5)$.
    \item $x^2 - 3 = x^2 - (\sqrt{3})^2 = (x - \sqrt{3})(x + \sqrt{3})$.
    \item $0,01 - 4y^2 = (0,1)^2 - (2y)^2 = (0,1 - 2y)(0,1 + 2y)$.
    \item $\dfrac{1}{4} - z^2 = \left(\dfrac{1}{2}\right)^2 - z^2 = \left(\dfrac{1}{2} - z\right)\left(\dfrac{1}{2} + z\right)$.
\end{enumerate}
\textbf{Astuce :} Pour repérer un carré, regarde si le coefficient est un carré parfait ($1, 4, 9, 16, 25, 36, 49, 64, 81, 100\dots$ ou $0,01, 0,04, 0,09\dots$) et si l'exposant de la variable est \textbf{pair}.
\end{exemplebox}

\begin{attention}[Erreur fatale : La somme de deux carrés ne se factorise pas dans $\mathbb{R}$ !]
\[
\boxed{a^2 + b^2 \neq (a+b)^2 \quad \text{et ne se factorise pas avec des nombres réels.}}
\]
\begin{itemize}
    \item $x^2 + 9$ \textbf{ne se factorise pas} (ce n'est pas $x^2 - 9$).
    \item $4x^2 + 25$ \textbf{ne se factorise pas}.
    \item $(x+3)^2 = x^2 + 6x + 9$ (c'est un développement, pas une factorisation de $x^2+9$).
\end{itemize}
\textbf{Seule la \emph{différence} de deux carrés se factorise : $a^2 - b^2$.}
\end{attention}

\courssub{Factorisations mixtes : la stratégie en deux temps}

C'est le cas le plus fréquent dans les exercices du BEPC : l'expression admet un facteur commun \textbf{ET} une identité remarquable à l'intérieur de la parenthèse. La règle d'or : \textbf{Toujours chercher le facteur commun EN PREMIER.}

\begin{methode}[Factorisation complète (Stratégie officielle)]
Pour factoriser complètement une expression :
\begin{enumerate}
    \item \textbf{Facteur commun :} Sors le plus grand facteur commun à tous les termes. Si l'expression dans la parenthèse est une somme de deux termes, va à l'étape 2.
    \item \textbf{Identité remarquable :} Regarde si l'expression dans la parenthèse est une différence de deux carrés ($A^2 - B^2$). Si oui, factorise-la en $(A-B)(A+B)$.
    \item \textbf{Vérification :} Développe mentalement (ou sur brouillon) pour retrouver l'expression de départ.
\end{enumerate}
\end{methode}

\begin{exoresolu}[Factorisations mixtes]
Factoriser complètement les expressions suivantes :
\begin{enumerate}
    \item $A = 2x^2 - 18$
    \item $B = 5x^3 - 20x$
    \item $C = (2x+1)(x-3) + (2x+1)(x+5)$
\end{enumerate}

\tcblower
\textbf{Solution détaillée :}

\textbf{1. Expression $A = 2x^2 - 18$}
\begin{itemize}
    \item \textbf{Étape 1 (Facteur commun) :} Les coefficients sont $2$ et $-18$. PGCD = $2$. Pas de facteur commun en $x$ (le second terme n'a pas de $x$).
    \item $A = 2(x^2 - 9)$.
    \item \textbf{Étape 2 (Identité) :} Dans la parenthèse : $x^2 - 9 = x^2 - 3^2$. C'est une différence de deux carrés !
    \item $x^2 - 9 = (x - 3)(x + 3)$.
    \item \textbf{Résultat final :} $\boxed{A = 2(x - 3)(x + 3)}$.
    \item \textit{Vérif :} $2 \times (x^2 - 9) = 2x^2 - 18$. \checkmark
\end{itemize}

\textbf{2. Expression $B = 5x^3 - 20x$}
\begin{itemize}
    \item \textbf{Étape 1 :} Coefficients $5$ et $20$ $\rightarrow$ PGCD $5$. Partie littérale $x^3$ et $x$ $\rightarrow$ $x$.
    \item Facteur commun = $5x$.
    \item $B = 5x(x^2 - 4)$.
    \item \textbf{Étape 2 :} $x^2 - 4 = x^2 - 2^2 = (x-2)(x+2)$.
    \item \textbf{Résultat final :} $\boxed{B = 5x(x - 2)(x + 2)}$.
    \item \textit{Vérif :} $5x \times (x^2 - 4) = 5x^3 - 20x$. \checkmark
\end{itemize}

\textbf{3. Expression $C = (2x+1)(x-3) + (2x+1)(x+5)$}
\begin{itemize}
    \item \textbf{Étape 1 :} Le facteur commun est \textbf{la parenthèse $(2x+1)$} elle-même.
    \item $C = (2x+1)\big[(x-3) + (x+5)\big]$.
    \item Simplification dans la parenthèse : $(x-3) + (x+5) = x + x - 3 + 5 = 2x + 2$.
    \item $C = (2x+1)(2x+2)$.
    \item \textbf{Étape 2 :} Regarde $(2x+2)$. Facteur commun $2$ ! $\rightarrow 2(x+1)$.
    \item \textbf{Résultat final :} $\boxed{C = 2(2x+1)(x+1)}$.
    \item \textit{Vérif :} $2(2x+1)(x+1) = (4x+2)(x+1) = 4x^2+4x+2x+2 = 4x^2+6x+2$.
    \item Développement original : $(2x+1)(x-3) + (2x+1)(x+5) = (2x^2-5x-3) + (2x^2+11x+5) = 4x^2+6x+2$. \checkmark
\end{itemize}
\end{exoresolu}

\courssub{Vérification, équations-produits et calcul mental}

\begin{aretenir}[Vérifier une factorisation]
Une factorisation n'est jamais finie tant qu'elle n'est pas \textbf{vérifiée}.
\begin{itemize}
    \item \textbf{Méthode 1 (Rapide) :} Développe mentalement le résultat trouvé. Si tu retrouves l'expression de départ, c'est gagné.
    \item \textbf{Méthode 2 (Test numérique) :} Choisis une valeur simple pour la variable (ex: $x=1$ ou $x=2$), calcule l'expression initiale et ton expression factorisée. Si les deux résultats sont égaux, c'est très bon signe (mais le développement reste la preuve absolue).
\end{itemize}
\end{aretenir}

\begin{propriete}[Résolution d'équation produit : Règle du « Produit Nul »]
Si un produit de facteurs est nul, alors \textbf{au moins un des facteurs est nul}.
\[
A \times B = 0 \quad \Longleftrightarrow \quad A = 0 \quad \text{ou} \quad B = 0
\]
C'est \textbf{LA} raison principale pour laquelle on apprend à factoriser : cela permet de résoudre des équations du second degré (et plus) sans formule magique.
\end{propriete}

\begin{exemplebox}[Résolution d'équation par factorisation]
Résoudre dans $\mathbb{R}$ : $(2x - 5)(x + 3) = 0$.
\begin{align*}
(2x - 5)(x + 3) &= 0 \\
\Leftrightarrow \quad 2x - 5 &= 0 \quad \text{ou} \quad x + 3 = 0 \\
\Leftrightarrow \quad 2x &= 5 \quad \text{ou} \quad x = -3 \\
\Leftrightarrow \quad x &= \frac{5}{2} = 2,5 \quad \text{ou} \quad x = -3
\end{align*}
L'ensemble des solutions est $\boxed{S = \left\{ -3 \,;\, 2,5 \right\}}$.
\end{exemplebox}

\begin{exemplebox}[Calcul mental astucieux grâce à la factorisation]
Calculer $49 \times 51$ sans calculatrice.
\begin{itemize}
    \item On repère $49 = 50 - 1$ et $51 = 50 + 1$.
    \item $49 \times 51 = (50 - 1)(50 + 1) = 50^2 - 1^2 = 2500 - 1 = \mathbf{2499}$.
\end{itemize}
Calculer $103^2 - 97^2$.
\begin{itemize}
    \item Différence de deux carrés : $103^2 - 97^2 = (103 - 97)(103 + 97) = 6 \times 200 = \mathbf{1200}$.
\end{itemize}
\end{exemplebox}

\begin{savaistu}[Le passeport pour la Seconde]
La factorisation n'est pas un exercice de style. En classe de Seconde, tu vas rencontrer systématiquement :
\begin{itemize}
    \item \textbullet\ La résolution d'équations du second degré : $ax^2+bx+c=0$ $\rightarrow$ factorisation en $a(x-x_1)(x-x_2)$.
    \item \textbullet\ L'étude du signe de trinômes (tableaux de signes) : il \textbf{faut} l'expression factorisée.
    \item \textbullet\ La simplification de fractions algébriques : $\dfrac{x^2-4}{x-2} = \dfrac{(x-2)(x+2)}{x-2} = x+2$.
    \item \textbullet\ Les fonctions homographiques, les limites, les dérivées...
\end{itemize}
\textbf{Maîtriser la factorisation en 3ème, c'est s'assurer 50\% des points faciles en Seconde.} Ne la néglige pas !

\begin{center}
\begin{tikzpicture}
\node[draw=popgold, thick, rounded corners, fill=popgoldL, inner sep=10pt, text width=0.85\linewidth, align=left, font=\small] {
    \textbf{\cle{RAPPEL EXAMEN (BEPC)}} \\
    \textbullet\ On te demande souvent : « Factoriser complètement... » $\rightarrow$ N'oublie \textbf{jamais} l'étape 1 (facteur commun). \\
    \textbullet\ « Résoudre l'équation... » $\rightarrow$ Factorise, puis applique la règle du produit nul. \\
    \textbullet\ « Simplifier la fraction... » $\rightarrow$ Factorise numérateur et dénominateur, puis simplifie.
};
\end{tikzpicture}
\end{center}
\end{savaistu}

\begin{exercice}[Entraînement ciblé]
Factoriser complètement les expressions suivantes :
\begin{enumerate}
    \item $E_1 = 14x^2 - 21x$
    \item $E_2 = 25 - 16x^2$
    \item $E_3 = 3x^3 - 12x$
    \item $E_4 = (x-4)^2 - 9$
    \item $E_5 = 2x(x+1) - 3(x+1)$
    \item $E_6 = 50x^2 - 8$
\end{enumerate}
Résoudre dans $\mathbb{R}$ :
\begin{enumerate}
    \item[7.] $(3x+2)(x-4) = 0$
    \item[8.] $x^2 - 16 = 0$
    \item[9.] $2x^2 - 18 = 0$
\end{enumerate}
\end{exercice}

\begin{corrige}[Exercice n°4 - Factorisation]
\begin{enumerate}
    \item $E_1 = 7x(2x - 3)$ \quad (Facteur commun $7x$)
    \item $E_2 = (5 - 4x)(5 + 4x)$ \quad ($5^2 - (4x)^2$, attention à l'ordre dans le premier facteur !)
    \item $E_3 = 3x(x^2 - 4) = 3x(x-2)(x+2)$ \quad (Mixte : $3x$ puis $a^2-b^2$)
    \item $E_4 = [(x-4)-3][(x-4)+3] = (x-7)(x-1)$ \quad (Identité avec $a=x-4, b=3$)
    \item $E_5 = (x+1)(2x-3)$ \quad (Facteur commun $(x+1)$)
    \item $E_6 = 2(25x^2 - 4) = 2((5x)^2 - 2^2) = 2(5x-2)(5x+2)$ \quad (Mixte : $2$ puis $a^2-b^2$)
\end{enumerate}
Résolutions :
7. $3x+2=0 \Rightarrow x=-\frac{2}{3}$ ou $x-4=0 \Rightarrow x=4$. $S=\left\{-\frac{2}{3}; 4\right\}$.
8. $x^2-16=0 \Leftrightarrow (x-4)(x+4)=0 \Rightarrow S=\{-4; 4\}$.
9. $2x^2-18=0 \Leftrightarrow 2(x^2-9)=0 \Leftrightarrow 2(x-3)(x+3)=0 \Rightarrow S=\{-3; 3\}$.
\end{corrige}

\courssub{Fractions rationnelles}

Dans la section précédente, nous avons appris à \emph{factoriser} des expressions littérales : écrire une somme sous forme de produit. Cette compétence est la \textbf{clé} qui va nous permettre d'entrer dans le monde des \emph{fractions rationnelles}. Une fraction rationnelle n'est rien d'autre qu'un quotient de deux polynômes, et pour la simplifier, il faut savoir factoriser le numérateur et le dénominateur. C'est la continuité naturelle de votre travail sur les identités remarquables et la factorisation.

\textbf{Définition et condition d'existence}\par

\begin{definition}[Fraction rationnelle et valeurs interdites]
Une \cle{fraction rationnelle} est une expression de la forme $\dfrac{P(x)}{Q(x)}$ où $P(x)$ et $Q(x)$ sont des polynômes et $Q(x)$ n'est pas le polynôme nul.

La fraction n'a de sens que pour les valeurs de $x$ qui \textbf{ne rendent pas le dénominateur nul}. Ces valeurs sont appelées \cle{valeurs interdites}. Elles sont les solutions de l'équation $Q(x) = 0$.

L'\cle{ensemble de définition} (ou domaine de définition) est l'ensemble des réels \emph{moins} les valeurs interdites : $\mathcal{D} = \mathbb{R} \setminus \{x \in \mathbb{R} \mid Q(x) = 0\}$.
\end{definition}

\begin{exemplebox}[Déterminer les valeurs interdites]
Soit la fraction rationnelle $F(x) = \dfrac{x^2 - 4}{x^2 + x - 6}$.
\begin{enumerate}
\item Le dénominateur est $Q(x) = x^2 + x - 6$.
\item On résout $Q(x) = 0 \iff x^2 + x - 6 = 0$.
\item $\Delta = 1^2 - 4 \times 1 \times (-6) = 25$. Les racines sont $x_1 = \dfrac{-1 - 5}{2} = -3$ et $x_2 = \dfrac{-1 + 5}{2} = 2$.
\item Les valeurs interdites sont \textbf{-3 et 2}. L'ensemble de définition est $\mathcal{D} = \mathbb{R} \setminus \{-3; 2\}$.
\end{enumerate}
\end{exemplebox}

\begin{attention}[Piège : confondre numérateur et dénominateur]
Seul le \textbf{dénominateur} détermine les valeurs interdites. Le numérateur peut s'annuler sans problème (la fraction vaut alors 0). Ne résolvez jamais $P(x)=0$ pour trouver les valeurs interdites !
\end{attention}

\textbf{Simplification d'une fraction rationnelle}\par

Simplifier une fraction rationnelle, c'est la rendre \emph{irréductible} : supprimer les facteurs communs au numérateur et au dénominateur. La règle d'or : \textbf{on ne simplifie que des \emph{facteurs} (éléments multipliés), jamais des \emph{termes} (éléments additionnés)}.

\begin{methode}[Simplifier une fraction rationnelle en 3 étapes]
\begin{enumerate}
\item \textbf{Factoriser} le numérateur $P(x)$ et le dénominateur $Q(x)$ le plus possible (identités remarquables, factorisation par $x$, trinômes, etc.).
\item \textbf{Simplifier} par les facteurs communs \textbf{non nuls} sur l'ensemble de définition.
\item \textbf{Préciser} la condition d'existence (valeurs interdites) et écrire la forme simplifiée en la faisant suivre de la mention « pour $x \neq \dots$ ».
\end{enumerate}
\end{methode}

\begin{popfigure}
\begin{tikzpicture}[
    node distance=1.2cm and 1.5cm,
    box/.style={draw=popblue, thick, rounded corners, fill=popblueL, text width=4.5cm, align=center, minimum height=1.8cm},
    arrow/.style={-Stealth, thick, popdark},
    labelbox/.style={draw=poporange, thick, rounded corners, fill=poporangeL, text width=5cm, align=center, minimum height=1.5cm}
]
\node[box, align=center] (start){Fraction rationnelle $\dfrac{P(x)}{Q(x)}$\\Ex: $\dfrac{x^2-9}{x^2+6x+9}$};
\node[box, right=of start, align=center] (factor){1. Factoriser\\$x^2-9 = (x-3)(x+3)$\\$x^2+6x+9 = (x+3)^2$};
\node[box, right=of factor, align=center] (simplify){2. Simplifier le facteur commun\\$(x+3)$ \textcolor{poporange}{\textbf{si $x+3 \neq 0$}}\\$\dfrac{(x-3)\cancel{(x+3)}}{\cancel{(x+3)}(x+3)} = \dfrac{x-3}{x+3}$};
\node[labelbox, below=of simplify] (cond) {Condition : $x+3 \neq 0 \iff x \neq -3$};
\node[box, right=of simplify, align=center] (end){3. Forme simplifiée\\$\dfrac{x-3}{x+3}$ pour $x \neq -3$};
\draw[arrow] (start) -- (factor) node[midway, above] {\small Factoriser};
\draw[arrow] (factor) -- (simplify) node[midway, above] {\small Simplifier};
\draw[arrow] (simplify) -- (end) node[midway, above] {\small Conclure};
\draw[arrow, poporange] (simplify) -- (cond);
\end{tikzpicture}
\legende{Les trois étapes de la simplification : factoriser, barrer le facteur commun (avec sa condition), conclure.}
\end{popfigure}

\begin{attention}[Erreur fatale : simplifier des termes au lieu de facteurs]
\textbf{Interdiction formelle :} $\dfrac{x+2}{x} \neq 2$ \quad et \quad $\dfrac{x^2+x}{x} \neq x^2+1$ (en barrant les $x$ « en l'air »).
\begin{itemize}
\item Dans $\dfrac{x+2}{x}$, $x$ est un \emph{terme} du numérateur (addition), pas un facteur. On ne peut pas le barrer.
\item Dans $\dfrac{x(x+1)}{x}$, $x$ est un \emph{facteur} (multiplication). On peut simplifier par $x$ (si $x \neq 0$) : $\dfrac{x(x+1)}{x} = x+1$ pour $x \neq 0$.
\end{itemize}
\textbf{Règle :} Mettez toujours entre parenthèses les sommes avant de simplifier. Factorisez \emph{avant} de barrer !
\end{attention}

\begin{exoresolu}[Simplification complète]
\textbf{Énoncé :} Simplifier la fraction rationnelle $G(x) = \dfrac{x^2 - 9}{x^2 + 6x + 9}$ et préciser son ensemble de définition.

\textbf{Solution détaillée :}
\begin{enumerate}
\item \textbf{Valeurs interdites (dénominateur) :} $x^2 + 6x + 9 = 0 \iff (x+3)^2 = 0 \iff x = -3$.
   L'ensemble de définition est $\mathcal{D} = \mathbb{R} \setminus \{-3\}$.
\item \textbf{Factorisation :}
   \begin{align*}
   \text{Numérateur : } & x^2 - 9 = (x-3)(x+3) \quad \text{(différence de deux carrés)} \\
   \text{Dénominateur : } & x^2 + 6x + 9 = (x+3)^2 \quad \text{(identité remarquable $(a+b)^2$)}
   \end{align*}
\item \textbf{Simplification :} Pour $x \neq -3$, le facteur $(x+3)$ est non nul, on peut le simplifier :
   \[
   G(x) = \frac{(x-3)(x+3)}{(x+3)^2} = \frac{(x-3)\cancel{(x+3)}}{\cancel{(x+3)}(x+3)} = \frac{x-3}{x+3}
   \]
\item \textbf{Réponse finale :} $G(x) = \dfrac{x-3}{x+3}$ pour $x \neq -3$.
\end{enumerate}
\end{exoresolu}

\textbf{Opérations sur les fractions rationnelles}\par

Les règles de calcul sont identiques à celles des fractions numériques, à condition de respecter les conditions d'existence de \emph{chaque} fraction intervenant dans le calcul.

\begin{propriete}[Produit et quotient]
Soient $F_1(x) = \dfrac{P_1(x)}{Q_1(x)}$ et $F_2(x) = \dfrac{P_2(x)}{Q_2(x)}$ deux fractions rationnelles.
\begin{itemize}
\item \textbf{Produit :} $F_1(x) \times F_2(x) = \dfrac{P_1(x)P_2(x)}{Q_1(x)Q_2(x)}$.
   \textit{Astuce :} Factorisez \emph{avant} de multiplier, simplifiez les facteurs communs, puis multipliez ce qui reste.
\item \textbf{Quotient :} $F_1(x) \div F_2(x) = F_1(x) \times \dfrac{Q_2(x)}{P_2(x)} = \dfrac{P_1(x)Q_2(x)}{Q_1(x)P_2(x)}$ (pour $P_2(x) \neq 0$).
   On multiplie par l'inverse de la seconde fraction.
\end{itemize}
Les valeurs interdites du résultat sont la réunion des valeurs interdites de $F_1$, $F_2$ (et de l'inverse pour le quotient).
\end{propriete}

\begin{propriete}[Somme et différence : dénominateur commun]
Pour additionner ou soustraire deux fractions rationnelles, on les ramène au \textbf{même dénominateur} (généralement le produit des dénominateurs, ou leur PPCM si on factorise).
\[
\frac{P_1(x)}{Q_1(x)} \pm \frac{P_2(x)}{Q_2(x)} = \frac{P_1(x)Q_2(x) \pm P_2(x)Q_1(x)}{Q_1(x)Q_2(x)}
\]
\textbf{Conseil de pro :} Factorisez les dénominateurs pour trouver le \textbf{PPCM} (Plus Petit Commun Multiple) et éviter des calculs trop lourds.
\end{propriete}

\begin{exemplebox}[Somme avec PPCM]
Calculer $H(x) = \dfrac{2}{x-1} + \dfrac{3}{x+2}$.
\begin{enumerate}
\item Dénominateurs : $x-1$ et $x+2$ (premiers entre eux). PPCM = $(x-1)(x+2)$.
\item Valeurs interdites : $x \neq 1$ et $x \neq -2$.
\item Mise au même dénominateur :
  \[
  H(x) = \frac{2(x+2)}{(x-1)(x+2)} + \frac{3(x-1)}{(x-1)(x+2)} = \frac{2x+4 + 3x-3}{(x-1)(x+2)} = \frac{5x+1}{(x-1)(x+2)}
  \]
\item Résultat : $H(x) = \dfrac{5x+1}{(x-1)(x+2)}$ pour $x \neq 1$ et $x \neq -2$.
\end{enumerate}
\end{exemplebox}

\textbf{Lien avec les équations et le partage proportionnel}\par

Les fractions rationnelles apparaissent naturellement quand on modélise des situations de \emph{vitesse moyenne}, de \emph{débit} ou de \emph{partage proportionnel}. Résoudre une équation contenant des fractions rationnelles revient à lever les dénominateurs (en multipliant par le PPCM) et à vérifier que les solutions trouvées ne sont pas des valeurs interdites.

\begin{exoresolu}[Problème concret : vitesse moyenne]
\textbf{Énoncé :} Un taxi quitte Yaoundé pour Douala (distance 300 km). Il parcourt la première moitié du trajet à 60 km/h et la seconde moitié à 90 km/h. Exprimer sa vitesse moyenne $v_m$ en fonction des vitesses $v_1$ et $v_2$, puis calculer sa valeur.

\textbf{Solution détaillée :}
\begin{enumerate}
\item \textbf{Modélisation :} Soit $d = 300$ km la distance totale.
   Temps 1ère moitié : $t_1 = \frac{d/2}{v_1} = \frac{d}{2v_1}$.
   Temps 2ème moitié : $t_2 = \frac{d/2}{v_2} = \frac{d}{2v_2}$.
   Temps total : $t = t_1 + t_2 = \frac{d}{2}\left(\frac{1}{v_1} + \frac{1}{v_2}\right) = \frac{d}{2} \cdot \frac{v_1+v_2}{v_1v_2}$.
\item \textbf{Vitesse moyenne :} $v_m = \frac{\text{distance totale}}{\text{temps total}} = \frac{d}{t} = \frac{d}{\frac{d}{2} \cdot \frac{v_1+v_2}{v_1v_2}} = \frac{2v_1v_2}{v_1+v_2}$.
   C'est la \cle{moyenne harmonique} des deux vitesses (fraction rationnelle en $v_1, v_2$).
\item \textbf{Calcul numérique :} $v_m = \dfrac{2 \times 60 \times 90}{60 + 90} = \dfrac{10800}{150} = 72$ km/h.
\item \textbf{Interprétation :} La vitesse moyenne (72 km/h) est plus proche de la vitesse la plus faible (60 km/h) car le taxi y a passé plus de temps. Ce n'est \textbf{pas} la moyenne arithmétique $\frac{60+90}{2}=75$ km/h !
\end{enumerate}
\end{exoresolu}

\textbf{Complément : domaine de définition d'une fonction rationnelle (horizon Seconde)}\par

\begin{savaistu}[Fonctions rationnelles et modélisation]
Au-delà du BEPC, en Seconde, on étudie les \textbf{fonctions rationnelles} $f : x \mapsto \frac{P(x)}{Q(x)}$. Leur \cle{domaine de définition} est exactement l'ensemble de définition de la fraction rationnelle : tous les réels \emph{sauf} les zéros du dénominateur.

Ces fonctions modélisent des phénomènes réels variés au Cameroun :
\begin{itemize}
\item \textbf{Télécoms (MTN/Orange) :} Le débit réel d'une connexion selon le nombre d'utilisateurs connectés sur une même antenne-relais suit souvent une loi rationnelle (saturation).
\item \textbf{Électrification rurale :} Le coût moyen par ménage pour raccorder un village au réseau (CAMTEL/ENEO) est une fraction rationnelle : $\frac{\text{coût fixe} + \text{coût variable} \times n}{n}$.
\item \textbf{Agriculture :} Le rendement moyen à l'hectare en fonction de la dose d'engrais suit une courbe rationnelle (loi de Mitscherlich) : augmentations rapides puis plateau.
\end{itemize}
Comprendre les fractions rationnelles en 3ème, c'est poser les bases pour analyser ces situations concrètes plus tard !
\end{savaistu}

\begin{popfigure}
\begin{center}
\begin{tabularx}{\linewidth}{|l|X|X|X|}
\hline
\rowcolor{popblueL}
\textbf{Fraction rationnelle} & \textbf{Valeurs interdites} & \textbf{Forme factorisée} & \textbf{Forme simplifiée} \\
\hline
$\dfrac{x^2-4}{x^2+x-6}$ & $x=-3,\; x=2$ & $\dfrac{(x-2)(x+2)}{(x-2)(x+3)}$ & $\dfrac{x+2}{x+3}$ pour $x \neq 2,\; x \neq -3$ \\
\hline
$\dfrac{x^2+5x+6}{x+2}$ & $x=-2$ & $\dfrac{(x+2)(x+3)}{x+2}$ & $x+3$ pour $x \neq -2$ \\
\hline
$\dfrac{2x}{x^2-9}$ & $x=-3,\; x=3$ & $\dfrac{2x}{(x-3)(x+3)}$ & Déjà irréductible \\
\hline
$\dfrac{x^2-9}{x^2+6x+9}$ & $x=-3$ & $\dfrac{(x-3)(x+3)}{(x+3)^2}$ & $\dfrac{x-3}{x+3}$ pour $x \neq -3$ \\
\hline
\end{tabularx}
\end{center}
\legende{Tableau récapitulatif : de la fraction brute à la forme simplifiée, en passant par la factorisation et l'identification des valeurs interdites.}
\end{popfigure}

\begin{exercice}[Entraînement]
\begin{enumerate}
\item Déterminer l'ensemble de définition et simplifier :
  \begin{itemize}
  \item $A(x) = \dfrac{x^2 - 25}{x^2 - 10x + 25}$
  \item $B(x) = \dfrac{x^2 + 4x + 4}{x^2 - 4}$
  \item $C(x) = \dfrac{3x^2 - 12x}{x^2 - 16}$
  \end{itemize}
\item Effectuer les opérations et simplifier :
  \begin{itemize}
  \item $\dfrac{x+1}{x-2} + \dfrac{2x-3}{x+2}$
  \item $\dfrac{x^2-1}{x+3} \times \dfrac{x+3}{x-1}$
  \item $\dfrac{x^2-4}{x^2+2x} \div \dfrac{x-2}{x}$
  \end{itemize}
\item \textbf{Problème :} Un camion parcourt un trajet aller-retour Yaoundé--Bafoussam. Il va à 50 km/h à l'aller et revient à 70 km/h. Calculer sa vitesse moyenne pour l'aller-retour. (Indication : utiliser la formule de la moyenne harmonique).
\end{enumerate}
\end{exercice}

\begin{corrige}[Exercice n°1 -- Extraits]
\begin{itemize}
\item $A(x) = \frac{(x-5)(x+5)}{(x-5)^2} = \frac{x+5}{x-5}$ pour $x \neq 5$.
\item $B(x) = \frac{(x+2)^2}{(x-2)(x+2)} = \frac{x+2}{x-2}$ pour $x \neq 2,\; x \neq -2$.
\item $C(x) = \frac{3x(x-4)}{(x-4)(x+4)} = \frac{3x}{x+4}$ pour $x \neq 4,\; x \neq -4$.
\item Vitesse moyenne = $\frac{2 \times 50 \times 70}{50+70} = \frac{7000}{120} \approx 58{,}3$ km/h.
\end{itemize}
\end{corrige}

\coursec{Synthèse et applications à la vie courante}

Après avoir maîtrisé les \cle{fractions rationnelles}, nous disposons maintenant d'une boîte à outils complète pour manipuler les \cle{expressions littérales}. Cette section finale a pour but de \textbf{relier entre eux} tous les concepts vus (\cle{expression}, \cle{monôme}, \cle{polynôme}, \cle{développement}, \cle{factorisation}, \cle{fraction}) et de vous montrer comment \textbf{choisir la bonne stratégie} face à un problème concret, comme on vous le demandera au BEPC.

\courssub{Bilan synthétique : la carte mentale du calcul littéral}

Toutes les notions de cette leçon s'articulent autour de l'idée centrale d'\cle{expression littérale}. Pour visualiser les liens, construisons ensemble la carte mentale suivante. Elle vous servira de \og{}tableau de bord\fg{} lors de vos révisions.

\begin{popfigure}
\begin{center}\tcbox[colback=popink,colframe=popink,coltext=white,arc=9pt,boxsep=4pt,left=6pt,right=6pt]{\bfseries\small Calcul Littéral}\end{center}
\vspace{-2pt}
\begin{tcbraster}[raster columns=3,raster equal height=rows,raster column skip=6pt,raster row skip=6pt,enhanced,blankest]
\begin{tcolorbox}[enhanced,colback=white,colframe=poporange,boxrule=0.9pt,arc=3pt,title={Expression littérale},fonttitle=\bfseries\scriptsize,coltitle=white,colbacktitle=poporange,left=4pt,right=4pt,top=2pt,bottom=2pt,boxsep=1pt,toptitle=1.5pt,bottomtitle=1.5pt,halign=left]
\begin{itemize}[nosep,leftmargin=*,itemsep=1pt,font=\scriptsize]
\item Définition : mélange nombres/lettres
\item Valeur numérique par substitution
\end{itemize}
\end{tcolorbox}
\begin{tcolorbox}[enhanced,colback=white,colframe=popgreen,boxrule=0.9pt,arc=3pt,title={Monôme},fonttitle=\bfseries\scriptsize,coltitle=white,colbacktitle=popgreen,left=4pt,right=4pt,top=2pt,bottom=2pt,boxsep=1pt,toptitle=1.5pt,bottomtitle=1.5pt,halign=left]
\begin{itemize}[nosep,leftmargin=*,itemsep=1pt,font=\scriptsize]
\item Forme : $a x^n$
\item Coefficient, degré, partie littérale
\item Opérations : produit, quotient
\end{itemize}
\end{tcolorbox}
\begin{tcolorbox}[enhanced,colback=white,colframe=poppurple,boxrule=0.9pt,arc=3pt,title={Polynôme},fonttitle=\bfseries\scriptsize,coltitle=white,colbacktitle=poppurple,left=4pt,right=4pt,top=2pt,bottom=2pt,boxsep=1pt,toptitle=1.5pt,bottomtitle=1.5pt,halign=left]
\begin{itemize}[nosep,leftmargin=*,itemsep=1pt,font=\scriptsize]
\item Somme de monômes
\item Degré, trier (ordre croissant/décroissant)
\item Réduction : regrouper termes semblables
\end{itemize}
\end{tcolorbox}
\begin{tcolorbox}[enhanced,colback=white,colframe=poppink,boxrule=0.9pt,arc=3pt,title={Développement},fonttitle=\bfseries\scriptsize,coltitle=white,colbacktitle=poppink,left=4pt,right=4pt,top=2pt,bottom=2pt,boxsep=1pt,toptitle=1.5pt,bottomtitle=1.5pt,halign=left]
\begin{itemize}[nosep,leftmargin=*,itemsep=1pt,font=\scriptsize]
\item Distributivité : $a(b+c)=ab+ac$
\item Identités remarquables
\item \textbf{But :} évaluer, simplifier
\end{itemize}
\end{tcolorbox}
\begin{tcolorbox}[enhanced,colback=white,colframe=popgold,boxrule=0.9pt,arc=3pt,title={Factorisation},fonttitle=\bfseries\scriptsize,coltitle=white,colbacktitle=popgold,left=4pt,right=4pt,top=2pt,bottom=2pt,boxsep=1pt,toptitle=1.5pt,bottomtitle=1.5pt,halign=left]
\begin{itemize}[nosep,leftmargin=*,itemsep=1pt,font=\scriptsize]
\item Mise en facteur commun
\item Identités remarquables inverses
\item \textbf{But :} résoudre, trouver zéros, simplifier fractions
\end{itemize}
\end{tcolorbox}
\begin{tcolorbox}[enhanced,colback=white,colframe=popdark,boxrule=0.9pt,arc=3pt,title={Fraction rationnelle},fonttitle=\bfseries\scriptsize,coltitle=white,colbacktitle=popdark,left=4pt,right=4pt,top=2pt,bottom=2pt,boxsep=1pt,toptitle=1.5pt,bottomtitle=1.5pt,halign=left]
\begin{itemize}[nosep,leftmargin=*,itemsep=1pt,font=\scriptsize]
\item Quotient de polynômes
\item Domaine de définition
\item Simplification par factorisation
\end{itemize}
\end{tcolorbox}
\end{tcbraster}

\legende{Carte mentale de synthèse : le calcul littéral en 3ème.}
\end{popfigure}

\courssub{Choisir la bonne forme : développée ou factorisée ?}

C'est LA compétence décisive pour le BEPC. Une même expression peut s'écrire sous forme \textbf{développée} (somme de monômes) ou sous forme \textbf{factorisée} (produit de facteurs). Le choix dépend \textbf{uniquement de la question posée}.

\begin{propriete}[Forme développée vs Forme factorisée]
Soit une expression littérale $E$.
\begin{itemize}
    \item \textbf{Forme développée (somme) :} $E = 3x^2 - 12$.
        \begin{itemize}
            \item \textbf{Utile pour :} \textbf{évaluer} rapidement une valeur numérique (remplacer $x$ par un nombre), \textbf{additionner/soustraire} d'autres polynômes.
        \end{itemize}
    \item \textbf{Forme factorisée (produit) :} $E = 3(x-2)(x+2)$.
        \begin{itemize}
            \item \textbf{Utile pour :} \textbf{résoudre une équation} $E=0$ (produit nul), \textbf{trouver les zéros} (racines), \textbf{étudier le signe} (tableau de signes), \textbf{simplifier une fraction rationnelle}.
        \end{itemize}
\end{itemize}
\end{propriete}

\begin{popfigure}
\begin{center}
\begin{tabularx}{\linewidth}{|l|X|X|}\hline
\rowcolor{popblueL}
\textbf{Expression} & \textbf{Forme développée} & \textbf{Forme factorisée} \\ \hline
$A = (x+3)^2 - 9$ & $x^2 + 6x$ & $x(x+6)$ \\ \hline
\textbf{Avantage principal} & Calculer $A$ pour $x=5$ : $25+30=55$ (rapide). & Résoudre $A=0$ : $x=0$ ou $x=-6$ (immédiat). \\ \hline
$B = \dfrac{x^2-4}{x-2}$ & Inutilisable directement. & $\dfrac{(x-2)(x+2)}{x-2} = x+2$ (pour $x \neq 2$). \\ \hline
\end{tabularx}
\end{center}
\legende{Tableau comparatif : même expression, deux formes, deux usages.}
\end{popfigure}

\begin{methode}[Choisir la forme adaptée à la question]
\begin{enumerate}
    \item \textbf{Lisez la question finale :} Cherche-t-on une \textbf{valeur numérique} ? $\rightarrow$ \textbf{Développez} (et réduisez). Cherche-t-on des \textbf{solutions}, des \textbf{zéros}, ou faut-il \textbf{simplifier une fraction} ? $\rightarrow$ \textbf{Factorisez}.
    \item \textbf{Regardez l'expression de départ :} Est-elle déjà factorisée ? Développée ? Un mélange ?
    \item \textbf{Transformez :} Appliquez la distributivité, les identités remarquables, ou la mise en facteur commun pour obtenir la forme voulue.
    \item \textbf{Vérifiez :} Re-développez votre factorisation pour retomber sur l'original (test du \og{}contrôle inverse\fg{}).
\end{enumerate}
\end{methode}

\courssub{Résolution de problèmes concrets : modéliser, calculer, interpréter}

Au BEPC, les exercices de calcul littéral sont souvent enveloppés dans un \textbf{problème de la vie courante}. La démarche rigoureuse attendue est toujours la même : \textbf{Modéliser $\rightarrow$ Calculer $\rightarrow$ Interpréter}.

\begin{exoresolu}[La vendeuse de pagnes au Marché Central de Yaoundé]
\textbf{Énoncé :} Mme Ngo achète des pagnes à $12\,000$ FCFA l'unité pour les revendre. Elle a des frais fixes de transport et de location de stand s'élevant à $50\,000$ FCFA par mois. Elle vend chaque pagne à $18\,000$ FCFA.
\begin{enumerate}
    \item Exprimer son bénéfice $B$ (en FCFA) en fonction du nombre $x$ de pagnes vendus.
    \item Développer et réduire l'expression de $B$.
    \item Factoriser l'expression obtenue.
    \item Combien de pagnes doit-elle vendre pour ne pas perdre d'argent (seuil de rentabilité) ?
    \item Si elle vend $30$ pagnes, quel est son bénéfice ?
\end{enumerate}
\tcblower
\textbf{Solution détaillée :}

\noindent \textbf{1. Modélisation (Poser les expressions).}
\begin{itemize}
    \item Recette (argent qui entre) : $R(x) = 18\,000 \times x = 18\,000x$.
    \item Coût total (argent qui sort) : $C(x) = \text{achat pagnes} + \text{frais fixes} = 12\,000x + 50\,000$.
    \item Bénéfice : $B(x) = R(x) - C(x) = 18\,000x - (12\,000x + 50\,000)$.
\end{itemize}

\noindent \textbf{2. Développement et réduction (Forme développée pour calculer).}
\begin{align*}
B(x) &= 18\,000x - 12\,000x - 50\,000 \\
     &= \textbf{6\,000x - 50\,000}.
\end{align*}

\noindent \textbf{3. Factorisation (Forme factorisée pour résoudre $B=0$).}
\begin{align*}
B(x) &= 6\,000x - 50\,000 \\
     &= 1\,000(6x - 50) \\
     &= \textbf{2\,000(3x - 25)}.
\end{align*}
\emph{(On a mis en facteur le PGCD des coefficients : $2\,000$).}

\noindent \textbf{4. Seuil de rentabilité (Résolution $B(x) = 0$).}
On utilise la forme factorisée : $2\,000(3x - 25) = 0$.
Puisque $2\,000 \neq 0$, on a $3x - 25 = 0 \iff 3x = 25 \iff x = \frac{25}{3} \approx 8,33$.
\emph{Interprétation :} Elle ne peut pas vendre un tiers de pagne. Elle doit vendre \textbf{au moins 9 pagnes} pour avoir un bénéfice positif (à 9 pagnes : $B(9) = 2\,000(27-25) = 4\,000$ FCFA).

\noindent \textbf{5. Bénéfice pour 30 pagnes (Évaluation).}
On utilise la forme développée : $B(30) = 6\,000 \times 30 - 50\,000 = 180\,000 - 50\,000 = \textbf{130\,000 FCFA}$.
\end{exoresolu}

\begin{attention}[Piège BEPC : Oublier l'interprétation finale]
Dans l'exemple ci-dessus, écrire seulement $x = 25/3$ ou $x \approx 8,33$ \textbf{ne suffit pas} pour avoir tous les points. Le correcteur attend la phrase : \og{}Comme on ne vend que des pagnes entiers, le seuil de rentabilité est atteint à partir de \textbf{9 pagnes}.\fg{}. \textbf{Conclure toujours par une phrase dans le contexte du problème.}
\end{attention}

\courssub{Autres situations types du programme}

Voici deux autres classiques que vous devez savoir traiter instantanément.

\begin{exemplebox}[Aire d'une parcelle agricole (Géométrie + Calcul littéral)]
Un agriculteur possède un champ rectangulaire de longueur $L = (2x + 5)$ mètres et de largeur $\ell = (x - 1)$ mètres. Il agrandit son champ en augmentant la longueur de $3$ m et la largeur de $2$ m.
\begin{enumerate}
    \item Exprimer l'aire $A_1$ du champ initial et l'aire $A_2$ du champ agrandi.
    \item Calculer l'augmentation de surface $\Delta A = A_2 - A_1$ sous forme développée.
    \item Si $x = 10$ m, de combien l'aire a-t-elle augmenté ?
\end{enumerate}
\textbf{Résolution guidée :}
\begin{enumerate}
    \item $A_1 = (2x+5)(x-1) = 2x^2 -2x +5x -5 = 2x^2 + 3x - 5$.
          $A_2 = (2x+5+3)(x-1+2) = (2x+8)(x+1) = 2x^2 + 2x + 8x + 8 = 2x^2 + 10x + 8$.
    \item $\Delta A = A_2 - A_1 = (2x^2 + 10x + 8) - (2x^2 + 3x - 5) = \textbf{7x + 13}$ m$^2$.
          \emph{(Remarquez la simplification des $x^2$ : l'augmentation est linéaire !)}
    \item Pour $x=10$ : $\Delta A = 7\times 10 + 13 = \textbf{83 m$^2$}$.
\end{enumerate}
\end{exemplebox}

\begin{exemplebox}[Bénéfice avec remise commerciale (Pourcentages + Fractions)]
Un commerçant achète un article $P$ FCFA. Il le met en vente avec une majoration de $20\%$. Ne vendant pas, il accorde ensuite une remise de $10\%$ sur le prix de vente.
\begin{enumerate}
    \item Exprimer le prix de vente initial $PV$ et le prix final $PF$ en fonction de $P$.
    \item Montrer que $PF = 1,08 P$. Le commerçant a-t-il fait bénéfice ou perte ?
    \item Exprimer le taux de bénéfice global sous forme de fraction rationnelle simplifiée.
\end{enumerate}
\textbf{Résolution :}
\begin{enumerate}
    \item $PV = P + 0,20P = 1,20P$. \quad $PF = PV - 0,10 PV = 0,90 \times PV = 0,90 \times 1,20 P = \textbf{1,08 P}$.
    \item Comme $1,08 > 1$, $PF > P$ : \textbf{bénéfice}.
    \item Taux = $\frac{PF - P}{P} = \frac{1,08P - P}{P} = \frac{0,08P}{P} = \frac{8}{100} = \mathbf{\frac{2}{25}}$.
\end{enumerate}
\end{exemplebox}

\courssub{Rédaction et justification : la méthode \og{}Annoncer - Calculer - Conclure\fg{}}

Au BEPC, la \textbf{présentation} compte pour près de la moitié des points. Une copie lisible, structurée et justifiée est une copie qui réussit.

\begin{methode}[Rédiger une démarche complète en calcul littéral]
\begin{enumerate}
    \item \textbf{ANNONCER :} Écrivez une phrase d'introduction. \og{}Soit $x$ le nombre de...\fg{}, \og{}On note $A$ l'expression de l'aire...\fg{}, \og{}Le bénéfice $B$ s'exprime par...\fg{}.
    \item \textbf{POSER L'EXPRESSION :} Écrivez la formule brute issue du texte (avec parenthèses si besoin). Ex : $B = 18000x - (12000x + 50000)$.
    \item \textbf{CALCULER (Développer / Factoriser / Simplifier) :} \textbf{Une ligne par étape}. Alignez les signes $=$.
        \begin{itemize}
            \item Justifiez brièvement : \og{}Distributivité\fg{}, \og{}Identité remarquable\fg{}, \og{}Mise en facteur commun par $3x$\fg{}, \og{}Simplification par $x-2$\fg{}.
        \end{itemize}
    \item \textbf{RÉPONDRE À LA QUESTION :} Utilisez la bonne forme (voir méthode précédente).
    \item \textbf{CONCLURE :} Phrase finale répondant exactement à la question, avec unité si besoin. \og{}Le bénéfice est de 130\,000 FCFA.\fg{}, \og{}Il faut vendre au moins 9 pagnes.\fg{}.
    \item \textbf{VÉRIFIER (si temps) :} Testez une valeur numérique dans l'expression de départ et dans votre forme finale.
\end{enumerate}
\end{methode}

\begin{aretenir}[Check-list de présentation pour le BEPC]
\begin{itemize}
    \item \textbf{Un trait de séparation} entre chaque grande étape (Modélisation / Calcul / Conclusion).
    \item \textbf{Les égalités alignées} sur le signe $=$ (utilisez les lignes de votre cahier).
    \item \textbf{Les justifications} entre parenthèses ou au-dessus de la flèche/égal : $\xrightarrow{\text{distributivité}}$, $\xrightarrow{\text{factorisation}}$.
    \item \textbf{Encadrez le résultat final} (cadre ou double trait).
    \item \textbf{Écrivez lisiblement} : distinguez bien le $x$ (lettre) du $\times$ (signe multiplication), soignez vos indices ($x_1$, $x_2$), vos exposants ($x^2$).
\end{itemize}
\end{aretenir}

\courssub{Vérification des résultats : le réflexe de substitution}

Ne rendez jamais une copie sans avoir fait au moins un \textbf{test numérique} (si le temps le permet). C'est la seule façon de détecter une erreur de signe ou de distributivité.

\begin{propriete}[Vérification par substitution]
Si vous avez transformé $E_1$ en $E_2$ (développement, factorisation, simplification), choisissez une valeur \textbf{simple} pour la variable (souvent $x=0$, $x=1$, $x=2$ ou $x=10$), interdite par le domaine de définition si c'est une fraction. Calculez $E_1$ et $E_2$. Si les résultats diffèrent, \textbf{il y a une erreur}.
\end{propriete}

\begin{exemplebox}[Test de vérification sur l'exemple des pagnes]
Expression de départ : $B(x) = 18\,000x - (12\,000x + 50\,000)$.
Forme finale développée : $B(x) = 6\,000x - 50\,000$.
Testons pour $x = 10$ :
\begin{itemize}
    \item Départ : $18\,000\times 10 - (12\,000\times 10 + 50\,000) = 180\,000 - 170\,000 = 10\,000$.
    \item Finale : $6\,000\times 10 - 50\,000 = 60\,000 - 50\,000 = 10\,000$.
\end{itemize}
\textcolor{popgreen}{OK} \textbf{.} Les résultats coïncident.
\end{exemplebox}

\courssub{Préparation au BEPC : gestion du temps et pièges classiques}

\begin{attention}[Les 5 erreurs qui coûtent le plus de points au BEPC]
\begin{enumerate}
    \item \textbf{Oublier les parenthèses lors d'une soustraction :} $A - (B + C) \neq A - B + C$. \textbf{Toujours} écrire $A - B - C$.
    \item \textbf{Confondre développement et factorisation :} Développer quand on demande de résoudre $E=0$ (perte de temps, risque d'erreur) ou factoriser quand on demande de calculer une valeur numérique.
    \item \textbf{Simplifier une fraction sans factoriser :} $\dfrac{x+2}{x} \neq 2$ ! On ne simplifie que des \textbf{facteurs} (produits), jamais des termes (sommes).
    \item \textbf{Oublier le domaine de définition :} Pour une fraction, \textbf{toujours} commencer par \og{}Existe si et seulement si dénominateur $\neq 0$\fg{}.
    \item \textbf{Réponse "nue" sans phrase :} \og{}9\fg{} ne vaut pas grand-chose. \og{}Il faut vendre 9 pagnes\fg{} vaut les points de conclusion.
\end{enumerate}
\end{attention}

\begin{methode}[Gestion du temps sur un exercice de calcul littéral (15-20 min)]
\begin{enumerate}
    \item \textbf{Lecture active (2 min) :} Surlignez les données, la question finale, les mots-clés (\og{}développer\fg{}, \og{}factoriser\fg{}, \og{}résoudre\fg{}, \og{}simplifier\fg{}).
    \item \textbf{Brouillon structuré (5 min) :} Faites les calculs "sales" ici. Testez la substitution.
    \item \textbf{Rédaction propre (8-10 min) :} Recopiez en suivant la méthode \og{}Annoncer-Calculer-Conclure\fg{}. Alignez les $=$. Justifiez.
    \item \textbf{Relecture (2 min) :} Ai-je répondu à la question précise ? L'unité est-elle là ? La conclusion est-elle une phrase ? Le domaine de définition est-il donné (si fraction) ?
\end{enumerate}
\end{methode}

\begin{savaistu}[Le calcul littéral, langage de la modélisation]
Ce que vous apprenez ici (modéliser une situation par une formule, choisir la forme algébrique adaptée à la question, vérifier la cohérence) est \textbf{exactement} ce que font les ingénieurs, économistes, data scientists et architectes au quotidien.
\begin{itemize}
    \item Un ingénieur CAMTEL modélise le débit d'une fibre optique par une fonction rationnelle pour dimensionner les équipements.
    \item Un économiste à la BEAC (Banque des États de l'Afrique Centrale) utilise des polynômes pour prévoir l'inflation.
    \item Un architecte à Douala calcule des aires de surfaces complexes (factorisation/développement) pour optimiser le béton.
\end{itemize}
Le calcul littéral n'est pas un jeu de symboles : c'est le \textbf{langage universel de la résolution de problèmes quantitatifs}.
\end{savaistu}

\begin{exercice}[Entraînement BEPC - Synthèse]
\textbf{Exercice 1.} Soit $E = (2x-3)(x+4) - (x-1)(2x+5)$.
\begin{enumerate}
    \item Développer et réduire $E$.
    \item Factoriser $E$.
    \item Résoudre l'équation $E = 0$.
    \item Calculer $E$ pour $x = -2$. Vérifier avec la forme factorisée.
\end{enumerate}

\textbf{Exercice 2.} Un taxi-brousse fait le trajet Yaoundé-Douala (300 km). Il consomme 8 L/100 km. Le carburant coûte 750 FCFA/L. Le chauffeur a des frais fixes (péage, entretien) de 15\,000 FCFA par trajet. Il transporte 12 passagers.
\begin{enumerate}
    \item Exprimer le coût total $C$ du trajet en fonction du prix du carburant $p$ (en FCFA/L).
    \item Déterminer le prix minimum du billet par passager pour que le chauffeur ne perde pas d'argent (seuil de rentabilité), en fonction de $p$.
    \item Application numérique : $p = 750$ FCFA/L.
\end{enumerate}
\end{exercice}

\begin{corrige}[Exercice 1]
\begin{enumerate}
    \item $E = (2x^2 + 8x - 3x - 12) - (2x^2 + 5x - 2x - 5) = (2x^2 + 5x - 12) - (2x^2 + 3x - 5) = 2x^2 + 5x - 12 - 2x^2 - 3x + 5 = \textbf{2x - 7}$.
    \item $E = 2x - 7$ (déjà factorisé au maximum sur $\mathbb{Z}$, ou $2(x - 3,5)$).
    \item $2x - 7 = 0 \iff 2x = 7 \iff x = \frac{7}{2} = 3,5$.
    \item $x=-2$ : Forme développée $2(-2)-7 = -11$. Forme factorisée (identique) $-11$. \textcolor{popgreen}{OK}.
\end{enumerate}
\end{corrige}

\begin{corrige}[Exercice 2]
\begin{enumerate}
    \item Conso = $8 \times 3 = 24$ L. Coût carburant = $24p$. $C(p) = 24p + 15\,000$.
    \item Recette = $12 \times \text{prix billet} = 12x$. Seuil : $12x = C(p) \iff 12x = 24p + 15\,000 \iff x = \frac{24p + 15\,000}{12} = \textbf{2p + 1\,250}$.
    \item $p=750$ : $x = 2\times 750 + 1\,250 = 1\,500 + 1\,250 = \textbf{2\,750 FCFA}$.
\end{enumerate}
\end{corrige}

\vspace{0.5cm}
\noindent \textbf{Félicitations !} Vous avez maintenant parcouru l'intégralité du programme de Calcul Littéral de 3ème. Vous savez \textbf{écrire}, \textbf{transformer}, \textbf{choisir la forme}, \textbf{résoudre}, \textbf{modéliser} et \textbf{communiquer} vos résultats. Ces compétences sont le socle de toutes les mathématiques du Secondaire (Seconde, Première, Terminale). Travaillez les exercices types, soignez votre rédaction, et le BEPC sera une formalité. \textbf{Bon courage pour la suite !}

\newpage

\coursec{Activité d'intégration}

\coursec{Calcul littéral : Activité d'intégration — Le cas de Mama Ngo Bell}

\courssub{Objectifs de l'activité}
Cette activité d'intégration vise à mobiliser l'ensemble des savoirs et savoir-faire de l'unité « Calcul littéral » dans une situation authentique de la vie économique camerounaise. Elle permet à l'élève de :
\begin{itemize}
    \item \textbf{Modéliser} une situation commerciale concrète par une expression littérale (recette, bénéfice).
    \item \textbf{Manipuler} les expressions littérales : développer, réduire, factoriser, en justifiant chaque étape par les propriétés de la distributivité et les identités remarquables.
    \item \textbf{Résoudre} une équation du premier degré sous forme factorisée (produit de facteurs nul) pour retrouver une donnée initiale (prix d'achat).
    \item \textbf{Interpréter} les résultats numériques obtenus pour prendre une décision argumentée (conseil de gestion).
    \item \textbf{Communiquer} mathématiquement : rédiger une démarche claire, structurée et justifiée, adaptée aux exigences de l'épreuve écrite du BEPC.
\end{itemize}

\begin{exoresolu}[Activité d'intégration : Mama Ngo Bell]
\textbf{Contexte :} Mama Ngo Bell est une commerçante bien connue du quartier Mokolo à Yaoundé. Elle vend des pagnes « wax » de qualité. Pour gérer son stock et fixer ses prix, elle utilise un raisonnement simple mais efficace.

\textbf{Données de la situation :}
\begin{itemize}
    \item Prix d'achat d'un pagne chez le grossiste : \textbf{x} FCFA (inconnu pour l'instant, mais constant).
    \item Prix de vente unitaire pratiqué par Mama Ngo Bell : \textbf{(x + 500)} FCFA. Elle se fixe ainsi une marge brute de 500 FCFA par pagne.
    \item \textbf{Politique commerciale :} Pour toute commande d'au moins 3 pagnes, elle accorde une remise de \textbf{10 \%} sur le montant total de la facture (prix de vente $\times$ quantité).
\end{itemize}

Un jour, un client professionnel (un tailleur du quartier) passe une commande de \textbf{4 pagnes}. Mama Ngo Bell applique donc la remise de 10\%.
On note :
\begin{itemize}
    \item $R(x)$ la recette totale (somme encaissée) pour cette commande de 4 pagnes, exprimée en FCFA.
    \item $B(x)$ le bénéfice net réalisé sur cette commande de 4 pagnes, exprimé en FCFA.
\end{itemize}

\textbf{Consignes}
Travailler par groupes de 3 ou 4 élèves. Chaque groupe rend une feuille de restitution collective, rédigée soigneusement, présentant les réponses aux questions suivantes :

\begin{enumerate}
    \item \textbf{Modélisation de la recette :} Exprimer $R(x)$ en fonction de $x$ en tenant compte de la remise de 10\%. Développer et réduire cette expression pour l'écrire sous la forme canonique $ax + b$ (avec $a$ et $b$ des nombres décimaux).
    \item \textbf{Modélisation du bénéfice :} Rappeler la relation entre Bénéfice, Recette et Coût d'achat. Exprimer $B(x)$ en fonction de $x$. Factoriser complètement l'expression obtenue pour $B(x)$.
    \item \textbf{Application numérique :} Calculer la valeur de la recette $R$ et du bénéfice $B$ dans les deux cas suivants :
    \begin{itemize}
        \item Cas 1 : Le prix d'achat est $x = 2\,500$ FCFA.
        \item Cas 2 : Le prix d'achat est $x = 4\,000$ FCFA.
    \end{itemize}
    \item \textbf{Résolution d'un problème inverse :} On sait qu'en vendant 4 pagnes avec la remise, Mama Ngo Bell a réalisé un bénéfice net de \textbf{3\,600 FCFA}. Retrouver le prix d'achat $x$ d'un pagne en résolvant l'équation $B(x) = 3\,600$. Utiliser la forme factorisée de $B(x)$ pour résoudre l'équation (méthode du produit nul).
    \item \textbf{Analyse et conseil :} En vous appuyant sur les résultats des questions 3 et 4, rédiger un court paragraphe argumenté (5 à 8 lignes) pour conseiller Mama Ngo Bell sur l'impact du prix d'achat sur sa rentabilité, et lui indiquer si sa marge de 500 FCFA par pagne (avant remise) lui semble adaptée pour atteindre un bénéfice cible de 3\,600 FCFA sur 4 pagnes.
\end{enumerate}
\tcblower
\textbf{Préambule :} Nous considérons une commande de $n = 4$ pagnes. Le prix de vente unitaire avant remise est $(x + 500)$ FCFA. La remise de 10\% s'applique sur le total.

\courssub{Question 1 : Expression et développement de la recette $R(x)$}
\textit{Étape 1 : Recette sans remise.}
Pour 4 pagnes vendus au prix unitaire $(x + 500)$ :
\[ \text{Recette sans remise} = 4 \times (x + 500) \]

\textit{Étape 2 : Application de la remise de 10\%.}
Une remise de 10\% signifie que le client paie 90\% du montant. Le multiplicateur est $\num{0,9}$ (ou $\frac{90}{100}$).
\[ R(x) = \num{0,9} \times \bigl[ 4 \times (x + 500) \bigr] \]

\textit{Étape 3 : Développement et réduction (forme canonique $ax+b$).}
On utilise l'associativité et la distributivité :
\begin{align*}
R(x) &= (\num{0,9} \times 4) \times (x + 500) \\
R(x) &= \num{3,6} \times (x + 500) \quad \text{(Calcul du coefficient numérique : $\num{0,9} \times 4 = \num{3,6}$)} \\
R(x) &= \num{3,6} \times x + \num{3,6} \times 500 \quad \text{(Distributivité)} \\
R(x) &= \num{3,6}x + 1\,800 \quad \text{($\num{3,6} \times 500 = 1\,800$)}
\end{align*}
\fbox{$R(x) = \num{3,6}x + 1\,800$} \quad (Expression développée et réduite)

\courssub{Question 2 : Expression et factorisation du bénéfice $B(x)$}
\textit{Étape 1 : Coût d'achat total.}
Mama Ngo Bell a acheté 4 pagnes à $x$ FCFA l'unité.
\[ \text{Coût total} = 4x \]

\textit{Étape 2 : Expression du bénéfice.}
\[ B(x) = \text{Recette} - \text{Coût} = R(x) - 4x \]
En remplaçant $R(x)$ par son expression développée :
\[ B(x) = (\num{3,6}x + 1\,800) - 4x \]

\textit{Étape 3 : Réduction.}
On regroupe les termes en $x$ :
\[ B(x) = (\num{3,6} - 4)x + 1\,800 = -\num{0,4}x + 1\,800 \]

\textit{Étape 4 : Factorisation.}
On cherche un facteur commun aux deux termes $-\num{0,4}x$ et $1\,800$.
$-\num{0,4}x = (-\num{0,4}) \times x$ et $1\,800 = (-\num{0,4}) \times (-4\,500)$ car $-\num{0,4} \times -4\,500 = 1\,800$.
On met $-\num{0,4}$ en facteur commun (on peut aussi mettre $\num{0,4}$ en changeant les signes à l'intérieur) :
\begin{align*}
B(x) &= -\num{0,4}x + 1\,800 \\
B(x) &= -\num{0,4} (x - 4\,500) \quad \text{ou bien} \quad B(x) = \num{0,4} (4\,500 - x)
\end{align*}
La forme $B(x) = \num{0,4} (4\,500 - x)$ est souvent plus lisible car le coefficient devant la parenthèse est positif.
\fbox{$B(x) = \num{0,4} (4\,500 - x)$} \quad (Expression factorisée)

\courssub{Question 3 : Applications numériques}
On utilise la forme développée pour $R(x)$ et la forme factorisée (ou développée) pour $B(x)$.

\begin{center}
\begin{tabularx}{\linewidth}{|c|X|X|}\hline
\textbf{Cas} & \textbf{Recette $R(x) = \num{3,6}x + 1\,800$} & \textbf{Bénéfice $B(x) = -\num{0,4}x + 1\,800$} \\ \hline
$x = 2\,500$ & $R = \num{3,6} \times 2\,500 + 1\,800 = 9\,000 + 1\,800 = \mathbf{10\,800 \text{ FCFA}}$ & $B = -\num{0,4} \times 2\,500 + 1\,800 = -1\,000 + 1\,800 = \mathbf{800 \text{ FCFA}}$ \\ \hline
$x = 4\,000$ & $R = \num{3,6} \times 4\,000 + 1\,800 = 14\,400 + 1\,800 = \mathbf{16\,200 \text{ FCFA}}$ & $B = -\num{0,4} \times 4\,000 + 1\,800 = -1\,600 + 1\,800 = \mathbf{200 \text{ FCFA}}$ \\ \hline
\end{tabularx}
\end{center}

\textit{Analyse rapide :} Lorsque le prix d'achat $x$ augmente, la recette augmente (coefficient $\num{3,6} > 0$) mais le bénéfice diminue (coefficient $-\num{0,4} < 0$). C'est logique : la marge fixe de 500 FCFA est érodée par la remise de 10\% qui porte sur un prix de vente plus élevé.

\courssub{Question 4 : Résolution de l'équation $B(x) = 3\,600$}
On cherche $x$ tel que le bénéfice sur 4 pagnes soit de 3\,600 FCFA.
On utilise la forme factorisée $B(x) = \num{0,4} (4\,500 - x)$.
\begin{align*}
B(x) &= 3\,600 \\
\num{0,4} (4\,500 - x) &= 3\,600
\end{align*}
\textit{Méthode 1 : Division par le coefficient (recommandée au 3ème).}
On divise les deux membres par $\num{0,4}$ (non nul) :
\[ 4\,500 - x = \frac{3\,600}{\num{0,4}} \]
Calcul de la division : $\frac{3\,600}{\num{0,4}} = \frac{36\,000}{4} = 9\,000$.
\[ 4\,500 - x = 9\,000 \]
\[ -x = 9\,000 - 4\,500 \]
\[ -x = 4\,500 \]
\[ x = -4\,500 \]

\textit{Méthode 2 : Développement puis résolution standard.}
\[ \num{0,4} \times 4\,500 - \num{0,4}x = 3\,600 \]
\[ 1\,800 - \num{0,4}x = 3\,600 \]
\[ -\num{0,4}x = 1\,800 \]
\[ x = \frac{1\,800}{-\num{0,4}} = -4\,500 \]

\textbf{Interprétation du résultat :} On trouve $x = -4\,500$ FCFA.
Un prix d'achat \textbf{négatif} est \textbf{impossible} dans la réalité économique.
\fbox{\textbf{Conclusion : Il n'existe aucun prix d'achat positif $x$ permettant d'atteindre un bénéfice de 3\,600 FCFA sur 4 pagnes avec cette politique de prix (marge 500 FCFA + remise 10\%).}}

\courssub{Question 5 : Analyse et conseil argumenté (Exemple de rédaction attendue)}
\textit{« Mama Ngo Bell, l'analyse de ton modèle mathématique montre que ton bénéfice sur 4 pagnes est donné par $B(x) = 1\,800 - \num{0,4}x$. Ce bénéfice diminue quand ton prix d'achat $x$ augmente. Pour $x = 2\,500$ FCFA, tu gagnes 800 FCFA ; pour $x = 4\,000$ FCFA, tu ne gagnes plus que 200 FCFA. Ton objectif de 3\,600 FCFA de bénéfice sur 4 pagnes est \textbf{inatteignable} avec ta marge actuelle de 500 FCFA et la remise de 10\%. En effet, le bénéfice maximal théorique (si tu achetais à 0 FCFA !) serait de 1\,800 FCFA, bien en dessous de 3\,600 FCFA. \textbf{Conseil :} Tu dois soit augmenter ta marge unitaire (vendre plus cher que $x+500$), soit réduire la remise commerciale, soit accepter un bénéfice plus modeste. Ta stratégie actuelle ne permet pas d'atteindre ton objectif financier. »}

\end{exoresolu}

\courssub{Ressources à mobiliser}
Pour mener à bien cette activité, les élèves devront réactiver les connaissances suivantes (cours de la leçon) :

\begin{propriete}[Distributivité simple et double]
Pour tous nombres relatifs $a, b, c, d$ :
\begin{itemize}
    \item $a \times (b + c) = a \times b + a \times c$ \quad (Distributivité simple)
    \item $(a + b) \times (c + d) = ac + ad + bc + bd$ \quad (Distributivité double)
\end{itemize}
C'est la base du \textbf{développement} d'une expression littérale.
\end{propriete}

\begin{propriete}[Réduction d'une expression littérale]
On regroupe les \textbf{termes semblables} (même partie littérale) en additionnant leurs coefficients.
Exemple : $3x + \num{2,5}x - 4 = \num{5,5}x - 4$.
\end{propriete}

\begin{propriete}[Factorisation par mise en facteur commun]
Si une expression est une somme de termes possédant un facteur commun $k$ :
$k \times a + k \times b = k \times (a + b)$.
On « met $k$ devant la parenthèse ». C'est l'opération inverse du développement.
\end{propriete}

\begin{propriete}[Résolution d'une équation produit nul — Niveau 3ème]
Soit $A$ et $B$ deux expressions littérales.
\[ A \times B = 0 \quad \Leftrightarrow \quad A = 0 \quad \text{ou} \quad B = 0 \]
\textbf{Attention :} Cette propriété s'applique uniquement si le produit est égal à \textbf{zéro}. Il faut donc d'abord mettre l'équation sous la forme « Produit = 0 ».
\end{propriete}

\begin{aretenir}[Vocabulaire économique simple]
\begin{itemize}
    \item \textbf{Coût d'achat (ou Prix de revient) :} Somme dépensée pour acheter les marchandises. Ici : $4 \times x$.
    \item \textbf{Recette (ou Chiffre d'affaires) :} Somme encaissée lors de la vente. Ici : $R(x)$.
    \item \textbf{Bénéfice (ou Profit) :} Différence Recette $-$ Coût d'achat. $B = R - \text{Coût}$.
    \item \textbf{Remise de 10\% :} Le client ne paie que 90\% du prix initial. Multiplicateur : $\num{0,9}$ ou $\frac{90}{100}$.
\end{itemize}
\end{aretenir}

\courssub{Bilan de l'activité}
Cette activité a permis de mettre en évidence les points clés suivants, essentiels pour le BEPC :

\begin{enumerate}
    \item \textbf{La modélisation algébrique :} Passer du texte à la formule ($R(x) = \num{0,9} \times 4(x+500)$) demande de bien identifier les opérations (multiplication pour la quantité, multiplication par $\num{0,9}$ pour la remise, soustraction pour le bénéfice).
    \item \textbf{La souplesse des formes d'écriture :} La forme \textbf{développée} ($\num{3,6}x + 1800$) est idéale pour le \textbf{calcul numérique} (remplacement rapide). La forme \textbf{factorisée} ($\num{0,4}(4500-x)$) est idéale pour \textbf{l'analyse du signe}, la \textbf{résolution d'équation} (produit nul) et la \textbf{compréhension de la variation} (ici : le bénéfice baisse quand $x$ monte).
    \item \textbf{La rigueur du langage mathématique :} L'usage du vocabulaire précis (terme, coefficient, factoriser, développer, réduire, terme constant, partie littérale) et la présentation en étapes justifiées (« Par distributivité », « On regroupe les termes en $x$ », « On divise par $\num{0,4}$ ») sont les critères de notation à l'examen.
    \item \textbf{L'interprétation du résultat :} Trouver $x = -4500$ n'est pas une « erreur de calcul », c'est une \textbf{information mathématique forte} qui invalide l'hypothèse économique de départ. Savoir dire « c'est impossible » avec une justification mathématique est une compétence de haut niveau (Savoir-être : esprit critique).
    \item \textbf{Le lien avec le réel :} Les mathématiques ne sont pas abstraites ; elles sont un outil d'aide à la décision pour Mama Ngo Bell comme pour n'importe quel entrepreneur camerounais (commerçant, agriculteur, artisan).
\end{enumerate}

\begin{attention}[Pièges fréquents relevés lors de cette activité]
\begin{itemize}
    \item \textbf{Oublier la remise :} Calculer $R(x) = 4(x+500)$ au lieu de $\num{0,9} \times 4(x+500)$.
    \item \textbf{Erreur sur le multiplicateur :} Croire que « remise de 10\% » signifie multiplier par $\num{0,1}$ (ce qui donne le montant de la remise, pas le prix payé). Le multiplicateur du prix payé est $1 - \num{0,1} = \num{0,9}$.
    \item \textbf{Confondre Recette et Bénéfice :} Oublier de soustraire le coût d'achat ($4x$) pour trouver $B(x)$.
    \item \textbf{Factorisation incorrecte :} Écrire $B(x) = \num{0,4}(4500 - x)$ mais ne pas vérifier en redéveloppant ($\num{0,4} \times 4500 = 1800$, correct). Ou factoriser par $-\num{0,4}$ et se tromper de signe à l'intérieur : $-\num{0,4}(x + 4500)$ au lieu de $-\num{0,4}(x - 4500)$.
    \item \textbf{Résolution d'équation :} Oublier de diviser par $\num{0,4}$ AVANT d'isoler $x$, ou faire $4500 - x = 3600$ (oubli de la division).
    \item \textbf{Non-interprétation :} Donner $x = -4500$ comme réponse finale sans commenter l'absurdité physique (prix négatif).
\end{itemize}
\end{attention}

\begin{savaistu}[Le calcul littéral au service de l'entrepreneuriat au Cameroun]
Au Cameroun, comme partout, la maîtrise du calcul littéral n'est pas réservée aux mathématiciens. Un gérant de \textbf{petite ou moyenne entreprise (PME)} à Douala, un cultivateur de cacao à \textbf{Ebolowa} qui calcule sa marge après transport et taxes, ou un revendeur de crédit \textbf{MTN/Orange} qui gère ses commissions, utilisent implicitement des fonctions affines ($ax+b$) et des raisonnements sur le seuil de rentabilité.
Savoir factoriser $B(x) = \num{0,4}(4500-x)$ permet de voir \textbf{immédiatement} que le bénéfice est nul si $x=4500$ (seuil de rentabilité) et négatif au-delà. C'est la puissance de l'algèbre : \textbf{voir l'invisible} (la structure, le seuil, la tendance) sans refaire 50 calculs numériques. C'est ce qu'on appelle la \textbf{modélisation mathématique}, compétence centrale du programme de 3ème et du BEPC.
\end{savaistu}

\newpage

\coursec{Méthodes et exercices résolus}

\courssub{Méthodes et exercices résolus}

\begin{methode}[Évaluer une expression littérale pour une valeur donnée]
    \textbf{Objectif :} Calculer la valeur numérique d'une expression quand la variable prend une valeur précise.
    \begin{enumerate}
        \item \textbf{Recopier} l'expression en remplaçant \textbf{chaque occurrence} de la variable par la valeur donnée, \textbf{entre parenthèses}.
        \item \textbf{Respecter l'ordre des priorités} : parenthèses $\rightarrow$ puissances $\rightarrow$ multiplications/divisions $\rightarrow$ additions/soustractions.
        \item \textbf{Calculer} étape par étape jusqu'au résultat final.
        \item \textbf{Conclure} par une phrase : « Pour $x = \dots$, l'expression vaut $\dots$ ».
    \end{enumerate}
    \begin{attention}[Piège classique] Ne jamais écrire $3x^2$ avec $x=-2$ comme $3 \times -2^2$ (donne $-12$). Écrire $3 \times (-2)^2 = 3 \times 4 = 12$.
    \end{attention}
\end{methode}

\begin{exoresolu}[Évaluation d'une expression au BEPC]
Soit l'expression $E(x) = 2x^2 - 5x + 3$. Calculer $E(-2)$ et $E(0,5)$.
\tcblower
\textbf{Solution détaillée.}

\underline{Calcul de $E(-2)$ :}
On remplace $x$ par $(-2)$ :
\begin{align*}
E(-2) &= 2 \times (-2)^2 - 5 \times (-2) + 3 \\
      &= 2 \times 4 - 5 \times (-2) + 3 \quad \text{(puissance en priorité)} \\
      &= 8 + 10 + 3 \quad \text{(multiplications)} \\
      &= 21
\end{align*}

\underline{Calcul de $E(0,5)$ :}
On remplace $x$ par $(0,5)$ :
\begin{align*}
E(0,5) &= 2 \times (0,5)^2 - 5 \times (0,5) + 3 \\
       &= 2 \times 0,25 - 2,5 + 3 \quad \text{(puissance et multiplication)} \\
       &= 0,5 - 2,5 + 3 \\
       &= 1
\end{align*}

\textbf{Conclusion :} $E(-2) = 21$ et $E(0,5) = 1$.
\end{exoresolu}

\begin{methode}[Réduire une expression littérale (monômes et polynômes)]
    \textbf{Objectif :} Écrire une somme/de différence de monômes sous la forme la plus simple (un seul polynôme réduit).
    \begin{enumerate}
        \item \textbf{Développer} toutes les parenthèses (distributivité) si nécessaire.
        \item \textbf{Identifier les termes semblables} : même partie littérale (mêmes variables avec mêmes exposants). Ex: $3x^2y$ et $-5x^2y$ sont semblables ; $3x^2y$ et $3xy^2$ ne le sont pas.
        \item \textbf{Regrouper} les coefficients des termes semblables (factoriser la partie littérale).
        \item \textbf{Effectuer les additions/soustractions} des coefficients.
        \item \textbf{Ranger} le polynôme par ordre décroissant des degrés (optionnel mais soigné).
    \end{enumerate}
    \begin{aretenir}[Vocabulaire] Un \cle{monôme} est un produit de nombres et de lettres (ex: $-4x^2y$). Son \cle{coefficient} est le nombre ($-4$). Son \cle{degré} est la somme des exposants ($2+1=3$). Un \cle{polynôme} est une somme de monômes réduits.
    \end{aretenir}
\end{methode}

\begin{exoresolu}[Réduction d'un polynôme]
Réduire l'expression $A = 3x^2 - 2x(4x - 5) + 7x^2 - 10$.
\tcblower
\textbf{Solution détaillée.}

\underline{Étape 1 : Développer le produit $-2x(4x - 5)$.}
\begin{align*}
-2x \times 4x &= -8x^2 \\
-2x \times (-5) &= +10x \\
\text{Donc } -2x(4x - 5) &= -8x^2 + 10x
\end{align*}

\underline{Étape 2 : Réécrire $A$ sans parenthèses.}
\[ A = 3x^2 - 8x^2 + 10x + 7x^2 - 10 \]

\underline{Étape 3 : Regrouper les termes semblables.}
\begin{itemize}
    \item Termes en $x^2$ : $3x^2 - 8x^2 + 7x^2 = (3 - 8 + 7)x^2 = 2x^2$.
    \item Termes en $x$ : $10x$.
    \item Constantes : $-10$.
\end{itemize}

\underline{Étape 4 : Écrire le résultat réduit.}
\[ A = 2x^2 + 10x - 10 \]

\textbf{Conclusion :} La forme réduite de $A$ est $2x^2 + 10x - 10$.
\end{exoresolu}

\begin{methode}[Développer une expression à l'aide des identités remarquables]
    \textbf{Objectif :} Transformer un produit en somme en utilisant les trois identités du programme de 3ème.
    \begin{enumerate}
        \item \textbf{Identifier la structure} : $(a+b)^2$, $(a-b)^2$ ou $(a+b)(a-b)$.
        \item \textbf{Repérer $a$ et $b$} (souvent des monômes, pas seulement des lettres).
        \item \textbf{Appliquer la formule} correspondante :
        \begin{itemize}
            \item $(a+b)^2 = a^2 + 2ab + b^2$
            \item $(a-b)^2 = a^2 - 2ab + b^2$
            \item $(a+b)(a-b) = a^2 - b^2$
        \end{itemize}
        \item \textbf{Simplifier} les monômes obtenus (calculer coefficients et puissances).
        \item \textbf{Réduire} si plusieurs développements sont à additionner.
    \end{enumerate}
    \begin{attention}[Piège du double produit] Dans $(a \pm b)^2$, le terme du milieu est $\pm \mathbf{2}ab$, pas $\pm ab$. Ex: $(3x-2)^2 = 9x^2 - 12x + 4$ (et non $9x^2 - 6x + 4$).
    \end{attention}
\end{methode}

\begin{exoresolu}[Développement et identités remarquables]
Développer et réduire $D = (2x - 3)^2 + (x+5)(x-5)$.
\tcblower
\textbf{Solution détaillée.}

\underline{Premier terme : $(2x - 3)^2$.}
C'est la forme $(a-b)^2$ avec $a = 2x$ et $b = 3$.
\begin{align*}
(2x - 3)^2 &= (2x)^2 - 2 \times (2x) \times 3 + 3^2 \\
           &= 4x^2 - 12x + 9
\end{align*}

\underline{Deuxième terme : $(x+5)(x-5)$.}
C'est la forme $(a+b)(a-b)$ avec $a = x$ et $b = 5$.
\begin{align*}
(x+5)(x-5) &= x^2 - 5^2 = x^2 - 25
\end{align*}

\underline{Somme et réduction :}
\begin{align*}
D &= (4x^2 - 12x + 9) + (x^2 - 25) \\
  &= 4x^2 + x^2 - 12x + 9 - 25 \\
  &= 5x^2 - 12x - 16
\end{align*}

\textbf{Conclusion :} $D = 5x^2 - 12x - 16$.
\end{exoresolu}

\begin{methode}[Factoriser une expression (factorisation simple, identités, trinômes)]
    \textbf{Objectif :} Écrire une somme sous forme de produit.
    \begin{enumerate}
        \item \textbf{Chercher un facteur commun} à tous les termes (PGCD des coefficients + plus petites puissances des variables communes). Le mettre devant une parenthèse.
        \item \textbf{Observer le nombre de termes restants dans la parenthèse} :
            \begin{itemize}
                \item \textbf{2 termes} : vérifier si c'est une différence de carrés $a^2 - b^2 = (a-b)(a+b)$.
                \item \textbf{3 termes} :
                      \begin{itemize}
                          \item Vérifier si c'est une identité remarquable du type $a^2 \pm 2ab + b^2 = (a \pm b)^2$ ($a^2$ et $b^2$ carrés parfaits, terme du milieu $= \pm 2ab$).
                          \item Sinon, pour un trinôme de la forme $x^2 + sx + p$, chercher deux nombres de \textbf{somme $s$} et de \textbf{produit $p$} pour écrire $(x + n_1)(x + n_2)$.
                      \end{itemize}
            \end{itemize}
        \item \textbf{Vérifier} en redéveloppant mentalement.
    \end{enumerate}
    \begin{aretenir}[Ordre de priorité] 1. Facteur commun $\rightarrow$ 2. Identités (3 termes : carré parfait) ou Trinôme $x^2+sx+p$ $\rightarrow$ 3. Différence de carrés (2 termes). On peut combiner : factoriser d'abord le commun, puis factoriser l'intérieur.
    \end{aretenir}
\end{methode}

\begin{exoresolu}[Factorisation complète]
Factoriser les expressions suivantes :
\begin{enumerate}
    \item $E = 12x^3 - 8x^2$
    \item $F = 9x^2 - 12x + 4$
    \item $G = 25x^2 - 49$
    \item $H = x^2 + 5x - 150$
\end{enumerate}
\tcblower
\textbf{Solution détaillée.}

\underline{1. Expression $E = 12x^3 - 8x^2$.}
\begin{itemize}
    \item PGCD des coefficients (12 et 8) : $4$.
    \item Plus petite puissance de $x$ : $x^2$.
    \item Facteur commun : $4x^2$.
\end{itemize}
\[ E = 4x^2 \left( \frac{12x^3}{4x^2} - \frac{8x^2}{4x^2} \right) = 4x^2 (3x - 2) \]
\textit{Vérification :} $4x^2 \times 3x = 12x^3$, $4x^2 \times (-2) = -8x^2$. \textbf{OK.}

\underline{2. Expression $F = 9x^2 - 12x + 4$.}
3 termes. Test identité $(a-b)^2 = a^2 - 2ab + b^2$.
\begin{itemize}
    \item $9x^2 = (3x)^2 \Rightarrow a = 3x$.
    \item $4 = 2^2 \Rightarrow b = 2$.
    \item Terme du milieu : $-2ab = -2 \times 3x \times 2 = -12x$. \textbf{Correspond !}
\end{itemize}
\[ F = (3x - 2)^2 \]

\underline{3. Expression $G = 25x^2 - 49$.}
2 termes, soustraction. Test différence de carrés $a^2 - b^2 = (a-b)(a+b)$.
\begin{itemize}
    \item $25x^2 = (5x)^2 \Rightarrow a = 5x$.
    \item $49 = 7^2 \Rightarrow b = 7$.
\end{itemize}
\[ G = (5x - 7)(5x + 7) \]

\underline{4. Expression $H = x^2 + 5x - 150$.}
3 termes, pas un carré parfait ($150$ n'est pas un carré parfait). Trinôme $x^2 + sx + p$ avec $s=5$, $p=-150$.
Recherche de deux nombres de somme $5$ et produit $-150$ : $+15$ et $-10$.
\[ H = (x + 15)(x - 10) \]

\textbf{Conclusion :} $E = 4x^2(3x-2)$ ; $F = (3x-2)^2$ ; $G = (5x-7)(5x+7)$ ; $H = (x+15)(x-10)$.
\end{exoresolu}

\begin{methode}[Simplifier une fraction rationnelle]
    \textbf{Objectif :} Rendre une fraction irréductible (numérateur et dénominateur premiers entre eux).
    \begin{enumerate}
        \item \textbf{Factoriser} le numérateur et le dénominateur \textbf{le plus complètement possible} (facteur commun + identités + trinômes).
        \item \textbf{Repérer les facteurs communs} (identiques) au numérateur et au dénominateur.
        \item \textbf{Simplifier} ces facteurs communs (les barrer ou les diviser).
        \item \textbf{Écrire le résultat} sous forme de fraction simplifiée (ou expression entière si dénominateur $=1$).
        \item \textbf{Préciser \textbf{systématiquement} les valeurs interdites} (annulent le dénominateur initial).
    \end{enumerate}
    \begin{attention}[Interdiction formelle] On ne simplifie \textbf{JAMAIS} des termes (séparés par $+$ ou $-$). On ne simplifie que des \textbf{facteurs} (multipliés).
    \begin{center}
    $\dfrac{x+3}{x} \neq 3$ \quad $\dfrac{2x}{x} = 2$ (pour $x \neq 0$)
    \end{center}
    \end{attention}
\end{methode}

\begin{exoresolu}[Simplification de fractions rationnelles]
Simplifier les fractions suivantes et donner les valeurs interdites :
\begin{enumerate}
    \item $\dfrac{4x^2 - 16}{2x + 4}$
    \item $\dfrac{x^2 - 5x + 6}{x^2 - 4}$
\end{enumerate}
\tcblower
\textbf{Solution détaillée.}

\underline{1. Fraction $\dfrac{4x^2 - 16}{2x + 4}$.}
\begin{itemize}
    \item \textbf{Numérateur :} $4x^2 - 16 = 4(x^2 - 4) = 4(x-2)(x+2)$ (facteur commun 4 puis différence de carrés).
    \item \textbf{Dénominateur :} $2x + 4 = 2(x+2)$ (facteur commun 2).
\end{itemize}
\begin{align*}
\dfrac{4x^2 - 16}{2x + 4} &= \dfrac{4(x-2)(x+2)}{2(x+2)} \\
&= \dfrac{2 \times 2 \cancel{(x-2)} \cancel{(x+2)}}{2 \cancel{(x+2)}} \quad \text{(simplification par $2$ et $x+2$, si $x \neq -2$)} \\
&= 2(x-2) = 2x - 4
\end{align*}
\textbf{Valeur interdite :} $x = -2$ (annule le dénominateur initial $2x+4$).
\textbf{Résultat :} $2x - 4$ pour $x \neq -2$.

\underline{2. Fraction $\dfrac{x^2 - 5x + 6}{x^2 - 4}$.}
\begin{itemize}
    \item \textbf{Numérateur :} $x^2 - 5x + 6$. Somme $-5$, produit $6 \Rightarrow -2$ et $-3$.
    $x^2 - 5x + 6 = (x-2)(x-3)$.
    \item \textbf{Dénominateur :} $x^2 - 4 = (x-2)(x+2)$ (différence de carrés).
\end{itemize}
\begin{align*}
\dfrac{x^2 - 5x + 6}{x^2 - 4} &= \dfrac{(x-2)(x-3)}{(x-2)(x+2)} \\
&= \dfrac{x-3}{x+2} \quad \text{(simplification par $x-2$, si $x \neq 2$)}
\end{align*}
\textbf{Valeurs interdites :} $x = 2$ et $x = -2$ (annulent le dénominateur initial $x^2-4$).
\textbf{Résultat :} $\dfrac{x-3}{x+2}$ pour $x \neq 2$ et $x \neq -2$.
\end{exoresolu}

\begin{methode}[Résoudre un problème concret par le calcul littéral]
    \textbf{Objectif :} Modéliser une situation de la vie courante (géométrie, commerce, nombres) par une expression, la développer/réduire/factoriser pour répondre à la question.
    \begin{enumerate}
        \item \textbf{Lire} et \textbf{surligner} les grandeurs connues et l'inconnue.
        \item \textbf{Choisir une variable} $x$ pour l'inconnue principale. Exprimer les autres grandeurs en fonction de $x$.
        \item \textbf{Traduire} la relation (égalité, aire, prix, périmètre) en une expression ou équation littérale.
        \item \textbf{Transformer} l'expression (développer, réduire, factoriser) selon ce qui est demandé.
        \item \textbf{Interpréter} le résultat mathématique dans le contexte du problème (unité, sens, positivité).
        \item \textbf{Rédiger la conclusion} en français courant.
    \end{enumerate}
\end{methode}

\begin{exoresolu}[Problème de géométrie et nombres (Type BEPC)]
Un terrain rectangulaire a une longueur supérieure de 5 m à sa largeur.
\begin{enumerate}
    \item Exprimer l'aire $A$ de ce terrain en fonction de la largeur $x$.
    \item Développer et réduire $A$.
    \item Factoriser l'expression obtenue.
    \item Si l'aire vaut $150 \text{ m}^2$, quelle est la largeur du terrain ? (On résoudra l'équation).
\end{enumerate}
\tcblower
\textbf{Solution détaillée.}

\underline{1. Modélisation.}
Largeur $= x$ (en mètres).
Longueur $= x + 5$ (en mètres).
Aire $A = \text{Longueur} \times \text{Largeur}$.
\[ A(x) = x(x + 5) \]

\underline{2. Développement et réduction.}
\[ A(x) = x \times x + x \times 5 = x^2 + 5x \]
L'expression est déjà réduite.

\underline{3. Factorisation.}
Recherche facteur commun : $x$.
\[ A(x) = x(x + 5) \]
(On retrouve la forme factorisée initiale, ce qui est normal).

\underline{4. Résolution pour $A = 150$.}
On résout l'équation : $x^2 + 5x = 150$
\[ x^2 + 5x - 150 = 0 \]
On cherche à factoriser le trinôme : Produit $-150$, Somme $+5$.
Les nombres sont $+15$ et $-10$ ($15 \times -10 = -150$, $15 + (-10) = 5$).
\[ (x + 15)(x - 10) = 0 \]
Un produit est nul si un des facteurs est nul :
\[ x + 15 = 0 \Rightarrow x = -15 \quad \text{(impossible, une largeur $>0$)} \]
\[ x - 10 = 0 \Rightarrow x = 10 \]

\textbf{Conclusion :} La largeur du terrain est de \textbf{10 mètres} (la longueur est de 15 m). Vérification : $10 \times 15 = 150 \text{ m}^2$.
\end{exoresolu}

\begin{methode}[Modéliser une évolution en pourcentage]
    \textbf{Objectif :} Utiliser le calcul littéral pour modéliser une augmentation ou une diminution en pourcentage (lien avec les fonctions affines).
    \begin{enumerate}
        \item \textbf{Identifier} la valeur initiale $V_i$ et le taux $t$ (en \%). $t > 0$ pour augmentation, $t < 0$ pour diminution.
        \item \textbf{Écrire le coefficient multiplicateur} $k = 1 + \frac{t}{100}$.
        \item \textbf{Exprimer la valeur finale} $V_f = k \times V_i$.
        \item \textbf{Si plusieurs variations successives} : multiplier les coefficients ($k_1 \times k_2 \times \dots$).
        \item \textbf{Si on cherche le taux global} : $k_{\text{global}} = k_1 k_2 \dots \Rightarrow t_{\text{global}} = (k_{\text{global}} - 1) \times 100$.
    \end{enumerate}
    \begin{aretenir}[Formule magique] Variation de $t\% \iff$ Multiplier par $\left(1 + \frac{t}{100}\right)$.
    Ex: $+20\% \iff \times 1,2$ ; $-30\% \iff \times 0,7$.
    \end{aretenir}
\end{methode}

\begin{exoresolu}[Évolution du prix du carburant (Contexte Cameroun)]
Le prix du litre de super à la pompe est de 730 FCFA.
\begin{enumerate}
    \item Le prix augmente de 10\%. Écrire littéralement le nouveau prix $P_1$ en fonction de l'ancien $P_0$ et calculer sa valeur.
    \item Quelques mois plus tard, ce nouveau prix diminue de 10\%. Calculer le prix final $P_2$.
    \item Comparer $P_2$ à $P_0$. Le prix est-il revenu à l'initial ? Justifier par le calcul littéral.
\end{enumerate}
\tcblower
\textbf{Solution détaillée.}

\underline{1. Augmentation de 10\%.}
Coefficient multiplicateur $k_1 = 1 + \frac{10}{100} = 1,1$.
Expression littérale : $P_1 = 1,1 \times P_0$.
Avec $P_0 = 730$ : $P_1 = 1,1 \times 730 = 803$ FCFA.

\underline{2. Diminution de 10\% sur $P_1$.}
Coefficient multiplicateur $k_2 = 1 - \frac{10}{100} = 0,9$.
Expression littérale : $P_2 = 0,9 \times P_1$.
Calcul : $P_2 = 0,9 \times 803 = 722,7$ FCFA.

\underline{3. Comparaison et justification littérale.}
\begin{align*}
P_2 &= 0,9 \times P_1 = 0,9 \times (1,1 \times P_0) = (0,9 \times 1,1) \times P_0 \\
    &= 0,99 \times P_0
\end{align*}
Coefficient global $k_g = 0,99$.
Taux global $t_g = (0,99 - 1) \times 100 = -1\%$.

\textbf{Conclusion :} Non, le prix n'est pas revenu à l'initial. Il a \textbf{baissé de 1\%} globalement. $P_2 = 722,7$ FCFA $< 730$ FCFA.
\textit{Note :} Une augmentation de $t\%$ suivie d'une diminution de $t\%$ ramène toujours à une valeur inférieure à la de départ (car $ (1+\frac{t}{100})(1-\frac{t}{100}) = 1 - \frac{t^2}{10000} < 1 $).
\end{exoresolu}

\begin{attention}[Les 5 erreurs qui coûtent des points au BEPC]
    \begin{enumerate}
        \item \textbf{Oublier les parenthèses lors de la substitution.}
        Ex: $x=-2$, calculer $x^2$. Écrire $-2^2 = -4$ (FAUX). Écrire $(-2)^2 = 4$ (VRAI). \textbf{Toujours des parenthèses pour les nombres négatifs.}
        
        \item \textbf{Confondre « terme » et « facteur » lors de la simplification de fractions.}
        $\dfrac{x+3}{3} = x+1$ (FAUX, on ne simplifie pas les termes). $\dfrac{3x}{3} = x$ (VRAI, on simplifie le facteur 3). \textbf{On ne barre que ce qui est multiplié.}
        
        \item \textbf{Oublier le « 2 » dans le double produit des identités.}
        $(3x+2)^2 = 9x^2 + 6x + 4$ (FAUX, il manque le 2). $(3x+2)^2 = 9x^2 + \mathbf{2}\times3x\times2 + 4 = 9x^2 + 12x + 4$ (VRAI).
        
        \item \textbf{Ne pas donner les valeurs interdites pour une fraction rationnelle.}
        Une fraction $\dfrac{N(x)}{D(x)}$ n'a de sens que si $D(x) \neq 0$. La réponse « $x \neq \dots$ » est obligatoire si l'exercice demande de simplifier ou de résoudre une équation rationnelle.
        
        \item \textbf{Mauvaise factorisation du trinôme $x^2 + sx + p$ (signe du produit).}
        Pour $x^2 - 5x - 6$, le produit est $-6$ (négatif) $\Rightarrow$ les deux nombres ont des signes \textbf{contraires}. Ne pas chercher deux nombres de même signe. Ici $+1$ et $-6$ (somme $-5$, produit $-6$). Factorisation : $(x+1)(x-6)$.
    \end{enumerate}
\end{attention}

\newpage

\coursec{Exercices}

\coursec{Calcul littéral}

\courssub{Expressions littérales}

\begin{definition}[Expression littérale]
Une \textbf{expression littérale} est une expression mathématique qui contient des lettres (appelées \textbf{variables} ou \textbf{inconnues}) liées par des opérations ($+$, $-$, $\times$, $\div$, puissances).\\
Exemples : $3x + 5$, $x^2 - 2xy + y^2$, $\dfrac{2x+1}{x-3}$.
\end{definition}

\begin{definition}[Valeur numérique d'une expression]
Donner une \textbf{valeur numérique} à une expression littérale, c'est remplacer la variable par un nombre donné et effectuer les calculs dans l'ordre des priorités (parenthèses, puissances, multiplications/divisions, additions/soustractions).
\end{definition}

\begin{exemplebox}[Évaluation d'une expression]
Soit $A = 2x^2 - 5x + 3$. Calculons la valeur de $A$ pour $x = -2$.
\[
A = 2(-2)^2 - 5(-2) + 3 = 2(4) + 10 + 3 = 8 + 10 + 3 = 21.
\]
\emph{Astuce :} Mettez toujours la valeur entre parenthèses pour éviter les erreurs de signe.
\end{exemplebox}

\begin{attention}[Piège : l'oubli des parenthèses]
Pour $x = -2$, $x^2 = (-2)^2 = 4$ mais $-2^2 = -4$.\\
Écrivez systématiquement : $A = 2(\textcolor{red}{-2})^2 - 5(\textcolor{red}{-2}) + 3$.
\end{attention}

\courssub{Monômes et polynômes}

\begin{definition}[Monôme]
Un \textbf{monôme} est un produit d'un nombre (le \textbf{coefficient}) et de lettres élevées à des puissances entières naturelles (la \textbf{partie littérale}).\\
Le \textbf{degré} d'un monôme est la somme des exposants de ses lettres.\\
Exemples : $5x^2$ (coefficient 5, degré 2), $-3ab^3$ (coefficient -3, degré $1+3=4$), $7$ (monôme constant, degré 0).
\end{definition}

\begin{definition}[Polynôme]
Un \textbf{polynôme} est une somme de monômes (appelés \textbf{termes}).\\
Le \textbf{degré} d'un polynôme est le plus grand degré de ses termes.\\
Vocabulaire : 2 termes $\rightarrow$ \textbf{binôme}, 3 termes $\rightarrow$ \textbf{trinôme}.
\end{definition}

\begin{propriete}[Réduction d'un polynôme]
Réduire un polynôme, c'est regrouper les \textbf{termes semblables} (même partie littérale) en additionnant leurs coefficients.\\
On range généralement le polynôme réduit par \textbf{puissances décroissantes}.
\end{propriete}

\begin{methode}[Réduire un polynôme]
\begin{enumerate}
    \item Repérer les termes semblables (mêmes lettres, mêmes exposants).
    \item Additionner leurs coefficients.
    \item Écrire le résultat en ordonnant les termes par puissances décroissantes.
\end{enumerate}
\end{methode}

\begin{exoresolu}[Réduction]
Réduire $A = 4x^2 - 3x + 7x^2 + 2x - 5$ et $B = 3a^2b - 2ab^2 + 5a^2b + ab^2$.
\tcblower
$A = (4x^2 + 7x^2) + (-3x + 2x) - 5 = 11x^2 - x - 5$.\\
$B = (3a^2b + 5a^2b) + (-2ab^2 + ab^2) = 8a^2b - ab^2$.
\end{exoresolu}

\courssub{Développement et réduction}

\begin{propriete}[Distributivité simple]
Pour tous nombres $k, a, b$ : $k(a + b) = ka + kb$ et $k(a - b) = ka - kb$.
\end{propriete}

\begin{propriete}[Distributivité double]
Pour tous nombres $a, b, c, d$ :
\[
(a + b)(c + d) = ac + ad + bc + bd.
\]
On retient souvent la méthode \textbf{FOIL} : \textbf{F}irst, \textbf{O}uter, \textbf{I}nner, \textbf{L}ast.
\end{propriete}

\begin{aretenir}[Identités remarquables (à connaître par cœur)]
\begin{enumerate}
    \item $(a + b)^2 = a^2 + 2ab + b^2$
    \item $(a - b)^2 = a^2 - 2ab + b^2$
    \item $(a - b)(a + b) = a^2 - b^2$
\end{enumerate}
\emph{Au BEPC, l'utilisation de ces identités doit être explicitée (ex: « d'après l'identité remarquable n°1 »)}.
\end{aretenir}

\begin{methode}[Développer et réduire]
\begin{enumerate}
    \item Appliquer la distributivité (ou une identité remarquable).
    \item Réduire l'expression obtenue en groupant les termes semblables.
    \item Ordonner par puissances décroissantes.
\end{enumerate}
\end{methode}

\begin{attention}[Erreurs fréquentes au développement]
\begin{itemize}
    \item $(a + b)^2 \neq a^2 + b^2$ (oublie le terme du milieu $2ab$).
    \item $(a - b)^2 \neq a^2 - b^2$ (signe du terme du milieu et signe de $b^2$).
    \item Signe moins devant une parenthèse : $-(x - 3) = -x + 3$ (changement de tous les signes).
\end{itemize}
\end{attention}

\begin{exoresolu}[Développement et identités]
Développer et réduire : $C = (2x - 3)(x + 4)$ et $D = (5 - y)^2$.
\tcblower
$C = 2x(x) + 2x(4) - 3(x) - 3(4) = 2x^2 + 8x - 3x - 12 = \textbf{2x}^2 \textbf{+ 5x - 12}$.\\
$D = (5 - y)^2 = 5^2 - 2\times5\times y + y^2 = \textbf{25 - 10y + y}^2$ (ou $y^2 - 10y + 25$).
\end{exoresolu}

\courssub{Factorisation}

\begin{definition}[Factorisation]
Factoriser une expression, c'est l'écrire sous forme d'un \textbf{produit} de facteurs. C'est l'opération inverse du développement.
\end{definition}

\begin{methode}[Méthodes de factorisation (dans l'ordre)]
\begin{enumerate}
    \item \textbf{Facteur commun :} $ab + ac = a(b + c)$. Chercher le plus grand facteur commun aux termes.
    \item \textbf{Identités remarquables :} Reconnaître $a^2+2ab+b^2$, $a^2-2ab+b^2$, $a^2-b^2$.
    \item \textbf{Regroupement / Groupement :} Regrouper des termes pour faire apparaître un facteur commun (souvent une parenthèse).
\end{enumerate}
\end{methode}

\begin{propriete}[Propriété du produit nul]
Soit $A$ et $B$ deux expressions littérales.
\[
A \times B = 0 \quad \Longleftrightarrow \quad A = 0 \quad \text{ou} \quad B = 0.
\]
C'est la méthode principale pour résoudre les équations du type $(ax+b)(cx+d)=0$ ou $x(ax+b)=0$.
\end{propriete}

\begin{attention}[Interdiction de diviser par la variable]
Pour résoudre $x(x-3)=0$, on n'écrit \textbf{JAMAIS} $x-3=0$ (on perd la solution $x=0$).\\
On utilise le produit nul : $x=0$ \textbf{ou} $x-3=0 \Rightarrow x=3$.
\end{attention}

\begin{exoresolu}[Factorisation variée]
Factoriser : $E = 12x^3 - 8x^2$, $F = (2x+1)(x-3) + (2x+1)(4x+5)$, $G = 9x^2 + 12x + 4$.
\tcblower
$E = 4x^2(3x - 2)$ (facteur commun $4x^2$).\\
$F = (2x+1)[(x-3)+(4x+5)] = (2x+1)(5x+2)$ (facteur commun $(2x+1)$).\\
$G = (3x)^2 + 2(3x)(2) + 2^2 = (3x+2)^2$ (identité remarquable n°1).
\end{exoresolu}

\courssub{Fractions rationnelles}

\begin{definition}[Fraction rationnelle]
Une \textbf{fraction rationnelle} est une expression de la forme $\dfrac{N(x)}{D(x)}$ où $N$ et $D$ sont des polynômes et $D \neq 0$.
\end{definition}

\begin{propriete}[Ensemble de définition]
Une fraction rationnelle n'a pas de sens (n'existe pas) pour les valeurs qui annulent le dénominateur.\\
L'\textbf{ensemble de définition} est $\mathbb{R}$ privé de ces \textbf{valeurs interdites}.
\end{propriete}

\begin{propriete}[Simplification]
Si $N(x)$ et $D(x)$ ont un facteur commun non constant, on peut le simplifier :
\[
\dfrac{A(x) \times B(x)}{A(x) \times C(x)} = \dfrac{B(x)}{C(x)} \quad \text{pour } x \text{ tel que } A(x) \neq 0 \text{ et } C(x) \neq 0.
\]
\emph{On factorise systématiquement numérateur et dénominateur avant de simplifier.}
\end{propriete}

\begin{attention}[Règles d'or des fractions]
\begin{itemize}
    \item On simplifie des \textbf{facteurs} (produits), \textbf{jamais des termes} (sommmes).
    \item $\dfrac{x+2}{x+3} \neq \dfrac{2}{3}$.
    \item Les valeurs interdites sont celles du dénominateur \textbf{initial} (on les conserve même après simplification).
\end{itemize}
\end{attention}

\begin{exoresolu}[Simplification de fractions]
Simplifier $F = \dfrac{x^2 - 4x + 4}{x^2 - 4}$ et donner son ensemble de définition.
\tcblower
Numérateur : $x^2 - 4x + 4 = (x-2)^2$.\\
Dénominateur : $x^2 - 4 = (x-2)(x+2)$.\\
Valeurs interdites : $x=2$ et $x=-2$. Ensemble de définition : $\mathbb{R} \setminus \{-2; 2\}$.\\
$F = \dfrac{(x-2)^2}{(x-2)(x+2)} = \dfrac{x-2}{x+2}$ pour $x \neq 2, -2$.
\end{exoresolu}

\courssub{Synthèse et applications à la vie courante}

\courssub{Je vérifie mes connaissances}

\begin{exercice}[QCM : Vocabulaire et notions de base]
Parmi les propositions suivantes, une seule est exacte. Recopie la bonne réponse.
\begin{enumerate}
    \item L'expression $A = 3x^2 - 5x + 7$ est :
    \begin{itemize}
        \item un monôme de degré 2
        \item un binôme de degré 2
        \item \textbf{un trinôme de degré 2}
        \item un polynôme de degré 3
    \end{itemize}
    \item Le coefficient de $x$ dans l'expression développée et réduite de $(2x - 3)(x + 4)$ est :
    \begin{itemize}
        \item \textbf{5}
        \item -5
        \item 11
        \item -11
    \end{itemize}
    \item L'expression $E = \dfrac{x+2}{x^2 - 4}$ n'a pas de sens pour :
    \begin{itemize}
        \item $x = 2$ seulement
        \item $x = -2$ seulement
        \item \textbf{$x = 2$ et $x = -2$}
        \item aucune valeur de $x$
    \end{itemize}
    \item La forme factorisée de $9x^2 - 16$ est :
    \begin{itemize}
        \item $(3x - 4)^2$
        \item $(3x + 4)^2$
        \item \textbf{$(3x - 4)(3x + 4)$}
        \item $(9x - 16)(9x + 16)$
    \end{itemize}
\end{enumerate}
\end{exercice}

\begin{exercice}[Vrai ou Faux ? Justifie ta réponse]
Pour chaque affirmation, réponds \textbf{Vrai} ou \textbf{Faux} et justifie par un calcul ou un contre-exemple.
\begin{enumerate}
    \item $(x + 5)^2 = x^2 + 25$
    \item $x^2 - 7 = (x - 7)(x + 7)$
    \item La fraction $\dfrac{x + 3}{x^2 - 9}$ existe pour tout nombre réel $x$.
    \item $2x^2 + 3x^2 = 5x^4$
    \item Si $A = (x - 1)(x + 1)$, alors $A = 0$ admet deux solutions : $1$ et $-1$.
\end{enumerate}
\end{exercice}

\begin{exercice}[Calculs directs et définitions]
\begin{enumerate}
    \item Donne la définition d'un \textbf{monôme} et d'un \textbf{polynôme}. Donne un exemple de chaque.
    \item Réduis les expressions suivantes :
    \begin{itemize}
        \item $A = 4x^2 - 3x + 7x^2 + 2x - 5$
        \item $B = 3a^2b - 2ab^2 + 5a^2b + ab^2$
    \end{itemize}
    \item Développe (sans réduire) :
    \begin{itemize}
        \item $C = (2x - 3)(x + 4)$
        \item $D = (5 - y)^2$
    \end{itemize}
    \item Factorise par mise en évidence d'un facteur commun :
    \begin{itemize}
        \item $E = 12x^3 - 8x^2$
        \item $F = (2x + 1)(x - 3) + (2x + 1)(4x + 5)$
    \end{itemize}
\end{enumerate}
\end{exercice}

\begin{exercice}[Identités remarquables : reconnaissance et application]
\begin{enumerate}
    \item Complète pour obtenir une identité remarquable :
    \begin{itemize}
        \item $x^2 + 10x + \dots = (x + \dots)^2$
        \item $4y^2 - \dots + 9 = (2y - \dots)^2$
        \item $9x^2 - 25 = (\dots - \dots)(\dots + \dots)$
    \end{itemize}
    \item Calcule astucieusement (sans calculatrice, en montrant l'identité utilisée) :
    \begin{itemize}
        \item $102^2$
        \item $97 \times 103$
        \item $49^2$
    \end{itemize}
\end{enumerate}
\end{exercice}

\courssub{Je m'entraîne}

\begin{exercice}[Développement, réduction et évaluation]
Soit l'expression $A = (3x - 2)(2x + 5) - (x - 4)^2$.
\begin{enumerate}
    \item Développe et réduis $A$. Ordonne le polynôme obtenu suivant les puissances décroissantes de $x$.
    \item Calcule la valeur de $A$ pour $x = -1$.
    \item Factorise l'expression réduite de $A$ si c'est possible.
\end{enumerate}
\end{exercice}

\begin{exercice}[Factorisation : méthodes variées]
Factorise au maximum les expressions suivantes :
\begin{enumerate}
    \item $B = 16x^2 - 25$
    \item $C = (2x - 3)(x + 1) + (2x - 3)(4x - 7)$
    \item $D = 9x^2 + 12x + 4$
    \item $E = x^2 - 6x + 9 - 4y^2$ \emph{(indice : regroupe les trois premiers termes)}
\end{enumerate}
\end{exercice}

\begin{exercice}[Choix de la forme adaptée : résolution d'équation]
On donne $E = (2x + 1)^2 - (2x + 1)(x - 3)$.
\begin{enumerate}
    \item Développe et réduis $E$.
    \item Factorise $E$.
    \item Résous l'équation $E = 0$ en choisissant la forme la plus adaptée (factorisée ou développée). Justifie ton choix.
\end{enumerate}
\end{exercice}

\begin{exercice}[Fractions rationnelles : valeurs interdites et simplification]
Pour chacune des fractions ci-dessous :
\begin{itemize}
    \item Détermine l'ensemble des valeurs interdites (pour lesquelles la fraction n'a pas de sens).
    \item Simplifie la fraction au maximum.
\end{itemize}
\begin{enumerate}
    \item $F = \dfrac{x^2 - 4x + 4}{x^2 - 4}$
    \item $G = \dfrac{3x^2 - 6x}{x^2 - 4}$
    \item $H = \dfrac{x^2 - 9}{x^2 + 6x + 9}$
\end{enumerate}
\end{exercice}

\courssub{Je me prépare au BEPC}

\begin{exercice}[Problème géométrique : Aire d'un terrain]
Un agriculteur de la région de l'Ouest possède un terrain rectangulaire. Sa longueur est $(2x + 5)$ mètres et sa largeur est $(x + 1)$ mètres, avec $x > 0$.
\begin{enumerate}
    \item Exprime l'aire $\mathcal{A}$ de ce terrain en fonction de $x$ sous forme développée et réduite.
    \item Calcule l'aire du terrain si $x = 10$ m.
    \item L'agriculteur souhaite que l'aire de son terrain soit exactement $90 \text{ m}^2$. Détermine la valeur de $x$ qui convient. Vérifie que cette valeur est compatible avec les dimensions du terrain.
    \item Pour cette valeur de $x$, calcule le périmètre du terrain.
\end{enumerate}
\end{exercice}

\begin{exercice}[Modélisation : Tarification d'un transporteur]
Un transporteur de Yaoundé facture ses courses selon le barème suivant :
\begin{itemize}
    \item Forfait de départ : 500 FCFA.
    \item Prix au kilomètre : 150 FCFA/km.
    \item Remise commerciale : 5\% sur le total de la course (forfait + kilométrique) si la distance dépasse 20 km.
\end{itemize}
On note $d$ la distance parcourue en kilomètres ($d > 0$) et $P(d)$ le prix total à payer en FCFA.
\begin{enumerate}
    \item Exprime $P(d)$ en fonction de $d$ pour $d \le 20$. Réduis l'expression.
    \item Exprime $P(d)$ en fonction de $d$ pour $d > 20$. Réduis l'expression. \emph{(Rappel : une remise de 5\% correspond à multiplier par 0,95).}
    \item Un client a payé 4 750 FCFA pour sa course. Détermine la distance $d$ qu'il a parcourue. Précise dans quel cas de figure on se trouve ($d \le 20$ ou $d > 20$).
    \item Le transporteur veut proposer un nouveau forfait unique sans remise : $P_{\text{new}}(d) = 800 + 140d$. À partir de quelle distance (arrondie au km supérieur) le \textbf{nouveau} tarif devient-il plus avantageux que l'ancien (avec remise au-delà de 20 km) ?
\end{enumerate}
\end{exercice}

\begin{exercice}[Synthèse algébrique : Identités et factorisation]
Soient les expressions :
\[
A = (x + 3)^2 - (x - 3)(x + 3) \qquad \text{et} \qquad B = \frac{x^2 - 9}{x^2 + 6x + 9}
\]
\begin{enumerate}
    \item Développe et réduis $A$. Que remarques-tu ? Factorise $A$.
    \item Détermine les valeurs interdites pour $B$. Simplifie $B$ au maximum.
    \item Soit $C = A \times B$. Exprime $C$ sous forme irréductible en fonction de $x$ (en tenant compte des valeurs interdites).
    \item Résous l'équation $C = 2$.
    \item Vérifie si $x = -3$ est une solution de l'équation $C = 2$. Explique pourquoi.
\end{enumerate}
\end{exercice}

\newpage

\coursec{Corrigés des exercices}

\coursec{Corrigés détaillés — Calcul littéral (3ème)}

\begin{corrige}[Exercice 1 — QCM : Vocabulaire et notions de base]
\begin{enumerate}
    \item \textbf{Réponse : un \cle{trinôme} de degré 2.} $A = 3x^2 - 5x + 7$ comporte trois termes (c'est un \cle{trinôme}) et le plus grand exposant de $x$ est 2.
    \item \textbf{Réponse : 5.} $(2x-3)(x+4) = 2x^2 + 8x - 3x - 12 = 2x^2 + 5x - 12$. Le coefficient de $x$ est $5$.
    \item \textbf{Réponse : $x=2$ et $x=-2$.} Le dénominateur $x^2 - 4 = (x-2)(x+2)$ s'annule pour $x=2$ et pour $x=-2$. Ce sont des \cle{valeurs interdites}.
    \item \textbf{Réponse : $(3x-4)(3x+4)$.} $9x^2 - 16 = (3x)^2 - 4^2$, \cle{identité remarquable} $a^2 - b^2 = (a-b)(a+b)$.
\end{enumerate}
\end{corrige}

\begin{corrige}[Exercice 2 — Vrai ou Faux ? Justifie ta réponse]
\begin{enumerate}
    \item \textbf{Faux.} $(x+5)^2 = x^2 + 2\times x\times 5 + 5^2 = x^2 + 10x + 25$ (\cle{identité remarquable n°1} : $(a+b)^2 = a^2 + 2ab + b^2$). Il manque le terme $10x$. Contre-exemple : pour $x=1$, $(1+5)^2 = 36 \neq 1 + 25 = 26$.
    \item \textbf{Faux.} $(x-7)(x+7) = x^2 - 49 \neq x^2 - 7$ (\cle{identité remarquable n°3} : $(a-b)(a+b)=a^2-b^2$). L'expression $x^2 - 7$ ne se factorise pas avec les identités remarquables du programme car $7$ n'est pas un carré parfait dans $\mathbb{Q}$.
    \item \textbf{Faux.} Le dénominateur $x^2 - 9 = (x-3)(x+3)$ s'annule pour $x=3$ et $x=-3$. La \cle{fraction rationnelle} n'existe pas pour ces deux valeurs (valeurs interdites).
    \item \textbf{Faux.} Les termes $2x^2$ et $3x^2$ sont \cle{semblables} (même partie littérale $x^2$) : $2x^2 + 3x^2 = 5x^2$ (on additionne les coefficients, les exposants ne changent pas).
    \item \textbf{Vrai.} $A = 0 \iff (x-1)(x+1) = 0$. Par la \cle{propriété du produit nul} : $x - 1 = 0$ ou $x + 1 = 0$, soit $x = 1$ ou $x = -1$.
\end{enumerate}
\end{corrige}

\begin{corrige}[Exercice 3 — Calculs directs et définitions]
\begin{enumerate}
    \item Un \cle{monôme} est un produit d'un nombre (coefficient) et de lettres élevées à des puissances entières naturelles. Exemple : $-4x^2y$. Un \cle{polynôme} est une somme de monômes. Exemple : $2x^2 - 3x + 1$.
    \item \textbf{Regroupement des termes semblables :}
    \begin{align*}
    A &= (4x^2 + 7x^2) + (-3x + 2x) - 5 = 11x^2 - x - 5,\\
    B &= (3a^2b + 5a^2b) + (-2ab^2 + ab^2) = 8a^2b - ab^2.
    \end{align*}
    \item \textbf{Développement (distributivité) :}
    \begin{align*}
    C &= (2x-3)(x+4) = 2x\times x + 2x\times 4 - 3\times x - 3\times 4 = 2x^2 + 8x - 3x - 12 = 2x^2 + 5x - 12,\\
    D &= (5-y)^2 = 5^2 - 2\times 5\times y + y^2 = 25 - 10y + y^2 \quad \text{(identité n°2 : $(a-b)^2$)}.
    \end{align*}
    \item \textbf{Factorisation (mise en facteur) :}
    \begin{align*}
    E &= 4x^2 \times 3x - 4x^2 \times 2 = 4x^2(3x - 2),\\
    F &= (2x+1)\big[(x-3) + (4x+5)\big] = (2x+1)(5x+2).
    \end{align*}
\end{enumerate}
\end{corrige}

\begin{corrige}[Exercice 4 — Identités remarquables : reconnaissance et application]
\begin{enumerate}
    \item \textbf{Complétion pour obtenir un carré parfait ou une différence de carrés :}
    \begin{itemize}
        \item $x^2 + 10x + \mathbf{25} = (x + \mathbf{5})^2$ car $10x = 2\times x\times 5$ et $5^2 = 25$.
        \item $4y^2 - \mathbf{12y} + 9 = (2y - \mathbf{3})^2$ car $2\times 2y\times 3 = 12y$ et $3^2 = 9$.
        \item $9x^2 - 25 = (\mathbf{3x} - \mathbf{5})(\mathbf{3x} + \mathbf{5})$ (identité $a^2-b^2$).
    \end{itemize}
    \item \textbf{Calcul mental à l'aide des identités :}
    \begin{itemize}
        \item $102^2 = (100 + 2)^2 = 100^2 + 2\times 100\times 2 + 2^2 = 10\,000 + 400 + 4 = \mathbf{10\,404}$.
        \item $97 \times 103 = (100 - 3)(100 + 3) = 100^2 - 3^2 = 10\,000 - 9 = \mathbf{9\,991}$.
        \item $49^2 = (50 - 1)^2 = 50^2 - 2\times 50\times 1 + 1^2 = 2\,500 - 100 + 1 = \mathbf{2\,401}$.
    \end{itemize}
\end{enumerate}
\end{corrige}

\begin{corrige}[Exercice 5 — Développement, réduction et évaluation]
\begin{enumerate}
    \item Développons chaque partie :
    \begin{align*}
    (3x-2)(2x+5) &= 6x^2 + 15x - 4x - 10 = 6x^2 + 11x - 10,\\
    (x-4)^2 &= x^2 - 8x + 16.
    \end{align*}
    Donc $A = 6x^2 + 11x - 10 - (x^2 - 8x + 16) = 6x^2 + 11x - 10 - x^2 + 8x - 16$.\\
    \textbf{Attention au signe moins devant la parenthèse :} tous les signes des termes de la parenthèse changent.
    \[A = 5x^2 + 19x - 26.\]
    \item Pour $x = -1$ : $A = 5(-1)^2 + 19(-1) - 26 = 5 - 19 - 26 = \mathbf{-40}$.
    \item $A = 5x^2 + 19x - 26$ : pas de facteur commun évident, et ce trinôme ne correspond à aucune identité remarquable ($5x^2$ et $-26$ ne sont pas des carrés parfaits compatibles). \textbf{On ne peut pas factoriser $A$ avec les méthodes de 3ème.}
\end{enumerate}
\end{corrige}

\begin{corrige}[Exercice 6 — Factorisation : méthodes variées]
\begin{enumerate}
    \item $B = 16x^2 - 25 = (4x)^2 - 5^2 = \mathbf{(4x-5)(4x+5)}$ (identité $a^2-b^2$).
    \item $C = (2x-3)\big[(x+1) + (4x-7)\big] = (2x-3)(5x - 6)$.\\
    Donc $\mathbf{C = (2x-3)(5x-6)}$ (factorisation par mise en facteur commun du facteur $(2x-3)$).
    \item $D = 9x^2 + 12x + 4 = (3x)^2 + 2\times 3x\times 2 + 2^2 = \mathbf{(3x+2)^2}$ (identité n°1 : $a^2+2ab+b^2$).
    \item On regroupe pour faire apparaître une différence de deux carrés :
    $E = (x^2 - 6x + 9) - 4y^2 = (x-3)^2 - (2y)^2$.\\
    C'est une différence de deux carrés :
    \[E = \big[(x-3) - 2y\big]\big[(x-3) + 2y\big] = \mathbf{(x - 3 - 2y)(x - 3 + 2y)}.\]
\end{enumerate}
\end{corrige}

\begin{corrige}[Exercice 7 — Choix de la forme adaptée : résolution d'équation]
\begin{enumerate}
    \item \textbf{Forme développée :}
    \begin{align*}
    (2x+1)^2 &= 4x^2 + 4x + 1,\\
    (2x+1)(x-3) &= 2x^2 - 6x + x - 3 = 2x^2 - 5x - 3.\\
    E &= 4x^2 + 4x + 1 - (2x^2 - 5x - 3) = 4x^2 + 4x + 1 - 2x^2 + 5x + 3.\\
    E &= 2x^2 + 9x + 4.
    \end{align*}
    \item \textbf{Forme factorisée (en conservant le facteur commun $(2x+1)$) :}
    \[E = (2x+1)\big[(2x+1) - (x-3)\big] = (2x+1)(2x + 1 - x + 3) = \mathbf{(2x+1)(x+4)}.\]
    \emph{Vérification :} $(2x+1)(x+4) = 2x^2 + 8x + x + 4 = 2x^2 + 9x + 4$. \checkmark
    \item On choisit la \textbf{forme factorisée}, car elle permet d'appliquer directement la \cle{propriété du produit nul} :
    \[E = 0 \iff (2x+1)(x+4) = 0 \iff 2x+1 = 0 \ \text{ou}\ x+4 = 0.\]
    Soit $x = -\dfrac{1}{2}$ ou $x = -4$. \quad $S = \left\{-4\ ;\ -\dfrac{1}{2}\right\}$.
\end{enumerate}
\textbf{Point de vigilance :} la forme développée $2x^2 + 9x + 4 = 0$ est une équation du second degré ; sa résolution n'est pas au programme de 3ème (seules les équations du premier degré le sont). C'est précisément la factorisation qui rend la résolution possible ici.
\end{corrige}

\begin{corrige}[Exercice 8 — Fractions rationnelles : valeurs interdites et simplification]
\begin{enumerate}
    \item $F = \dfrac{x^2 - 4x + 4}{x^2 - 4}$.
    Dénominateur : $x^2 - 4 = (x-2)(x+2) = 0$ pour $x = 2$ ou $x = -2$.
    \cle{Valeurs interdites} : $\mathbf{2}$ et $\mathbf{-2}$.
    $F = \dfrac{(x-2)^2}{(x-2)(x+2)} = \mathbf{\dfrac{x-2}{x+2}}$ pour $x \neq 2$ et $x \neq -2$.
    \item $G = \dfrac{3x^2 - 6x}{x^2 - 4}$.
    Valeurs interdites : $x = 2$ et $x = -2$ (même dénominateur).
    Numérateur : $3x^2 - 6x = 3x(x-2)$.
    $G = \dfrac{3x(x-2)}{(x-2)(x+2)} = \mathbf{\dfrac{3x}{x+2}}$ pour $x \neq 2$ et $x \neq -2$.
    \item $H = \dfrac{x^2 - 9}{x^2 + 6x + 9}$.
    Dénominateur : $x^2 + 6x + 9 = (x+3)^2 = 0$ pour $x = -3$.
    Valeur interdite : $\mathbf{-3}$.
    $H = \dfrac{(x-3)(x+3)}{(x+3)^2} = \mathbf{\dfrac{x-3}{x+3}}$ pour $x \neq -3$.
\end{enumerate}
\textbf{Point de vigilance :} les valeurs interdites se déterminent sur le dénominateur \emph{avant} simplification et restent interdites même après simplification (la fraction simplifiée n'est équivalente à l'initiale que hors de ces valeurs).
\end{corrige}

\begin{corrige}[Exercice 9 — Problème géométrique : Aire d'un terrain]
\begin{enumerate}
    \item $\mathcal{A} = \text{longueur} \times \text{largeur} = (2x+5)(x+1) = 2x^2 + 2x + 5x + 5$.
    \[\mathcal{A} = 2x^2 + 7x + 5 \ \text{(en m}^2\text{)}.\]
    \item Pour $x = 10$ : $\mathcal{A} = 2(10)^2 + 7(10) + 5 = 200 + 70 + 5 = \mathbf{275 \text{ m}^2}$.
    \item On résout $\mathcal{A} = 90$, soit $2x^2 + 7x + 5 = 90 \iff 2x^2 + 7x - 85 = 0$.
    En 3ème, on privilégie la forme factorisée de l'aire pour chercher une solution entière :
    $(2x+5)(x+1) = 90$.
    Testons des valeurs entières positives pour $x$ :
    $x = 5 \Rightarrow (2\times5+5)(5+1) = 15 \times 6 = 90$. \checkmark
    Donc $\mathbf{x = 5}$ m.
    \emph{Vérification des dimensions :} longueur $= 2(5)+5 = 15$ m, largeur $= 5+1 = 6$ m, aire $= 15 \times 6 = 90$ m$^2$. Les dimensions sont bien positives : la valeur convient.
    \item Périmètre : $P = 2(L + l) = 2(15 + 6) = \mathbf{42 \text{ m}}$.
\end{enumerate}
\textbf{Barème indicatif (4 pts) :} 1 pt pour l'aire développée, 0,5 pt pour le calcul en $x=10$, 1,5 pt pour la résolution avec vérification, 1 pt pour le périmètre.\\
\textbf{Points de vigilance :} citer les unités (m$^2$, m) ; vérifier explicitement que la solution trouvée donne des dimensions positives.
\end{corrige}

\begin{corrige}[Exercice 10 — Modélisation : Tarification d'un transporteur]
\begin{enumerate}
    \item Pour $d \le 20$ : $P(d) = \mathbf{500 + 150d}$ (FCFA).
    \item Pour $d > 20$ : remise de $5\%$, donc on multiplie par $0{,}95$ (ou $\num{0,95}$) :
    \[P(d) = 0{,}95(500 + 150d) = 0{,}95\times 500 + 0{,}95\times 150d = \mathbf{475 + 142{,}5d}.\]
    \item Cherchons d'abord dans le cas $d \le 20$ : $500 + 150d = 4\,750 \iff 150d = 4\,250 \iff d = \dfrac{4\,250}{150} \approx 28{,}3$ km. Or $28{,}3 > 20$ : contradiction avec l'hypothèse.
    On est donc dans le cas $d > 20$ : $475 + 142{,}5d = 4\,750 \iff 142{,}5d = 4\,275 \iff d = \dfrac{4\,275}{142{,}5} = 30$.
    Le client a parcouru $\mathbf{d = 30}$ km (cas $d > 20$, cohérent).
    \emph{Vérification :} $0{,}95\times(500 + 150\times 30) = 0{,}95\times 5\,000 = 4\,750$ FCFA. \checkmark
    \item On compare pour $d > 20$ : $P_{\text{new}}(d) < P(d) \iff 800 + 140d < 475 + 142{,}5d$.
    \[800 - 475 < 142{,}5d - 140d \iff 325 < 2{,}5d \iff d > \dfrac{325}{2{,}5} = 130.\]
    Le nouveau tarif devient plus avantageux à partir de $\mathbf{131}$ km (arrondi au km supérieur).
\end{enumerate}
\textbf{Barème indicatif (4 pts) :} 0,5 pt par expression (questions 1 et 2), 1,5 pt pour la question 3 (avec justification du cas), 1,5 pt pour l'inéquation et sa résolution.\\
\textbf{Points de vigilance :} à la question 3, il faut \emph{tester les deux cas} ou vérifier la cohérence de la solution avec le cas choisi ; à la question 4, résoudre une inéquation du 1er degré en justifiant chaque étape.
\end{corrige}

\begin{corrige}[Exercice 11 — Synthèse algébrique : Identités et factorisation]
\begin{enumerate}
    \item $(x+3)^2 = x^2 + 6x + 9$ et $(x-3)(x+3) = x^2 - 9$.
    $A = x^2 + 6x + 9 - (x^2 - 9) = x^2 + 6x + 9 - x^2 + 9 = \mathbf{6x + 18}$.
    \emph{Remarque :} les termes en $x^2$ s'éliminent, $A$ est du premier degré.
    Factorisation : $A = 6x + 18 = \mathbf{6(x + 3)}$.
    \emph{Autre méthode (factorisation directe) :} $A = (x+3)\big[(x+3)-(x-3)\big] = (x+3)(6) = 6(x+3)$. \checkmark
    \item $B = \dfrac{x^2 - 9}{x^2 + 6x + 9}$.
    Dénominateur : $(x+3)^2 = 0 \iff x = -3$. \cle{Valeur interdite} : $\mathbf{-3}$.
    $B = \dfrac{(x-3)(x+3)}{(x+3)^2} = \mathbf{\dfrac{x-3}{x+3}}$ pour $x \neq -3$.
    \item $C = A \times B = 6(x+3)\times\dfrac{x-3}{x+3} = \mathbf{6(x-3)}$ pour $x \neq -3$, soit $C = 6x - 18$ ($x \neq -3$).
    \item $C = 2 \iff 6x - 18 = 2 \iff 6x = 20 \iff x = \dfrac{20}{6} = \dfrac{10}{3}$.
    $\dfrac{10}{3} \neq -3$, donc la solution est valable : $S = \left\{\dfrac{10}{3}\right\}$.
    \item $x = -3$ \textbf{n'est pas une solution}, car $-3$ est une \cle{valeur interdite} : la fraction $B$ (et donc $C$) n'existe pas pour $x = -3$, même si l'écriture simplifiée $6x - 18$ donnerait $-36 \neq 2$ de toute façon. La valeur interdite du dénominateur initial doit toujours être exclue.
\end{enumerate}
\textbf{Barème indicatif (5 pts) :} 1 pt pour $A$ développé et factorisé, 1 pt pour les valeurs interdites et la simplification de $B$, 1 pt pour $C$, 1 pt pour la résolution, 1 pt pour la justification de la question 5.\\
\textbf{Points de vigilance :} citer la condition $x \neq -3$ à chaque étape où l'on simplifie ; ne jamais conclure sans vérifier que la solution n'est pas une valeur interdite.
\end{corrige}

\newpage

\coursec{Fiche bilan}

\begin{popfigure}
\begin{center}\tcbox[colback=popblue,colframe=popblue,coltext=white,arc=9pt,boxsep=4pt,left=6pt,right=6pt]{\bfseries\small Calcul littéral}\end{center}
\vspace{-2pt}
\begin{tcbraster}[raster columns=3,raster equal height=rows,raster column skip=6pt,raster row skip=6pt,enhanced,blankest]
\begin{tcolorbox}[enhanced,colback=white,colframe=poporange,boxrule=0.9pt,arc=3pt,title={Expressions littérales},fonttitle=\bfseries\scriptsize,coltitle=white,colbacktitle=poporange,left=4pt,right=4pt,top=2pt,bottom=2pt,boxsep=1pt,toptitle=1.5pt,bottomtitle=1.5pt,halign=left]
\begin{itemize}[nosep,leftmargin=*,itemsep=1pt,font=\scriptsize]
\item Variable, valeur numérique
\item Traduire une situation
\end{itemize}
\end{tcolorbox}
\begin{tcolorbox}[enhanced,colback=white,colframe=popgreen,boxrule=0.9pt,arc=3pt,title={Monômes, polynômes},fonttitle=\bfseries\scriptsize,coltitle=white,colbacktitle=popgreen,left=4pt,right=4pt,top=2pt,bottom=2pt,boxsep=1pt,toptitle=1.5pt,bottomtitle=1.5pt,halign=left]
\begin{itemize}[nosep,leftmargin=*,itemsep=1pt,font=\scriptsize]
\item Coefficient, degré
\item Somme, produit
\end{itemize}
\end{tcolorbox}
\begin{tcolorbox}[enhanced,colback=white,colframe=popgold,boxrule=0.9pt,arc=3pt,title={Développer, réduire},fonttitle=\bfseries\scriptsize,coltitle=white,colbacktitle=popgold,left=4pt,right=4pt,top=2pt,bottom=2pt,boxsep=1pt,toptitle=1.5pt,bottomtitle=1.5pt,halign=left]
\begin{itemize}[nosep,leftmargin=*,itemsep=1pt,font=\scriptsize]
\item $k(a+b)=ka+kb$
\item Identités remarquables
\end{itemize}
\end{tcolorbox}
\begin{tcolorbox}[enhanced,colback=white,colframe=poppurple,boxrule=0.9pt,arc=3pt,title={Factoriser},fonttitle=\bfseries\scriptsize,coltitle=white,colbacktitle=poppurple,left=4pt,right=4pt,top=2pt,bottom=2pt,boxsep=1pt,toptitle=1.5pt,bottomtitle=1.5pt,halign=left]
\begin{itemize}[nosep,leftmargin=*,itemsep=1pt,font=\scriptsize]
\item Facteur commun
\item $a^2-b^2$, $a^2\pm2ab+b^2$
\end{itemize}
\end{tcolorbox}
\begin{tcolorbox}[enhanced,colback=white,colframe=poppink,boxrule=0.9pt,arc=3pt,title={Fractions rationnelles},fonttitle=\bfseries\scriptsize,coltitle=white,colbacktitle=poppink,left=4pt,right=4pt,top=2pt,bottom=2pt,boxsep=1pt,toptitle=1.5pt,bottomtitle=1.5pt,halign=left]
\begin{itemize}[nosep,leftmargin=*,itemsep=1pt,font=\scriptsize]
\item Condition d'existence
\item Simplifier, opérer
\end{itemize}
\end{tcolorbox}
\begin{tcolorbox}[enhanced,colback=white,colframe=popblue!60,boxrule=0.9pt,arc=3pt,title={Vie courante},fonttitle=\bfseries\scriptsize,coltitle=white,colbacktitle=popblue!60,left=4pt,right=4pt,top=2pt,bottom=2pt,boxsep=1pt,toptitle=1.5pt,bottomtitle=1.5pt,halign=left]
\begin{itemize}[nosep,leftmargin=*,itemsep=1pt,font=\scriptsize]
\item Aires, périmètres, budgets
\end{itemize}
\end{tcolorbox}
\end{tcbraster}

\legende{Carte mentale de la leçon : le calcul littéral en un coup d'œil.}
\end{popfigure}

\begin{bilan}
\textbf{1. Définitions clés}
\begin{itemize}
  \item Une \cle{expression littérale} est une expression contenant une ou plusieurs lettres (variables) représentant des nombres. Sa \cle{valeur numérique} s'obtient en remplaçant les lettres par des nombres donnés.
  \item Un \cle{monôme} est un produit d'un nombre (le \cle{coefficient}) par des puissances de lettres : $5x^2$ a pour coefficient $5$ et pour degré $2$.
  \item Un \cle{polynôme} est une somme de monômes ; il est \cle{réduit} lorsqu'on a regroupé les termes de même degré.
  \item \cle{Développer} : transformer un produit en somme. \cle{Factoriser} : transformer une somme en produit.
  \item Une \cle{fraction rationnelle} est un quotient de deux expressions littérales ; elle n'existe que si son dénominateur est non nul.
\end{itemize}

\textbf{2. Formules à connaître par cœur}
\begin{center}
\begin{tabularx}{\linewidth}{|l|X|}
\hline
\rowcolor{popblueL} \textbf{Nom} & \textbf{Formule} \\
\hline
Distributivité simple & $k(a+b)=ka+kb$ \quad et \quad $k(a-b)=ka-kb$ \\
\hline
Double distributivité & $(a+b)(c+d)=ac+ad+bc+bd$ \\
\hline
Identité remarquable 1 & $(a+b)^2=a^2+2ab+b^2$ \\
\hline
Identité remarquable 2 & $(a-b)^2=a^2-2ab+b^2$ \\
\hline
Identité remarquable 3 & $(a+b)(a-b)=a^2-b^2$ \\
\hline
Fractions rationnelles & $\dfrac{A}{B}$ existe si et seulement si $B\neq 0$ ; $\dfrac{A\times C}{B\times C}=\dfrac{A}{B}$ \\
\hline
\end{tabularx}
\end{center}

\textbf{3. Méthodes express}
\begin{itemize}
  \item \textbf{Réduire} : regrouper les termes en $x^2$, puis en $x$, puis les constantes. Ex. : $3x^2-2x+5x^2+x=8x^2-x$.
  \item \textbf{Développer} $(x+3)(2x-5)$ : $2x^2-5x+6x-15=2x^2+x-15$.
  \item \textbf{Factoriser} : chercher d'abord un \cle{facteur commun} ($6x^2+9x=3x(2x+3)$), sinon repérer une identité remarquable ($x^2-25=(x-5)(x+5)$).
  \item \textbf{Fraction rationnelle} : factoriser numérateur et dénominateur, écrire la condition d'existence, puis simplifier.
  \item \textbf{Vérifier} un développement ou une factorisation : tester avec une valeur simple ($x=0$ ou $x=1$).
\end{itemize}
\end{bilan}

\begin{aretenir}[Checklist avant le BEPC]
\begin{itemize}
  \item[$\square$] Je sais calculer la valeur numérique d'une expression littérale pour une valeur donnée de la variable.
  \item[$\square$] Je reconnais un monôme, son coefficient et son degré, et je sais réduire un polynôme.
  \item[$\square$] Je maîtrise la distributivité simple : $k(a+b)=ka+kb$.
  \item[$\square$] Je sais développer $(a+b)(c+d)$ sans oublier aucun terme.
  \item[$\square$] Je connais par cœur les trois identités remarquables.
  \item[$\square$] Je sais factoriser en mettant un facteur commun en évidence.
  \item[$\square$] Je sais factoriser avec $a^2-b^2=(a-b)(a+b)$ et $a^2\pm 2ab+b^2=(a\pm b)^2$.
  \item[$\square$] Je pense toujours à écrire la condition d'existence d'une fraction rationnelle (dénominateur non nul).
  \item[$\square$] Je factorise avant de simplifier une fraction rationnelle.
  \item[$\square$] Je vérifie mes résultats en testant une valeur numérique simple.
  \item[$\square$] Je sais traduire un problème concret (aire, périmètre, dépense) par une expression littérale, puis la développer ou la factoriser pour conclure.
  \item[$\square$] Je rédige ma démarche proprement : étapes alignées, justification, conclusion en français.
\end{itemize}
\end{aretenir}

\findecours

\end{document}
